% Fonctions dérivées — Mathématiques 1ère A — Cours signé OnBuch+
% Programme officiel MINESEC. Compiler avec : tectonic NA03-fonctions-derivees-a.tex
\newcommand{\DOCMATIERE}{Mathématiques}
\newcommand{\DOCNIVEAU}{1ère A}
\newcommand{\DOCMODULE}{Relations et opérations fondamentales dans l'ensemble des nombres réels}
\newcommand{\DOCLECON}{NA03}
\newcommand{\DOCTITRE}{Fonctions dérivées}
% OnBuch+ — préambule « cours signés » ENRICHI (compilé avec tectonic / XeLaTeX)
% Système visuel « Pop » d'OnBuch+ : fond crème (jamais blanc), bordures encre
% épaisses, ombres « dures » décalées, titres Archivo Black en capitales.
% Version MATHÉMATIQUES : graphes (pgfplots/tikz), boîtes pédagogiques
% (définition, propriété, méthode, attention, le savais-tu, exercice résolu,
% exercice, TP, bilan), page de garde et signature OnBuch+ ancrée en bas.
%
% Macros attendues dans main.tex AVANT \input{preamble} :
%   \DOCMATIERE  \DOCNIVEAU  \DOCMODULE  \DOCLECON (numéro)  \DOCTITRE
\documentclass[11pt]{article}

\usepackage{fontspec}
\usepackage[french]{babel}
\usepackage[a4paper,margin=2cm,top=3.4cm,bottom=2.8cm,headheight=1.4cm,footskip=1.3cm]{geometry}
\usepackage{amsmath,amssymb,amsfonts}
\usepackage{mathtools} % \xmapsto, \coloneqq... souvent utilisés par les modèles
\usepackage{cancel} % simplifications barrées
% \abs{z} (module, valeur absolue) et \norm{u} (norme), utilisés par les modèles
\providecommand{\abs}{}\let\abs\relax\DeclarePairedDelimiter{\abs}{\lvert}{\rvert}
\providecommand{\norm}{}\let\norm\relax\DeclarePairedDelimiter{\norm}{\lVert}{\rVert}
\usepackage[table]{xcolor} % [table] : fournit aussi \rowcolors (alternance de lignes)
\usepackage{pagecolor}
\usepackage{tikz}
\usetikzlibrary{arrows.meta,positioning,calc,shapes.geometric,shapes.misc,shapes.arrows,decorations.pathmorphing,decorations.pathreplacing,decorations.markings,patterns,shadows,mindmap,trees,fit,backgrounds,matrix,intersections,angles,quotes}
\usepackage{pgfplots}
\usepackage{pgf-pie} % diagrammes circulaires \pie (statistiques)
\pgfplotsset{compat=1.18}
\usepgfplotslibrary{fillbetween,groupplots} % aires entre courbes (calcul intégral)
\usepackage{enumitem}
\usepackage{fancyhdr}
% [most] charge toutes les bibliothèques tcolorbox (listings, minted, raster,
% documentation...) et gonfle inutilement la mémoire TeX sur un document long
% (~300 boîtes) : on ne charge que ce qui est réellement utilisé.
\usepackage[skins,breakable]{tcolorbox}
\usepackage{graphicx}
\usepackage{colortbl}
\usepackage{booktabs}
\usepackage{tabularx}
\usepackage{diagbox} % case d'en-tête diagonale des tableaux à double entrée
% Colonnes extensibles centrée / gauche / droite (C, L, R), souvent utilisées par les modèles
\newcolumntype{C}{>{\centering\arraybackslash}X}
\newcolumntype{L}{>{\raggedright\arraybackslash}X}
\newcolumntype{R}{>{\raggedleft\arraybackslash}X}
\usepackage{array}
\usepackage{multicol}
\usepackage{siunitx}
% parse-numbers=false : accepte aussi \SI{2,5 \times 10^{-3}}{...}
\sisetup{locale=FR,output-decimal-marker={,},group-minimum-digits=4,parse-numbers=false}
\DeclareSIUnit\ohm{\ensuremath{\symup{\Omega}}} % U+2126 absent de la police
\usepackage{qrcode}
% \degree : commande fréquemment utilisée par les modèles pour un angle ou une
% température en dehors de \SI{}{} (non fournie par les paquets chargés).
\providecommand{\degree}{^{\circ}}
% \qed (fin de démonstration), utilisé par les modèles sans amsthm
\providecommand{\qed}{\hfill$\square$}
% \barre : double barre des valeurs interdites dans un tableau de variations (tkz-tab)
\providecommand{\barre}{\ensuremath{\|}}
% environnement proof (amsthm non chargé) utilisé par les modèles
\ifdefined\proof\else\newenvironment{proof}[1][Démonstration]{\par\noindent\textit{#1.}\ }{\qed\par}\fi
\usepackage{lastpage}
\usepackage{hyperref}
\hypersetup{hidelinks}

% ============================================================
% Palette « Pop » OnBuch+ — mêmes hex que la classe P (plus_ui.dart)
% ============================================================
\definecolor{poporange}{HTML}{F59321}
\definecolor{popdark}{HTML}{E07A0C}
\definecolor{popcream}{HTML}{FFF6E8}
\definecolor{popink}{HTML}{1C1714}
\definecolor{poppurple}{HTML}{7A5AE0}
\definecolor{popgreen}{HTML}{2E9E6A}
\definecolor{poppink}{HTML}{E0457B}
\definecolor{popblue}{HTML}{2F7DE1}
\definecolor{popgold}{HTML}{C98A12}
\definecolor{popmuted}{HTML}{8A7F76}
\definecolor{popline}{HTML}{EADFD2}
\definecolor{popgray}{HTML}{8A7F76}
\colorlet{popgrey}{popgray}
% Alias tolérants (le modèle nomme parfois les couleurs différemment)
\colorlet{popwhite}{white}
\colorlet{popblack}{popink}
\colorlet{popred}{poppink}
\colorlet{popyellow}{popgold}
% Teintes claires pour fonds de tableaux / zones
\colorlet{popblueL}{popblue!10!white}
\colorlet{poporangeL}{poporange!14!white}
\colorlet{popgreenL}{popgreen!12!white}
\colorlet{poppurpleL}{poppurple!12!white}
\colorlet{poppinkL}{poppink!12!white}
\colorlet{popgoldL}{popgold!14!white}

\pagecolor{popcream}

% ============================================================
% Typographie
% ============================================================
\defaultfontfeatures{Ligatures=TeX}
\setmainfont{PlusJakartaSans-Medium.ttf}[Path=../../../fonts/,BoldFont=PlusJakartaSans-Bold.ttf]
% Maths en sans-serif (Fira Math) avec chiffres et lettres droites de Plus
% Jakarta, pour que les formules se fondent dans le texte.
\usepackage[math-style=ISO,bold-style=ISO]{unicode-math}
\setmathfont{FiraMath-Regular.otf}[Path=../../../fonts/]
\setmathfont{PlusJakartaSans-Medium.ttf}[Path=../../../fonts/,range={up/num,up/{Latin,latin},\mathup}]
\usepackage{icomma} % virgule décimale sans espace en mode math
% Caractères mathématiques Unicode absents des polices de texte (Plus Jakarta,
% Archivo Black) : rendus par leur équivalent mathématique, en texte comme en
% formule — sinon ils s'affichent en carré vide (ex. « Arithmétique dans ℤ »).
\usepackage{newunicodechar}
\newunicodechar{ℝ}{\ensuremath{\mathbb{R}}}
\newunicodechar{ℤ}{\ensuremath{\mathbb{Z}}}
\newunicodechar{ℂ}{\ensuremath{\mathbb{C}}}
\newunicodechar{ℕ}{\ensuremath{\mathbb{N}}}
\newunicodechar{ℚ}{\ensuremath{\mathbb{Q}}}
\newunicodechar{𝔻}{\ensuremath{\mathbb{D}}}
\newunicodechar{α}{\ensuremath{\alpha}}
\newunicodechar{β}{\ensuremath{\beta}}
\newunicodechar{γ}{\ensuremath{\gamma}}
\newunicodechar{δ}{\ensuremath{\delta}}
\newunicodechar{ε}{\ensuremath{\varepsilon}}
\newunicodechar{θ}{\ensuremath{\theta}}
\newunicodechar{λ}{\ensuremath{\lambda}}
\newunicodechar{μ}{\ensuremath{\mu}}
\newunicodechar{σ}{\ensuremath{\sigma}}
\newunicodechar{τ}{\ensuremath{\tau}}
\newunicodechar{φ}{\ensuremath{\varphi}}
\newunicodechar{ψ}{\ensuremath{\psi}}
\newunicodechar{ω}{\ensuremath{\omega}}
\newunicodechar{Σ}{\ensuremath{\Sigma}}
% majuscules produites par \MakeUppercase dans les titres (α-aminés -> Α)
\newunicodechar{Α}{\ensuremath{\alpha}}
\newunicodechar{Β}{\ensuremath{\beta}}
\newunicodechar{Γ}{\ensuremath{\Gamma}}
\newunicodechar{Δ}{\ensuremath{\Delta}}
\newunicodechar{Ε}{\ensuremath{\varepsilon}}
\newunicodechar{Θ}{\ensuremath{\Theta}}
\newunicodechar{Λ}{\ensuremath{\Lambda}}
\newunicodechar{Μ}{\ensuremath{\mu}}
\newunicodechar{Τ}{\ensuremath{\tau}}
\newunicodechar{Φ}{\ensuremath{\Phi}}
\newunicodechar{Ψ}{\ensuremath{\Psi}}
\newunicodechar{Ω}{\ensuremath{\Omega}}
\newunicodechar{↦}{\ensuremath{\mapsto}}
\newunicodechar{∈}{\ensuremath{\in}}
\newunicodechar{∉}{\ensuremath{\notin}}
\newunicodechar{∧}{\ensuremath{\wedge}}
\newunicodechar{⇔}{\ensuremath{\Leftrightarrow}}
\newunicodechar{⟺}{\ensuremath{\Longleftrightarrow}}
\newunicodechar{⇒}{\ensuremath{\Rightarrow}}
\newunicodechar{⟹}{\ensuremath{\Longrightarrow}}
\newunicodechar{∘}{\ensuremath{\circ}}
\newunicodechar{∪}{\ensuremath{\cup}}
\newunicodechar{∩}{\ensuremath{\cap}}
\newunicodechar{≡}{\ensuremath{\equiv}}
\newunicodechar{⊂}{\ensuremath{\subset}}
\newunicodechar{⊥}{\ensuremath{\perp}}
\newunicodechar{∃}{\ensuremath{\exists}}
\newunicodechar{∀}{\ensuremath{\forall}}
\newunicodechar{∛}{\ensuremath{\sqrt[3]{\,}}}
\newunicodechar{∜}{\ensuremath{\sqrt[4]{\,}}}
\newunicodechar{⃗}{\ensuremath{{}^{\to}}}
\newunicodechar{ⁿ}{\ifmmode^{n}\else\textsuperscript{\ensuremath{n}}\fi}
\newunicodechar{ˣ}{\ifmmode^{x}\else\textsuperscript{\ensuremath{x}}\fi}
\newunicodechar{ᵇ}{\ifmmode^{b}\else\textsuperscript{\ensuremath{b}}\fi}
\newunicodechar{ʳ}{\ifmmode^{r}\else\textsuperscript{\ensuremath{r}}\fi}
\newunicodechar{ᵘ}{\ifmmode^{u}\else\textsuperscript{\ensuremath{u}}\fi}
\newunicodechar{ʸ}{\ifmmode^{y}\else\textsuperscript{\ensuremath{y}}\fi}
\newunicodechar{ᵃ}{\ifmmode^{a}\else\textsuperscript{\ensuremath{a}}\fi}
\newunicodechar{ᵉ}{\ifmmode^{e}\else\textsuperscript{\ensuremath{e}}\fi}
\newunicodechar{ᵏ}{\ifmmode^{k}\else\textsuperscript{\ensuremath{k}}\fi}
\newunicodechar{ᵖ}{\ifmmode^{p}\else\textsuperscript{\ensuremath{p}}\fi}
\newunicodechar{ᴺ}{\ifmmode^{N}\else\textsuperscript{\ensuremath{N}}\fi}
\newunicodechar{ᶜ}{\ifmmode^{c}\else\textsuperscript{\ensuremath{c}}\fi}
\newunicodechar{ʰ}{\ifmmode^{h}\else\textsuperscript{\ensuremath{h}}\fi}
\newunicodechar{ᵠ}{\ifmmode^{q}\else\textsuperscript{\ensuremath{q}}\fi}
\newunicodechar{ₙ}{\ifmmode_{n}\else\textsubscript{\ensuremath{n}}\fi}
\newunicodechar{ₐ}{\ifmmode_{a}\else\textsubscript{\ensuremath{a}}\fi}
\newunicodechar{ᵢ}{\ifmmode_{i}\else\textsubscript{\ensuremath{i}}\fi}
\newunicodechar{ᵣ}{\ifmmode_{r}\else\textsubscript{\ensuremath{r}}\fi}
\newunicodechar{ₖ}{\ifmmode_{k}\else\textsubscript{\ensuremath{k}}\fi}
\newunicodechar{ᵦ}{\ifmmode_{\beta}\else\textsubscript{\ensuremath{\beta}}\fi}
\newunicodechar{ₓ}{\ifmmode_{x}\else\textsubscript{\ensuremath{x}}\fi}
\newfontfamily\popdisplay{ArchivoBlack-Regular.ttf}[Path=../../../fonts/]
\newfontfamily\poplogo{SpaceGrotesk-Bold.ttf}[Path=../../../fonts/]
\newfontfamily\popsemi{PlusJakartaSans-SemiBold.ttf}[Path=../../../fonts/]
\newfontfamily\popxbold{PlusJakartaSans-ExtraBold.ttf}[Path=../../../fonts/]
\newfontfamily\popxboldls{PlusJakartaSans-ExtraBold.ttf}[Path=../../../fonts/,LetterSpace=3.0]

\setlist[enumerate]{leftmargin=1.7em,itemsep=5pt,topsep=4pt}
\setlist[itemize]{leftmargin=1.7em,itemsep=4pt,topsep=4pt,label={\color{poporange}\textbullet}}
\setlength{\parindent}{0pt}
\setlength{\parskip}{5pt}
\renewcommand{\arraystretch}{1.25}

% pgfplots : style Pop par défaut
\pgfplotsset{
  every axis/.append style={axis line style={popink,line width=1pt},tick style={popink},
    grid style={popline,line width=0.5pt},label style={font=\small},tick label style={font=\footnotesize}},
  popcurve/.style={poporange,line width=1.8pt,no markers,smooth},
}
% Même style disponible dans un \draw TikZ simple (hors pgfplots)
\tikzset{popcurve/.style={poporange,line width=1.8pt}}

% ============================================================
% Pill / squiggle
% ============================================================
\newcommand{\poppill}[3]{%
  \tikz[baseline=-0.55ex]\node[draw=popink,line width=1.2pt,rounded corners=2.6pt,%
    fill=#2,inner xsep=6.5pt,inner ysep=3pt]{%
    \popxboldls\fontsize{7}{7}\selectfont\color{#3}\MakeUppercase{#1}};%
}
\newcommand{\popsquiggle}[1]{%
  \tikz[baseline=-0.4ex]{
    \draw[#1,line width=2.6pt,line cap=round]
      plot [smooth] coordinates {(0,0.09) (0.34,0.01) (0.68,0.21) (1.02,0.01) (1.36,0.10)};
  }%
}
% Mot-clé surligné (marqueur orange sous le texte)
\newcommand{\cle}[1]{{\popxbold\color{popdark}#1}}

% ============================================================
% En-tête / pied
% ============================================================
\pagestyle{fancy}
\fancyhf{}
\renewcommand{\headrulewidth}{0pt}
\renewcommand{\footrulewidth}{0pt}
\lhead{\raisebox{-0.15ex}{\poplogo\fontsize{13}{13}\selectfont\color{popink}OnBuch\color{poporange}+}}
\rhead{\poppill{\DOCMATIERE\ \textbullet\ \DOCNIVEAU\ \textbullet\ Leçon \DOCLECON}{white}{popink}}
\lfoot{\popxbold\fontsize{7.6}{7.6}\selectfont\color{popmuted}\MakeUppercase{Cours signé OnBuch+}\ \textbullet\ {\popsemi\DOCTITRE}}
\rfoot{\tikz[baseline=-0.6ex]\node[rounded corners=8.5pt,fill=popink,minimum height=17pt,minimum width=17pt,inner xsep=5pt,inner sep=0]{\popxbold\fontsize{8}{8}\selectfont\color{popcream}\thepage\,/\,\pageref*{LastPage}};}

% ============================================================
% Titres
% ============================================================
\newcommand{\chaptitre}[2]{%
  \begin{tcolorbox}[enhanced,colback=poporange,colframe=popink,boxrule=2.6pt,arc=11pt,
      drop shadow=popink,left=15pt,right=15pt,top=12pt,bottom=11pt,before skip=3pt,after skip=18pt]
    \poppill{#2}{popink}{popcream}\par\vspace{9pt}
    {\raggedright\hyphenpenalty=10000\exhyphenpenalty=10000\popdisplay\fontsize{22}{24}\selectfont\color{popcream}\MakeUppercase{#1}\par}
  \end{tcolorbox}
}
% Section numérotée : pastille orange avec le numéro + titre Archivo
\newcounter{popsec}
\newcommand{\coursec}[1]{%
  \stepcounter{popsec}\setcounter{popsub}{0}%
  \par\vspace{16pt}\noindent
  \begin{minipage}{\linewidth}
  \tikz[baseline=-0.7ex]\node[draw=popink,line width=1.6pt,fill=poporange,rounded corners=4pt,minimum size=22pt,inner sep=0]{\popdisplay\fontsize{12}{12}\selectfont\color{popcream}\thepopsec};%
  \hspace{8pt}{\popdisplay\fontsize{15}{17}\selectfont\color{popink}\MakeUppercase{#1}}\par
  \vspace{4pt}\noindent{\color{poporange}\rule{36pt}{3.6pt}}
  \end{minipage}\par\nopagebreak\vspace{8pt}
}
\newcounter{popsub}
\newcommand{\courssub}[1]{%
  \stepcounter{popsub}%
  \par\vspace{9pt}\noindent{\popxbold\fontsize{11.5}{13}\selectfont\color{popink}{\color{poporange}\thepopsec.\thepopsub}\hspace{6pt}#1}\par\nopagebreak\vspace{4pt}
}

% ============================================================
% Boîtes pédagogiques — même recette : fond blanc, bordure épaisse colorée,
% ombre dure encre, badge plein. Titre optionnel [#1].
% ============================================================
\tcbset{popbase/.style={breakable,enhanced,colback=white,boxrule=2.3pt,arc=9pt,
  shadow={3pt}{-3pt}{0pt}{popink},left=14pt,right=14pt,top=11pt,bottom=11pt,before skip=11pt,after skip=15pt,
  fontupper=\popsemi\fontsize{10.6}{14.5}\selectfont\color{popink},
  fontlower=\popsemi\fontsize{10.6}{14.5}\selectfont\color{popink}}}
% Coche dessinée (le glyphe ✓ n'existe pas dans les polices du document)
\newcommand{\popcheck}[1]{\tikz[baseline=-0.5ex]\draw[#1,line width=1.6pt,line cap=round,line join=round](-0.13,0.0)--(-0.04,-0.09)--(0.13,0.1);}
\newcommand{\popboxhead}[3]{\poppill{#1}{#2}{white}\ifx\relax#3\relax\else\hspace{6pt}{\popxbold\fontsize{10.6}{12}\selectfont\color{popink}#3}\fi\par\vspace{7pt}}

\newtcolorbox{aretenir}[1][]{popbase,colframe=popgreen,before upper={\popboxhead{À retenir}{popgreen}{#1}}}
\newtcolorbox{exemplebox}[1][]{popbase,colframe=poporange,before upper={\popboxhead{Exemple}{poporange}{#1}}}
\newtcolorbox{definition}[1][]{popbase,colframe=popblue,before upper={\popboxhead{Définition}{popblue}{#1}}}
\newtcolorbox{propriete}[1][]{popbase,colframe=poppurple,before upper={\popboxhead{Propriété}{poppurple}{#1}}}
\newtcolorbox{methode}[1][]{popbase,colframe=popgold,before upper={\popboxhead{Méthode}{popgold}{#1}}}
\newtcolorbox{attention}[1][]{popbase,colframe=poppink,before upper={\popboxhead{Attention}{poppink}{#1}}}
\newtcolorbox{experience}[1][]{popbase,colframe=popblue,colback=popblueL,before upper={\popboxhead{Expérience}{popblue}{#1}}}
% « Le savais-tu ? » : carte violette pleine (contexte, culture, Cameroun)
\newtcolorbox{savaistu}[1][]{popbase,colback=poppurple,colframe=popink,
  fontupper=\popsemi\fontsize{10.4}{14.2}\selectfont\color{white},
  before upper={\poppill{Le savais-tu\,?}{white}{poppurple}\ifx\relax#1\relax\else\hspace{6pt}{\popxbold\color{white}#1}\fi\par\vspace{7pt}}}
% Exercice résolu : énoncé en haut, \tcblower puis la solution sur fond crème
\newcounter{popexo}\newcounter{popexr}
\newtcolorbox{exoresolu}[1][]{popbase,colframe=poporange,colbacklower=popcream,
  segmentation style={popink,line width=1.2pt,dashed},
  before upper={\stepcounter{popexr}\popboxhead{Exercice résolu \thepopexr}{poporange}{#1}},
  before lower={\poppill{Solution}{popink}{popcream}\par\vspace{6pt}}}
% Exercice (non résolu en place) — la correction est dans la partie Corrigés
\newtcolorbox{exercice}[1][]{popbase,colframe=popink,
  before upper={\stepcounter{popexo}\popboxhead{Exercice \thepopexo}{popink}{#1}}}
% Corrigé
\newtcolorbox{corrige}[1][]{popbase,colframe=popgreen,colback=popgreenL,
  before upper={\popboxhead{Corrigé}{popgreen}{#1}}}
% Objectifs / prérequis / bilan
\newtcolorbox{objectifs}{popbase,colframe=popink,colback=poporangeL,
  before upper={\popboxhead{Objectifs de la leçon}{popink}{}}}
\newtcolorbox{prerequis}{popbase,colframe=popink,colback=popblueL,
  before upper={\popboxhead{Prérequis}{popblue}{}}}
\newtcolorbox{bilan}{popbase,colframe=popink,colback=white,boxrule=2.8pt,
  before upper={\popboxhead{Fiche bilan}{poporange}{L'essentiel à connaître pour l'examen}}}
% Figure encadrée avec légende
\newtcolorbox{popfigure}[1][]{enhanced,breakable=false,colback=white,colframe=popline,boxrule=1.4pt,arc=8pt,
  left=10pt,right=10pt,top=10pt,bottom=8pt,before skip=10pt,after skip=12pt,halign=center,
  lower separated=false,halign lower=center,
  fontlower=\popsemi\fontsize{9}{11.5}\selectfont\color{popmuted},
  #1}
\newcommand{\legende}[1]{\par\vspace{6pt}{\popsemi\fontsize{9}{11.5}\selectfont\color{popmuted}#1}}

% ============================================================
% Signature OnBuch+
% ============================================================
\newcommand{\poplogoinline}[1]{{\poplogo\fontsize{#1}{#1}\selectfont\color{popink}OnBuch\color{poporange}+}}
\newcommand{\signatureonbuch}{%
  \begin{tcolorbox}[enhanced,colback=popink,colframe=popink,boxrule=2.2pt,arc=10pt,
      drop shadow=poporange,left=14pt,right=14pt,top=10pt,bottom=10pt,before skip=0pt,after skip=0pt]
    \tikz[baseline=-0.7ex]\node[circle,fill=popgreen,draw=popcream,line width=1.3pt,minimum size=22pt,inner sep=0]{\popcheck{white}};%
    \hspace{9pt}\begin{minipage}[c]{0.62\linewidth}
      {\popxbold\fontsize{12}{13}\selectfont\color{popcream}Signé\ \textbullet\ {\poplogo OnBuch\color{poporange}+}}\par\vspace{2pt}
      {\raggedright\popsemi\fontsize{8}{10.5}\selectfont\color{popline}\MakeUppercase{Équipe pédagogique OnBuch+}\\\MakeUppercase{Conforme au programme officiel MINESEC}\par}
    \end{minipage}\hfill
    {\poplogo\fontsize{22}{22}\selectfont\color{popcream}OnBuch\color{poporange}+}
  \end{tcolorbox}%
}

% ============================================================
% Page de garde
% #1 titre · #2 matière · #3 niveau · #4 module · #5 n° leçon · #6 durée ·
% #7 description / objectifs (texte ou itemize)
% ============================================================
\newcommand{\pagegarde}[7]{%
  \thispagestyle{empty}
  \newgeometry{a4paper,margin=1.7cm,top=1.6cm,bottom=1.4cm}%
  \begin{tikzpicture}[remember picture,overlay]
    \fill[poporange!18] ([xshift=-3.2cm,yshift=-3.6cm]current page.north east) circle (4.6cm);
    \fill[poppurple!14] ([xshift=2.4cm,yshift=7.6cm]current page.south west) circle (3.8cm);
  \end{tikzpicture}%
  {\poplogo\fontsize{30}{30}\selectfont\color{popink}OnBuch\color{poporange}+}\hfill
  \raisebox{0.6ex}{\poppill{Cours signé \textbullet\ Premium}{popink}{popcream}}\par
  \vspace{1pt}\popsquiggle{poppurple}\par
  \vspace{8pt}
  \begin{tcolorbox}[enhanced,colback=poporange,colframe=popink,boxrule=2.8pt,arc=13pt,
      drop shadow=popink,left=18pt,right=18pt,top=12pt,bottom=13pt,after skip=10pt]
    \poppill{#2 \textbullet\ #3}{popink}{popcream}\hspace{5pt}\poppill{#4}{white}{popink}\par\vspace{8pt}
    {\popdisplay\fontsize{14}{14}\selectfont\color{popink}LEÇON #5}\par\vspace{4pt}
    {\raggedright\hyphenpenalty=10000\exhyphenpenalty=10000\popdisplay\fontsize{25}{27}\selectfont\color{popcream}\MakeUppercase{#1}\par}
    \vspace{8pt}
    {\popxbold\fontsize{9.5}{9.5}\selectfont\color{popink}\MakeUppercase{Durée conseillée : #6}}
  \end{tcolorbox}
  \begin{tcolorbox}[enhanced,colback=white,colframe=popink,boxrule=2.2pt,arc=10pt,
      drop shadow=popink,left=16pt,right=16pt,top=10pt,bottom=10pt,after skip=8pt]
    \poppill{Dans cette leçon}{poppurple}{white}\par\vspace{5pt}
    \popsemi\fontsize{9.3}{12.6}\selectfont\color{popink}#7
  \end{tcolorbox}
  \noindent\begin{minipage}[t]{0.487\linewidth}
    \begin{tcolorbox}[enhanced,colback=popgreen,colframe=popink,boxrule=2.2pt,arc=10pt,
        drop shadow=popink,left=11pt,right=11pt,top=10pt,bottom=10pt,height=3.1cm,valign=center]
      \tcbox[colback=white,colframe=popink,boxrule=1.3pt,arc=5pt,boxsep=0pt,nobeforeafter,
        left=3pt,right=3pt,top=3pt,bottom=3pt,valign=center]{\qrcode[height=1.9cm]{https://play.google.com/store/apps/details?id=com.luvvix.onbuch&hl=fr}}%
      \hspace{8pt}\begin{minipage}[c]{0.5\linewidth}
        {\popdisplay\fontsize{11}{12.5}\selectfont\color{white}\MakeUppercase{Télécharge OnBuch}}\par\vspace{4pt}
        {\raggedright\popsemi\fontsize{8.4}{11}\selectfont\color{white}Gratuit \textbullet\ Google Play\\Tuteur IA, annales, concours, résultats\par}
      \end{minipage}
    \end{tcolorbox}
  \end{minipage}\hfill
  \begin{minipage}[t]{0.487\linewidth}
    \begin{tcolorbox}[enhanced,colback=poppurple,colframe=popink,boxrule=2.2pt,arc=10pt,
        drop shadow=popink,left=11pt,right=11pt,top=10pt,bottom=10pt,height=3.1cm,valign=center]
      \tikz[baseline=-0.7ex]\node[circle,fill=white,draw=popink,line width=1.3pt,minimum size=1.6cm,inner sep=0]{\popdisplay\fontsize{24}{24}\selectfont\color{poppurple}+};%
      \hspace{8pt}\begin{minipage}[c]{0.6\linewidth}
        {\popdisplay\fontsize{11}{12.5}\selectfont\color{white}\MakeUppercase{Passe à OnBuch+}}\par\vspace{4pt}
        {\raggedright\popsemi\fontsize{8.4}{11}\selectfont\color{white}Tous les cours signés, Tuteur IA illimité, fiches PDF — onglet OnBuch+ de l'app\par}
      \end{minipage}
    \end{tcolorbox}
  \end{minipage}
  \vspace{8pt}
  \popcontenu
  \vfill
  \signatureonbuch
  \restoregeometry
  \newpage
}

% Rangée « Ce cours contient » (page de garde)
\newcommand{\poptile}[3]{% #1 couleur #2 gros texte #3 légende
  \begin{tcolorbox}[enhanced,colback=#1,colframe=popink,boxrule=2pt,arc=9pt,shadow={2.5pt}{-2.5pt}{0pt}{popink},
      left=6pt,right=6pt,top=9pt,bottom=9pt,halign=center,valign=center,height=1.75cm,width=\linewidth,nobeforeafter]
    {\popdisplay\fontsize{12.5}{13}\selectfont\color{white}\MakeUppercase{#2}}\par\vspace{3pt}
    {\centering\popsemi\fontsize{7.8}{9.5}\selectfont\color{white}#3\par}
  \end{tcolorbox}}
\newcommand{\popcontenu}{%
  \par\noindent{\popxboldls\fontsize{7.5}{7.5}\selectfont\color{popmuted}CE COURS SIGNÉ CONTIENT}\par\vspace{7pt}
  \noindent\begin{minipage}[t]{0.235\linewidth}\poptile{popblue}{Cours}{complet, illustré, pas à pas}\end{minipage}\hfill
  \begin{minipage}[t]{0.235\linewidth}\poptile{poporange}{Méthodes}{et exercices résolus}\end{minipage}\hfill
  \begin{minipage}[t]{0.235\linewidth}\poptile{poppink}{Exercices}{type examen + corrigés}\end{minipage}\hfill
  \begin{minipage}[t]{0.235\linewidth}\poptile{popgreen}{Fiche bilan}{l'essentiel à retenir}\end{minipage}\par}

% Sommaire
\newtcolorbox{sommaire}{enhanced,colback=white,colframe=popink,boxrule=2.2pt,arc=10pt,
  drop shadow=popink,left=18pt,right=18pt,top=13pt,bottom=13pt,before skip=6pt,after skip=20pt,
  before upper={\poppill{Sommaire}{poppurple}{white}\par\vspace{9pt}\popsemi\fontsize{10.6}{14.5}\selectfont\color{popink}}}

% Signature de fin de cours — poussée tout en bas de la dernière page
\newcommand{\findecours}{%
  \par\vfill\nopagebreak
  \begin{center}\popsquiggle{poporange}\end{center}\vspace{4pt}
  \signatureonbuch
}
\AtBeginDocument{\def\llbracket{\mathopen{[\mkern-3.5mu[}}\def\rrbracket{\mathclose{]\mkern-3.5mu]}}} % intervalles d'entiers (glyphe ⟦ absent de la police)
% Glyphes absents de Fira Math : redéfinis à partir de glyphes disponibles
\AtBeginDocument{%
  \def\setminus{\mathbin{\backslash}}\let\smallsetminus\setminus
  \def\vdots{\mathord{\vbox{\baselineskip3.2pt\lineskiplimit0pt\kern4pt\hbox{.}\hbox{.}\hbox{.}}}}%
  \def\ddots{\mathinner{\mkern1mu\raise6.5pt\hbox{.}\mkern2mu\raise3.6pt\hbox{.}\mkern2mu\raise0.7pt\hbox{.}\mkern1mu}}%
  \def\triangle{\mathord{\scalebox{0.9}{$\symup{\Delta}$}}}%
  \def\nabla{\mathord{\raisebox{\depth}{\scalebox{1}[-1]{$\symup{\Delta}$}}}}%
  \def\checkmark{\popcheck{popdark}}%
  \def\star{\mathbin{\ast}}\def\bigstar{\mathord{\ast}}%
  \def\diamond{\mathbin{\scalebox{0.6}{$\Diamond$}}}%
  \def\bigcup{\mathop{\mathchoice{\raisebox{-0.3em}{\scalebox{1.7}{$\cup$}}}{\scalebox{1.25}{$\cup$}}{\cup}{\cup}}\displaylimits}%
  \def\bigcap{\mathop{\mathchoice{\raisebox{-0.3em}{\scalebox{1.7}{$\cap$}}}{\scalebox{1.25}{$\cap$}}{\cap}{\cap}}\displaylimits}%
  \def\lesseqqgtr{\mathrel{\lessgtr}}\def\bigoplus{\mathop{\oplus}\displaylimits}%
}
\providecommand{\ssi}{si et seulement si} % abréviation fréquente
\AtBeginDocument{\providecommand{\wideparen}{\overparen}} % arc de cercle

%% FIGFIX-BEGIN v4 — correctifs globaux des figures (docs/onbuch-generation/tools/figfix)
\usepackage{adjustbox}\usepackage{etoolbox}
% toute figure TikZ/pgfplots plus large que la ligne est réduite à la largeur disponible
\BeforeBeginEnvironment{tikzpicture}{\begin{adjustbox}{max width=\linewidth}}
\AfterEndEnvironment{tikzpicture}{\end{adjustbox}}
% légendes pgfplots lisibles par-dessus les courbes
\pgfplotsset{every axis/.append style={legend style={fill=white,fill opacity=0.92,text opacity=1,draw=black!15,inner sep=3pt}}}
% étiquettes d'arêtes (syntaxe edge["..."]) avec fond blanc
\tikzset{every edge quotes/.append style={fill=white,inner sep=1.2pt}}
%% FIGFIX-END
\begin{document}

\pagegarde{Fonctions dérivées}{Mathématiques}{1ère A}{Relations et opérations fondamentales dans l'ensemble des nombres réels}{NA03}{10 h}{Cette leçon introduit la notion fondamentale d'analyse : la dérivée d'une fonction, d'abord en un point (nombre dérivé, tangente), puis globalement (fonction dérivée). Elle fournit les formules de dérivation des fonctions usuelles et des opérations, outils indispensables pour l'étude des variations et l'optimisation.
\vspace{4pt}
\begin{multicols}{2}\small
\begin{itemize}
\item À la fin de cette leçon, je sais définir le nombre dérivé d'une fonction en un réel x0 comme limite du taux d'accroissement
\item À la fin de cette leçon, je sais interpréter graphiquement la dérivabilité en x0 (existence d'une tangente non verticale)
\item À la fin de cette leçon, je sais déterminer le nombre dérivé d'une fonction en x0 par le calcul de limite
\item À la fin de cette leçon, je sais déterminer une équation cartésienne de la tangente à une courbe en un point d'abscisse x0
\item À la fin de cette leçon, je sais construire géométriquement la tangente à une courbe en un point
\item À la fin de cette leçon, je sais dériver les fonctions élémentaires (constantes, puissances, racine carrée, valeur absolue)
\end{itemize}\end{multicols}}

\begin{sommaire}
\begin{multicols}{2}
\begin{enumerate}
\item Taux d'accroissement et nombre dérivé
\item Interprétation graphique et tangente
\item Fonction dérivée et dérivées des fonctions élémentaires
\item Dérivées et opérations sur les fonctions
\item Applications et synthèse
\item Activité d'intégration
\item Méthodes et exercices résolus
\item Exercices
\item Corrigés des exercices
\item Fiche bilan
\end{enumerate}
\end{multicols}
\end{sommaire}

\begin{objectifs}
\begin{itemize}
\item À la fin de cette leçon, je sais définir le nombre dérivé d'une fonction en un réel x0 comme limite du taux d'accroissement
\item À la fin de cette leçon, je sais interpréter graphiquement la dérivabilité en x0 (existence d'une tangente non verticale)
\item À la fin de cette leçon, je sais déterminer le nombre dérivé d'une fonction en x0 par le calcul de limite
\item À la fin de cette leçon, je sais déterminer une équation cartésienne de la tangente à une courbe en un point d'abscisse x0
\item À la fin de cette leçon, je sais construire géométriquement la tangente à une courbe en un point
\item À la fin de cette leçon, je sais dériver les fonctions élémentaires (constantes, puissances, racine carrée, valeur absolue)
\item À la fin de cette leçon, je sais dériver une somme, un produit par un scalaire, un produit, un inverse et un quotient de fonctions dérivables
\item À la fin de cette leçon, je sais justifier la non-dérivabilité d'une fonction en un point (cas de |x| et de √x en 0)
\end{itemize}
\end{objectifs}

\begin{prerequis}
\begin{itemize}
\item Notion de fonction, ensemble de définition, courbe représentative
\item Calcul de limites (limites finies en un point, formes indéterminées simples)
\item Taux d'accroissement d'une fonction entre deux points
\item Équations de droites : coefficient directeur, équation y = ax + b
\item Fonctions usuelles : x ↦ xⁿ, x ↦ √x, x ↦ 1/x, x ↦ |x|
\end{itemize}
\end{prerequis}

\newpage

\coursec{Taux d'accroissement et nombre dérivé}

\emph{Un moto-taxi quitte le carrefour Total de Bafoussam et accélère sur la route de Foumban. À un instant précis, à quelle vitesse roule-t-il exactement ?} Son compteur affiche une valeur, mais d'où vient-elle ? La \cle{vitesse moyenne} sur 10 kilomètres — distance parcourue divisée par le temps écoulé — ne suffit pas : elle masque les accélérations et les freinages. Il faut un outil capable de donner la vitesse \emph{à un instant donné}, la \cle{vitesse instantanée}. De même, un commerçant au marché Mokolo veut savoir à quel rythme son bénéfice varie lorsqu'il augmente légèrement ses prix : non pas la variation totale, mais la variation \emph{au voisinage} de son prix actuel.

Cet outil mathématique, c'est la \cle{dérivation}. Dans cette section, nous allons le construire pas à pas : d'abord le taux d'accroissement, qui mesure une variation moyenne, puis le nombre dérivé, qui mesure une variation instantanée.

\courssub{Le taux d'accroissement : mesurer une variation moyenne}

\textbf{Intuition.} Lorsqu'une quantité $f(x)$ dépend d'une variable $x$, une question naturelle est : « quand $x$ augmente d'une certaine quantité, de combien $f(x)$ varie-t-elle, et à quel rythme ? ». Si le prix du kilogramme de tomates passe de $f(x)$ à $f(x+h)$ francs CFA quand la demande passe de $x$ à $x+h$ tonnes, le quotient $\dfrac{f(x+h)-f(x)}{h}$ indique de combien de francs le prix varie \emph{en moyenne par tonne supplémentaire}. C'est exactement l'idée du taux d'accroissement.

\begin{definition}[Taux d'accroissement]
Soit $f$ une fonction définie sur un intervalle $I$, $x_0$ un réel de $I$ et $h$ un réel non nul tel que $x_0 + h \in I$. On appelle \cle{taux d'accroissement} (ou taux de variation) de $f$ entre $x_0$ et $x_0 + h$ le réel :
\[ \tau(h) = \frac{f(x_0 + h) - f(x_0)}{h}. \]
\end{definition}

\begin{attention}[Le dénominateur n'est jamais nul]
Le taux d'accroissement n'est défini que pour $h \neq 0$. On ne peut pas « poser $h = 0$ » directement : c'est précisément toute la difficulté — et toute la richesse — de ce qui va suivre. De plus, l'accroissement $h$ peut être \textbf{négatif} : on regarde alors ce qui se passe « à gauche » de $x_0$.
\end{attention}

\textbf{Interprétation graphique : le coefficient directeur de la sécante.} Plaçons-nous dans un repère et considérons la courbe $\mathcal{C}_f$ de $f$. Notons $A$ le point d'abscisse $x_0$ et $M$ le point d'abscisse $x_0 + h$ :
\[ A\big(x_0 \,;\, f(x_0)\big) \qquad \text{et} \qquad M\big(x_0+h \,;\, f(x_0+h)\big). \]
La droite $(AM)$ est une \cle{sécante} de la courbe : elle la coupe en deux points. Son coefficient directeur vaut :
\[ a_{(AM)} = \frac{y_M - y_A}{x_M - x_A} = \frac{f(x_0+h) - f(x_0)}{(x_0+h) - x_0} = \frac{f(x_0+h) - f(x_0)}{h} = \tau(h). \]

\begin{aretenir}[Lecture graphique du taux d'accroissement]
Le taux d'accroissement $\tau(h) = \dfrac{f(x_0+h)-f(x_0)}{h}$ est le \cle{coefficient directeur de la sécante} $(AM)$, où $A\big(x_0\,;\,f(x_0)\big)$ et $M\big(x_0+h\,;\,f(x_0+h)\big)$.
\end{aretenir}

\begin{exemplebox}[Calcul d'un taux d'accroissement]
Soit $f(x) = x^2$, $x_0 = 1$ et $h = 2$. Alors :
\[ \tau(2) = \frac{f(3) - f(1)}{2} = \frac{9 - 1}{2} = 4. \]
Entre les abscisses $1$ et $3$, la fonction $x \mapsto x^2$ augmente en moyenne de $4$ unités par unité d'abscisse. La sécante joignant $A(1\,;\,1)$ à $M(3\,;\,9)$ a pour coefficient directeur $4$.
\end{exemplebox}

\courssub{De la sécante à la tangente : l'idée de la limite}

\textbf{Expérience de pensée.} Fixons le point $A$ et faisons glisser le point $M$ sur la courbe, de plus en plus près de $A$ : cela revient à faire tendre $h$ vers $0$. La sécante $(AM)$ pivote alors autour de $A$ et semble se rapprocher d'une position limite : une droite qui « effleure » la courbe en $A$, la \cle{tangente}. Le coefficient directeur de cette droite limite, s'il existe, est la limite des taux d'accroissement. C'est cette idée que formalise la définition suivante.

\begin{popfigure}
\begin{center}
\begin{tikzpicture}
\begin{axis}[
    width=0.85\linewidth, height=7cm,
    axis lines=middle,
    xlabel={$x$}, ylabel={$y$},
    xmin=-0.6, xmax=3.6, ymin=-0.6, ymax=10.5,
    xtick={1,2,3}, ytick={1,4,9},
    legend style={at={(0.03,0.97)}, anchor=north west, font=\small},
    clip=true
]
\addplot[popcurve, domain=-0.5:3.2, samples=200] {x^2};
\addlegendentry{$y = x^2$}
\addplot[poporange, thick, domain=-0.2:3.2, samples=2] {4*x - 3};
\addlegendentry{sécante $h=2$ : $\tau=4$}
\addplot[popgreen, thick, domain=-0.2:3.2, samples=2] {3*x - 2};
\addlegendentry{sécante $h=1$ : $\tau=3$}
\addplot[poppurple, thick, domain=-0.2:3.2, samples=2] {2.5*x - 1.5};
\addlegendentry{sécante $h=0.5$ : $\tau=2.5$}
\addplot[popink, very thick, domain=-0.2:3.2, samples=2] {2*x - 1};
\addlegendentry{tangente : pente $f'(1)=2$}
\addplot[only marks, mark=*, mark size=2pt, popdark] coordinates {(1,1) (2,4) (3,9) (1.5,2.25)};
\node[popdark, anchor=south west] at (axis cs:1.05,1.05) {$A$};
\node[popdark, anchor=south west] at (axis cs:3.05,9.1) {$M$};
\end{axis}
\end{tikzpicture}
\end{center}
\legende{Parabole $y = x^2$ au voisinage de $A(1\,;\,1)$ : quand $h$ diminue, la sécante $(AM)$ pivote et se rapproche de la tangente en $A$, de coefficient directeur $2$.}
\end{popfigure}

\courssub{Dérivabilité en un point et nombre dérivé}

Nous voici au cœur de la leçon. Lorsque $h$ tend vers $0$, deux cas peuvent se présenter : soit le taux $\tau(h)$ se rapproche d'un nombre réel bien déterminé, soit il ne se stabilise pas (il explose vers l'infini, ou il hésite entre plusieurs valeurs). Seul le premier cas nous intéresse.

\begin{definition}[Fonction dérivable en $x_0$, nombre dérivé]
Soit $f$ une fonction définie sur un intervalle ouvert contenant $x_0$. On dit que $f$ est \cle{dérivable en} $x_0$ si le taux d'accroissement
\[ \tau(h) = \frac{f(x_0 + h) - f(x_0)}{h} \]
admet une \textbf{limite finie} lorsque $h$ tend vers $0$. Cette limite est alors appelée \cle{nombre dérivé} de $f$ en $x_0$ et est notée $f'(x_0)$ :
\[ f'(x_0) = \lim_{h \to 0} \frac{f(x_0 + h) - f(x_0)}{h}. \]
\end{definition}

\begin{attention}[La limite doit être FINIE]
C'est l'erreur la plus fréquente au Probatoire : conclure « $f'(x_0) = +\infty$ ». Un nombre dérivé est un \textbf{nombre réel}. Si le taux d'accroissement tend vers $+\infty$ ou $-\infty$, ou si les limites à gauche et à droite sont différentes, alors $f$ \textbf{n'est pas dérivable} en $x_0$. Graphiquement, une limite infinie correspond à une tangente \emph{verticale}, qui n'a pas de coefficient directeur.
\end{attention}

\begin{methode}[Calculer $f'(x_0)$ par la définition]
Pour déterminer le nombre dérivé de $f$ en $x_0$ :
\begin{enumerate}
\item \textbf{Former le taux} : écrire $\tau(h) = \dfrac{f(x_0+h) - f(x_0)}{h}$ en développant $f(x_0+h)$.
\item \textbf{Lever l'indétermination} : pour $h = 0$, on obtient $\frac{0}{0}$, forme indéterminée. Simplifier le numérateur pour \textbf{factoriser par $h$}, puis simplifier la fraction par $h$ (licite car $h \neq 0$).
\item \textbf{Passer à la limite} : faire tendre $h$ vers $0$ dans l'expression simplifiée. Si la limite est un réel $\ell$, alors $f$ est dérivable en $x_0$ et $f'(x_0) = \ell$.
\end{enumerate}
\end{methode}

\begin{experience}[Exemple fondateur : $f(x) = x^2$ en $x_0 = 1$]
Déroulons la méthode en détail, comme au tableau.
\begin{enumerate}
\item On calcule $f(1+h) = (1+h)^2 = 1 + 2h + h^2$ et $f(1) = 1$.
\item Le taux d'accroissement est, pour $h \neq 0$ :
\[ \tau(h) = \frac{(1 + 2h + h^2) - 1}{h} = \frac{2h + h^2}{h} = \frac{h(2 + h)}{h} = 2 + h. \]
L'indétermination est levée : le facteur $h$ s'est simplifié.
\item On passe à la limite : $\displaystyle\lim_{h \to 0} (2 + h) = 2$, limite \textbf{finie}.
\end{enumerate}
\textbf{Conclusion :} $f$ est dérivable en $1$ et $f'(1) = 2$. La tangente à la parabole au point $A(1\,;\,1)$ a pour coefficient directeur $2$ — c'est exactement ce que montre la figure précédente.
\end{experience}

\begin{exoresolu}[Nombre dérivé de $f(x) = x^2 - 3x$ en $x_0 = 2$]
Soit $f$ la fonction définie sur $\mathbb{R}$ par $f(x) = x^2 - 3x$. Montrer que $f$ est dérivable en $2$ et déterminer $f'(2)$.
\tcblower
\textbf{Étape 1 : former le taux.} On a $f(2) = 4 - 6 = -2$ et :
\[ f(2+h) = (2+h)^2 - 3(2+h) = 4 + 4h + h^2 - 6 - 3h = h^2 + h - 2. \]
\textbf{Étape 2 : simplifier.} Pour $h \neq 0$ :
\[ \tau(h) = \frac{f(2+h) - f(2)}{h} = \frac{(h^2 + h - 2) - (-2)}{h} = \frac{h^2 + h}{h} = \frac{h(h+1)}{h} = h + 1. \]
\textbf{Étape 3 : limite.} $\displaystyle\lim_{h \to 0} (h + 1) = 1$, limite finie.

\textbf{Conclusion :} $f$ est dérivable en $2$ et $\boxed{f'(2) = 1}$.
\end{exoresolu}

\courssub{Des fonctions non dérivables : deux contre-exemples à connaître}

Toutes les fonctions ne sont pas dérivables partout. Deux situations classiques doivent être reconnues immédiatement.

\textbf{Premier cas : un « point anguleux ».} Considérons la fonction valeur absolue $f(x) = |x|$ en $x_0 = 0$. Pour $h > 0$ :
\[ \tau(h) = \frac{|h| - 0}{h} = \frac{h}{h} = 1 \qquad \text{donc} \qquad \lim_{\substack{h \to 0 \\ h > 0}} \tau(h) = 1. \]
Pour $h < 0$ :
\[ \tau(h) = \frac{|h|}{h} = \frac{-h}{h} = -1 \qquad \text{donc} \qquad \lim_{\substack{h \to 0 \\ h < 0}} \tau(h) = -1. \]
Les limites à gauche ($-1$) et à droite ($1$) sont \textbf{différentes} : le taux n'a pas de limite en $0$. La fonction $x \mapsto |x|$ n'est donc \textbf{pas dérivable en} $0$. Graphiquement, la courbe présente en $0$ deux \cle{demi-tangentes} distinctes, de pentes $-1$ et $1$ : c'est un point anguleux.

\begin{popfigure}
\begin{center}
\begin{tikzpicture}
\begin{axis}[
    width=0.8\linewidth, height=6.5cm,
    axis lines=middle,
    xlabel={$x$}, ylabel={$y$},
    xmin=-3.2, xmax=3.2, ymin=-0.6, ymax=3.2,
    xtick={-2,-1,1,2}, ytick={1,2,3},
    legend style={at={(0.03,0.97)}, anchor=north west, font=\small}
]
\addplot[popcurve, domain=-3:3, samples=200] {abs(x)};
\addlegendentry{$y = |x|$}
\addplot[poporange, thick, dashed, domain=0:3, samples=2] {x};
\addlegendentry{demi-tangente droite : pente $1$}
\addplot[poppurple, thick, dashed, domain=-3:0, samples=2] {-x};
\addlegendentry{demi-tangente gauche : pente $-1$}
\addplot[only marks, mark=*, mark size=2pt, popdark] coordinates {(0,0)};
\node[popdark, anchor=north] at (axis cs:0,-0.15) {$O$};
\end{axis}
\end{tikzpicture}
\end{center}
\legende{La courbe de $x \mapsto |x|$ présente un point anguleux en $0$ : deux demi-tangentes de pentes différentes, donc pas de nombre dérivé en $0$.}
\end{popfigure}

\textbf{Deuxième cas : une tangente verticale.} Considérons $f(x) = \sqrt{x}$ en $x_0 = 0$ (on ne peut prendre que $h > 0$, car $f$ n'est définie que sur $[0\,;\,+\infty[$) :
\[ \tau(h) = \frac{\sqrt{h} - 0}{h} = \frac{\sqrt{h}}{h} = \frac{1}{\sqrt{h}} \xrightarrow[h \to 0^+]{} +\infty. \]
La limite est \textbf{infinie} : $f$ n'est pas dérivable en $0$. Graphiquement, les sécantes se redressent de plus en plus et la courbe admet en $0$ une tangente \emph{verticale} — qui n'a pas de coefficient directeur.

\begin{aretenir}[Les trois situations possibles en $x_0$]
\begin{itemize}
\item $\tau(h)$ tend vers un réel $\ell$ : $f$ est dérivable en $x_0$, $f'(x_0) = \ell$, la courbe admet une tangente de coefficient directeur $\ell$.
\item $\tau(h)$ tend vers l'infini : $f$ n'est pas dérivable en $x_0$ ; la courbe admet une tangente (ou demi-tangente) verticale.
\item Les limites à gauche et à droite diffèrent : $f$ n'est pas dérivable en $x_0$ ; la courbe présente un point anguleux.
\end{itemize}
\end{aretenir}

\courssub{Lien avec la physique : la vitesse instantanée}

Revenons à notre moto-taxi de Bafoussam. Notons $d(t)$ la distance (en km) qu'il a parcourue à l'instant $t$ (en heures). Entre les instants $t_0$ et $t_0 + h$, sa \cle{vitesse moyenne} est :
\[ v_{\text{moy}}(h) = \frac{d(t_0 + h) - d(t_0)}{h}. \]
C'est exactement un taux d'accroissement ! Pour obtenir la vitesse \emph{à l'instant} $t_0$, on resserre l'intervalle de temps autour de $t_0$, c'est-à-dire qu'on fait tendre $h$ vers $0$ :
\[ v(t_0) = \lim_{h \to 0} \frac{d(t_0 + h) - d(t_0)}{h} = d'(t_0). \]
La \cle{vitesse instantanée} est le nombre dérivé de la fonction position. C'est précisément pour résoudre ce problème de physique que Newton et Leibniz ont inventé le calcul différentiel au \textsc{xvii}\textsuperscript{e} siècle.

\begin{exemplebox}[Vitesse instantanée d'un moto-taxi]
La distance parcourue par le moto-taxi est modélisée par $d(t) = 40t^2$ (en km, $t$ en heures) pendant sa phase d'accélération. Sa vitesse à l'instant $t_0 = 0{,}1$ h est :
\[ \tau(h) = \frac{40(0{,}1 + h)^2 - 40(0{,}1)^2}{h} = \frac{40(0{,}01 + 0{,}2h + h^2) - 0{,}4}{h} = \frac{8h + 40h^2}{h} = 8 + 40h, \]
qui tend vers $8$ quand $h \to 0$. La vitesse instantanée à $t = 0{,}1$ h est donc $d'(0{,}1) = 8$ km/h. Le moto-taxi accélère encore !
\end{exemplebox}

\begin{savaistu}[Lagrange, Leibniz et les notations]
La notation $f'(x_0)$ que nous utilisons est due au mathématicien franco-italien \textbf{Joseph-Louis Lagrange} (1736--1813), qui parlait de « fonction dérivée » — d'où le mot \emph{dérivée}. Son contemporain \textbf{Gottfried Wilhelm Leibniz} (1646--1716) écrivait quant à lui $\dfrac{\mathrm{d}f}{\mathrm{d}x}$, une notation qui suggère le quotient de deux variations infiniment petites et qui est encore massivement utilisée en physique : la vitesse s'y note $v = \dfrac{\mathrm{d}x}{\mathrm{d}t}$. Newton, de son côté, notait $\dot{x}$. Trois notations, une même idée géniale : mesurer l'instantané.
\end{savaistu}

\begin{exercice}[À toi de jouer]
\begin{enumerate}
\item Soit $f(x) = 3x^2 + x$. En utilisant la définition, montrer que $f$ est dérivable en $x_0 = 1$ et calculer $f'(1)$.
\item Soit $g(x) = \dfrac{1}{x}$. Calculer le taux d'accroissement de $g$ entre $2$ et $2 + h$, puis en déduire $g'(2)$.
\item La fonction $h(x) = |x - 1|$ est-elle dérivable en $1$ ? Justifier en calculant les limites à gauche et à droite du taux d'accroissement.
\end{enumerate}
\end{exercice}

\begin{corrige}[Exercice n°1]
\textbf{1.} $f(1) = 4$ et $f(1+h) = 3(1+h)^2 + (1+h) = 3 + 6h + 3h^2 + 1 + h = 3h^2 + 7h + 4$. Donc :
\[ \tau(h) = \frac{3h^2 + 7h + 4 - 4}{h} = \frac{h(3h + 7)}{h} = 3h + 7 \xrightarrow[h \to 0]{} 7. \]
$f$ est dérivable en $1$ et $f'(1) = 7$.

\textbf{2.} $g(2) = \dfrac{1}{2}$ et $g(2+h) = \dfrac{1}{2+h}$. Pour lever l'indétermination, on réduit au même dénominateur :
\[ \tau(h) = \frac{\dfrac{1}{2+h} - \dfrac{1}{2}}{h} = \frac{\dfrac{2 - (2+h)}{2(2+h)}}{h} = \frac{-h}{2h(2+h)} = \frac{-1}{2(2+h)} \xrightarrow[h \to 0]{} -\frac{1}{4}. \]
Donc $g'(2) = -\dfrac{1}{4}$.

\textbf{3.} Pour $h > 0$ : $\tau(h) = \dfrac{|h|}{h} = 1$. Pour $h < 0$ : $\tau(h) = \dfrac{|h|}{h} = -1$. Les limites à gauche et à droite sont différentes, donc $h$ n'est \textbf{pas dérivable en} $1$ (point anguleux, comme $|x|$ en $0$, translaté).
\end{corrige}

\begin{aretenir}[Bilan de la section]
\begin{itemize}
\item Le \cle{taux d'accroissement} $\tau(h) = \dfrac{f(x_0+h) - f(x_0)}{h}$ est le coefficient directeur de la sécante $(AM)$.
\item $f$ est \cle{dérivable en} $x_0$ si $\tau(h)$ admet une \textbf{limite finie} quand $h \to 0$ ; cette limite est le \cle{nombre dérivé} $f'(x_0)$.
\item Pour le calculer : former le taux, factoriser par $h$ pour lever l'indétermination $\frac{0}{0}$, passer à la limite.
\item Limite infinie $\Rightarrow$ non-dérivabilité (tangente verticale) ; limites gauche/droite différentes $\Rightarrow$ non-dérivabilité (point anguleux).
\item En physique, la vitesse instantanée est le nombre dérivé de la position.
\end{itemize}
\end{aretenir}

\coursec{Interprétation graphique et tangente}

Dans la section précédente, nous avons défini le \cle{nombre dérivé} $f'(x_0)$ comme la limite du taux d'accroissement $\dfrac{f(x_0+h)-f(x_0)}{h}$ lorsque $h$ tend vers $0$. Cette définition est analytique : elle repose sur des calculs de limites. Mais que \emph{voit-on} sur la courbe ? C'est toute la question de cette section : donner un \cle{sens géométrique} au nombre dérivé. Nous allons découvrir que $f'(x_0)$ n'est rien d'autre que le \cle{coefficient directeur de la tangente} à la courbe au point d'abscisse $x_0$. Ce résultat, simple en apparence, est l'un des plus féconds des mathématiques : c'est lui qui permettra, dans les leçons suivantes, d'étudier les variations des fonctions et de résoudre des problèmes d'optimisation concrets (coût minimal d'une entreprise de Douala, recette maximale d'un commerçant du marché Mokolo, etc.).

\courssub{Position limite des sécantes : naissance de la tangente}

\textbf{Intuition.} Place-toi devant la courbe représentative $\mathcal{C}_f$ d'une fonction $f$. Fixe un point $A$ de la courbe, d'abscisse $x_0$, et choisis un second point $M$ de la courbe, d'abscisse $x_0 + h$ avec $h \neq 0$. La droite $(AM)$ est une \cle{sécante} : elle coupe la courbe en deux points. Or le coefficient directeur de cette sécante est exactement
\[
\frac{y_M - y_A}{x_M - x_A} = \frac{f(x_0+h) - f(x_0)}{h},
\]
c'est-à-dire le \cle{taux d'accroissement} étudié dans la section précédente ! Faisons maintenant glisser le point $M$ le long de la courbe, de plus en plus près de $A$ (c'est-à-dire $h \to 0$). La sécante $(AM)$ pivote autour de $A$ et semble se rapprocher d'une \emph{position limite} : une droite qui \og épouse \fg{} la courbe au point $A$. Cette droite limite, c'est la \cle{tangente}.

\begin{definition}[Tangente à une courbe]
Soit $f$ une fonction définie sur un intervalle ouvert contenant $x_0$, de courbe représentative $\mathcal{C}_f$, et $A$ le point de $\mathcal{C}_f$ d'abscisse $x_0$. Si, lorsque $M$ tend vers $A$ le long de la courbe, la sécante $(AM)$ admet une position limite, cette droite limite est appelée \cle{tangente} à $\mathcal{C}_f$ au point $A$.
\end{definition}

\begin{savaistu}[D'où vient le mot « tangente » ?]
Le mot vient du latin \emph{tangere}, « toucher ». Les Grecs (Euclide, Archimède) connaissaient déjà la tangente au cercle : la droite qui le touche en un seul point et qui est perpendiculaire au rayon. Mais pour une courbe quelconque, il a fallu attendre le XVII\textsuperscript{e} siècle et les travaux de Fermat, Descartes, puis Newton et Leibniz, pour définir la tangente comme \emph{position limite des sécantes} — exactement l'idée que nous venons de développer. C'est la naissance du calcul différentiel.
\end{savaistu}

\courssub{Le théorème fondamental : dérivabilité et existence de la tangente}

Le raisonnement géométrique précédent se transforme en un théorème précis, qui fait le pont entre le calcul (le nombre dérivé) et la géométrie (la tangente).

\begin{propriete}[Dérivabilité et tangente]
Soit $f$ une fonction définie sur un intervalle ouvert contenant $x_0$, et $A\bigl(x_0 \,;\, f(x_0)\bigr)$ le point de sa courbe $\mathcal{C}_f$.
\begin{itemize}
\item Si $f$ est \cle{dérivable en} $x_0$, alors $\mathcal{C}_f$ admet au point $A$ une tangente \textbf{non verticale}, dont le coefficient directeur est $f'(x_0)$.
\item Réciproquement (admise), si $\mathcal{C}_f$ admet en $A$ une tangente non verticale, alors $f$ est dérivable en $x_0$ et le coefficient directeur de cette tangente est $f'(x_0)$.
\end{itemize}
Autrement dit : \textbf{$f$ est dérivable en $x_0$ si et seulement si $\mathcal{C}_f$ admet en $A$ une tangente non verticale.}
\end{propriete}

\begin{experience}[Pourquoi le coefficient directeur de la tangente est-il $f'(x_0)$ ?]
Suivons pas à pas le passage à la limite.
\begin{enumerate}
\item Pour $h \neq 0$, la sécante $(AM)$ joint $A\bigl(x_0\,;\,f(x_0)\bigr)$ et $M\bigl(x_0+h\,;\,f(x_0+h)\bigr)$. Son coefficient directeur est $a_h = \dfrac{f(x_0+h)-f(x_0)}{h}$.
\item Par hypothèse, $f$ est dérivable en $x_0$, donc $\lim\limits_{h \to 0} a_h = f'(x_0)$ : les coefficients directeurs des sécantes ont une limite finie.
\item Or une droite passant par le point fixe $A$ est entièrement déterminée par son coefficient directeur. Les sécantes $(AM)$ approchent donc la droite passant par $A$ et de coefficient directeur $f'(x_0)$.
\item Cette droite est, par définition, la position limite des sécantes : c'est la tangente en $A$. Comme son coefficient directeur $f'(x_0)$ est un nombre fini, elle n'est pas verticale. $\blacksquare$
\end{enumerate}
\end{experience}

\begin{attention}[La tangente n'est pas « la droite qui ne touche qu'un point »]
Contrairement au cas du cercle, la tangente à une courbe quelconque \textbf{peut recouper la courbe} ailleurs, et même au point de contact lui-même ! Exemple célèbre : la courbe de $f(x) = x^3$ au point $O(0\,;\,0)$. On a $f'(0) = 0$, donc la tangente en $O$ est l'axe des abscisses $y = 0$... qui \emph{traverse} la courbe en $O$ (la cubique passe de $y<0$ à $y>0$). La bonne définition est celle de la \emph{position limite des sécantes}, jamais celle du « point unique de contact ».
\end{attention}

\courssub{Équation cartésienne de la tangente}

Connaître le coefficient directeur de la tangente ne suffit pas : pour tracer la droite ou faire des calculs, il faut son \cle{équation}. C'est l'objet du résultat central de cette section, à connaître par cœur pour le Probatoire.

\begin{aretenir}[Équation de la tangente]
Soit $f$ une fonction dérivable en $x_0$. La tangente à $\mathcal{C}_f$ au point $A\bigl(x_0\,;\,f(x_0)\bigr)$ a pour équation :
\[
\boxed{\;y = f'(x_0)\,(x - x_0) + f(x_0)\;}
\]
\end{aretenir}

\begin{experience}[Justification guidée de la formule]
Démontrons cette formule — c'est une démonstration courte et formatrice, souvent demandée sous forme de question de cours.
\begin{enumerate}
\item La tangente $T$ passe par $A\bigl(x_0\,;\,f(x_0)\bigr)$ et a pour coefficient directeur $f'(x_0)$ (théorème précédent).
\item Toute droite non verticale admet une équation de la forme $y = mx + p$. Ici $m = f'(x_0)$, donc $T : y = f'(x_0)\,x + p$ pour un certain réel $p$ à déterminer.
\item Le point $A$ appartient à $T$, donc ses coordonnées vérifient l'équation : $f(x_0) = f'(x_0)\,x_0 + p$, d'où $p = f(x_0) - f'(x_0)\,x_0$.
\item En remplaçant : $y = f'(x_0)\,x + f(x_0) - f'(x_0)\,x_0 = f'(x_0)(x - x_0) + f(x_0)$. $\blacksquare$
\end{enumerate}
Retiens la structure : \emph{ordonnée du point} $+$ \emph{pente} $\times$ \emph{écart à} $x_0$.
\end{experience}

\begin{methode}[Écrire l'équation d'une tangente en 3 étapes]
Pour déterminer l'équation de la tangente à $\mathcal{C}_f$ au point d'abscisse $x_0$ :
\begin{enumerate}
\item \textbf{Calculer} $f(x_0)$ : c'est l'ordonnée du point de contact $A$.
\item \textbf{Calculer} $f'(x_0)$ : c'est le coefficient directeur de la tangente (par la définition avec le taux d'accroissement, ou avec les formules de dérivation vues en section 3).
\item \textbf{Appliquer} la formule $y = f'(x_0)(x - x_0) + f(x_0)$, puis développer et réduire si nécessaire.
\end{enumerate}
Vérification rapide : la tangente obtenue doit passer par $A$ — en remplaçant $x$ par $x_0$, on doit retrouver $y = f(x_0)$.
\end{methode}

\begin{exemplebox}[Tangente à l'hyperbole $y = \dfrac{1}{x}$ en $x_0 = 2$]
Soit $f(x) = \dfrac{1}{x}$, dérivable en $x_0 = 2$. Appliquons la méthode.
\begin{enumerate}
\item $f(2) = \dfrac{1}{2}$ : le point de contact est $A\left(2\,;\,\dfrac{1}{2}\right)$.
\item Le taux d'accroissement en $2$ est $\dfrac{f(2+h)-f(2)}{h} = \dfrac{\frac{1}{2+h}-\frac{1}{2}}{h} = \dfrac{\frac{2-(2+h)}{2(2+h)}}{h} = \dfrac{-1}{2(2+h)}$, qui tend vers $-\dfrac{1}{4}$ quand $h \to 0$. Donc $f'(2) = -\dfrac{1}{4}$.
\item Équation : $y = -\dfrac{1}{4}(x - 2) + \dfrac{1}{2}$, c'est-à-dire après développement :
\[ y = -\frac{x}{4} + 1. \]
\end{enumerate}
Vérification : pour $x = 2$, $y = -\dfrac{2}{4} + 1 = \dfrac{1}{2} = f(2)$. La tangente passe bien par $A$.
\end{exemplebox}

\begin{attention}[Ne confonds jamais $f(x_0)$ et $f'(x_0)$ !]
L'erreur la plus fréquente au Probatoire est d'écrire $y = f(x_0)(x - x_0) + f'(x_0)$, c'est-à-dire d'intervertir les deux nombres. Retiens bien les rôles :
\begin{itemize}
\item $f(x_0)$ est une \textbf{ordonnée} (la hauteur du point de contact) ;
\item $f'(x_0)$ est une \textbf{pente} (le coefficient directeur de la tangente).
\end{itemize}
Dans la formule $y = f'(x_0)(x - x_0) + f(x_0)$, la pente multiplie $(x - x_0)$, et l'ordonnée s'ajoute à la fin. Autre piège : oublier les parenthèses et écrire $f'(x_0) \cdot x - x_0 + f(x_0)$.
\end{attention}

\courssub{Exemple complet : la tangente à la parabole en $A(1\,;\,1)$}

Traitons en détail l'exemple de référence : la parabole d'équation $y = x^2$ au point d'abscisse $1$. Soit $f(x) = x^2$.
\begin{enumerate}
\item $f(1) = 1^2 = 1$, donc $A(1\,;\,1)$.
\item $f'(1) = 2 \times 1 = 2$ (nous l'avons établi dans la section 1 par le taux d'accroissement ; la formule $(x^2)' = 2x$ sera justifiée en section 3).
\item Tangente : $y = 2(x - 1) + 1$, soit $\boxed{y = 2x - 1}$.
\end{enumerate}

\begin{popfigure}
\begin{center}
\begin{tikzpicture}
\begin{axis}[
  width=0.9\linewidth, height=7.5cm,
  axis lines=middle, axis line style={-{Stealth}},
  xlabel={$x$}, ylabel={$y$},
  xmin=-1.6, xmax=2.6, ymin=-1.6, ymax=4.6,
  xtick={-1,1,2}, ytick={-1,1,2,3,4},
  grid=both, grid style={popline!40, very thin},
  tick label style={font=\small},
  clip=true
]
\addplot[popcurve,domain=-1.5:2.1,samples=200]{x^2};
\addplot[poporange,very thick,domain=-0.6:2.6,samples=2]{2*x-1};
\addplot[only marks,mark=*,mark size=2pt,popdark] coordinates {(1,1)};
\node[popdark,above left] at (axis cs:1,1) {$A(1\,;\,1)$};
% triangle de pente
\draw[popgreen,thick] (axis cs:1,1) -- (axis cs:2,1) node[midway,below,font=\small]{$1$};
\draw[popgreen,thick] (axis cs:2,1) -- (axis cs:2,3) node[midway,right,font=\small]{$f'(1)=2$};
\addplot[only marks,mark=*,mark size=1.5pt,popgreen] coordinates {(2,3)};
\node[poporange,font=\small] at (axis cs:2.35,2.6) {$T : y=2x-1$};
\node[popblue,font=\small] at (axis cs:-1.1,3.2) {$y=x^2$};
\end{axis}
\end{tikzpicture}
\end{center}
\legende{La parabole $y = x^2$ et sa tangente en $A(1\,;\,1)$. Le triangle vert montre la pente : pour un déplacement de $1$ en abscisse, la tangente monte de $f'(1) = 2$ en ordonnée.}
\end{popfigure}

\courssub{Construction géométrique de la tangente}

La formule $y = f'(x_0)(x - x_0) + f(x_0)$ fournit un mode d'emploi pour \textbf{tracer} la tangente sans calculer d'équation, directement sur le graphique.

\begin{methode}[Construire la tangente connaissant $f'(x_0)$]
\begin{enumerate}
\item \textbf{Placer} le point de contact $A\bigl(x_0\,;\,f(x_0)\bigr)$ sur la courbe.
\item \textbf{Reporter la pente} : à partir de $A$, effectuer un déplacement de $1$ unité en abscisse (vers la droite), puis de $f'(x_0)$ unités en ordonnée (vers le haut si $f'(x_0) > 0$, vers le bas si $f'(x_0) < 0$). On obtient un second point $B$.
\item \textbf{Tracer} la droite $(AB)$ : c'est la tangente cherchée.
\end{enumerate}
Sur la figure précédente, c'est exactement le triangle vert : depuis $A(1\,;\,1)$, on avance de $1$ puis on monte de $2$, ce qui donne le point $(2\,;\,3)$, et la droite joignant ces deux points est la tangente $y = 2x - 1$.
\end{methode}

\courssub{Lecture graphique d'un nombre dérivé}

Le théorème se lit aussi \emph{à l'envers}, et c'est une compétence exigible : si l'on te donne une courbe avec sa tangente déjà tracée, tu peux \textbf{lire} la valeur de $f'(x_0)$ sans aucun calcul de limite.

\begin{methode}[Lire graphiquement $f'(x_0)$]
\begin{enumerate}
\item Repérer le point $A$ d'abscisse $x_0$ et la tangente $T$ tracée en ce point.
\item Choisir deux points de $T$ aux coordonnées lisibles sur le quadrillage, par exemple $A$ et un point $B$.
\item Calculer le coefficient directeur : $f'(x_0) = \dfrac{y_B - y_A}{x_B - x_A}$.
\end{enumerate}
On lit de même $f(x_0)$ : c'est simplement l'ordonnée de $A$.
\end{methode}

\begin{popfigure}
\begin{center}
\begin{tikzpicture}[scale=1.05]
% repère
\draw[-{Stealth},popdark] (-0.6,0) -- (6.4,0) node[below left]{$x$};
\draw[-{Stealth},popdark] (0,-0.8) -- (0,4.4) node[below left]{$y$};
\foreach \x in {1,2,3,4,5,6} \draw[popline!50,very thin] (\x,-0.8) -- (\x,4.4);
\foreach \y in {1,2,3,4} \draw[popline!50,very thin] (-0.6,\y) -- (6.4,\y);
\foreach \x in {1,2,3,4,5} \node[below,font=\small] at (\x,0) {$\x$};
\foreach \y in {1,2,3} \node[left,font=\small] at (0,\y) {$\y$};
\node[below left,font=\small] at (0,0) {$O$};
% courbe
\draw[popblue,very thick] plot[smooth] coordinates {(0.4,0.6) (1,1.5) (2,2.6) (3,3.0) (4,2.7) (5,2.0) (6,1.5)};
\node[popblue,font=\small] at (5.3,2.6) {$\mathcal{C}_f$};
% tangente en x0=2 : pente 1/2, passe par (2;2.6) approx -> on choisit points exacts (2,2.5) et (4,3.5)? pente 0.5
\draw[poporange,very thick] (0.5,1.75) -- (5.5,4.25);
\node[poporange,font=\small,rotate=22] at (5.0,3.65) {$T$};
% points de lecture
\fill[popdark] (2,2.5) circle (1.8pt) node[below left,font=\small]{$A(2\,;\,2{,}5)$};
\fill[popgreen] (4,3.5) circle (1.8pt) node[above right,font=\small]{$B(4\,;\,3{,}5)$};
% pointillés
\draw[dashed,popmuted] (2,0) -- (2,2.5);
\draw[dashed,popmuted] (0,2.5) -- (2,2.5);
\node[below,font=\small,popmuted] at (2,-0.35) {$x_0$};
\end{tikzpicture}
\end{center}
\legende{Lecture graphique : la tangente $T$ en $A$ passe par $A(2\,;\,2{,}5)$ et $B(4\,;\,3{,}5)$. On lit $f(2) = 2{,}5$ et $f'(2) = \dfrac{3{,}5 - 2{,}5}{4 - 2} = \dfrac{1}{2}$.}
\end{popfigure}

\begin{exoresolu}[Lecture graphique complète]
Sur la figure ci-dessus, on donne la courbe $\mathcal{C}_f$ et sa tangente $T$ au point $A$ d'abscisse $2$. Déterminer graphiquement $f(2)$ et $f'(2)$, puis donner l'équation de $T$.
\tcblower
\textbf{Solution.} Le point $A$ a pour coordonnées $(2\,;\,2{,}5)$, donc $f(2) = 2{,}5$. La tangente $T$ passe aussi par $B(4\,;\,3{,}5)$, donc son coefficient directeur est
\[
f'(2) = \frac{y_B - y_A}{x_B - x_A} = \frac{3{,}5 - 2{,}5}{4 - 2} = \frac{1}{2}.
\]
L'équation de la tangente est donc $y = \dfrac{1}{2}(x - 2) + 2{,}5$, c'est-à-dire $y = \dfrac{x}{2} + 1{,}5$.
\end{exoresolu}

\courssub{Cas de non-dérivabilité : demi-tangentes et tangente verticale}

Que se passe-t-il lorsque $f$ n'est \textbf{pas} dérivable en $x_0$ ? Le théorème (dérivabilité $\Longleftrightarrow$ tangente non verticale) nous dit que la courbe ne peut pas avoir de tangente non verticale en $A$. Deux situations typiques se présentent.

\textbf{1. Le point anguleux : deux demi-tangentes.} Le taux d'accroissement peut admettre une limite \emph{à gauche} et une limite \emph{à droite} différentes. On note alors
\[
f'_g(x_0) = \lim_{h \to 0^-} \frac{f(x_0+h)-f(x_0)}{h}
\qquad \text{et} \qquad
f'_d(x_0) = \lim_{h \to 0^+} \frac{f(x_0+h)-f(x_0)}{h},
\]
appelés \cle{nombre dérivé à gauche} et \cle{nombre dérivé à droite}. La courbe présente en $A$ deux \cle{demi-tangentes} de pentes différentes : elle forme un « angle ». Exemple classique : $f(x) = |x|$ en $x_0 = 0$. Pour $h > 0$, $\dfrac{|h|}{h} = 1$ ; pour $h < 0$, $\dfrac{|h|}{h} = -1$. Donc $f'_d(0) = 1$ et $f'_g(0) = -1$ : les limites diffèrent, $f$ n'est pas dérivable en $0$, et la courbe (un « V ») présente un point anguleux à l'origine.

\textbf{2. La tangente verticale.} Le taux d'accroissement peut aussi avoir une limite \emph{infinie} : $\lim\limits_{h \to 0} \dfrac{f(x_0+h)-f(x_0)}{h} = +\infty$ (ou $-\infty$). Les sécantes deviennent de plus en plus « raides » et leur position limite est une droite \textbf{verticale} d'équation $x = x_0$. La courbe admet bien une tangente en $A$, mais verticale, donc $f$ n'est pas dérivable en $x_0$. Exemple : $f(x) = \sqrt{x}$ en $0$ : pour $h > 0$, $\dfrac{\sqrt{h}}{h} = \dfrac{1}{\sqrt{h}} \to +\infty$.

\begin{aretenir}[À retenir sur la non-dérivabilité]
\begin{itemize}
\item $f$ dérivable en $x_0$ $\Longleftrightarrow$ $\mathcal{C}_f$ admet en $A$ une tangente \textbf{non verticale}, de pente $f'(x_0)$.
\item Deux demi-tangentes de pentes différentes $\Rightarrow$ \textbf{point anguleux} $\Rightarrow$ non-dérivabilité (ex. $|x|$ en $0$).
\item Taux d'accroissement de limite infinie $\Rightarrow$ \textbf{tangente verticale} $x = x_0$ $\Rightarrow$ non-dérivabilité (ex. $\sqrt{x}$ en $0$).
\end{itemize}
Graphiquement, une fonction dérivable en $x_0$ a une courbe « lisse » en $A$ : ni angle, ni verticalité.
\end{aretenir}

\begin{exercice}[Pour t'entraîner]
\begin{enumerate}
\item Soit $g(x) = x^2 - 3x$. Déterminer l'équation de la tangente à $\mathcal{C}_g$ au point d'abscisse $x_0 = 2$.
\item Sur un graphique, la tangente à $\mathcal{C}_f$ au point $A(1\,;\,3)$ passe par le point $B(3\,;\,-1)$. Lire $f(1)$ et $f'(1)$.
\item Montrer, en revenant aux taux d'accroissement à gauche et à droite, que $f(x) = |x - 1|$ n'est pas dérivable en $1$. Interpréter graphiquement.
\end{enumerate}
\end{exercice}

\begin{corrige}[Exercice]
\begin{enumerate}
\item $g(2) = 4 - 6 = -2$. Le taux d'accroissement en $2$ est $\dfrac{g(2+h)-g(2)}{h} = \dfrac{(2+h)^2 - 3(2+h) + 2}{h} = \dfrac{h^2 + h}{h} = h + 1 \to 1$, donc $g'(2) = 1$. Tangente : $y = 1 \times (x - 2) - 2$, soit $y = x - 4$.
\item $f(1) = 3$ (ordonnée de $A$) et $f'(1) = \dfrac{-1 - 3}{3 - 1} = -2$.
\item Pour $h > 0$ : $\dfrac{|h|}{h} = 1 \to 1 = f'_d(1)$. Pour $h < 0$ : $\dfrac{|h|}{h} = -1 \to -1 = f'_g(1)$. Les limites étant différentes, $f$ n'est pas dérivable en $1$ : la courbe présente un point anguleux en $A(1\,;\,0)$, avec deux demi-tangentes de pentes $-1$ et $1$.
\end{enumerate}
\end{corrige}

En résumé, le nombre dérivé $f'(x_0)$ possède désormais deux visages : un visage \emph{analytique} (limite du taux d'accroissement) et un visage \emph{géométrique} (pente de la tangente). C'est cette double lecture qui fait toute la puissance de la dérivation. Dans la section suivante, nous quitterons le point $x_0$ fixé pour définir la \cle{fonction dérivée} $f'$, et nous établirons les formules permettant de dériver rapidement toutes les fonctions usuelles — fini les calculs de limites à chaque fois !

\coursec{Fonction dérivée et dérivées des fonctions élémentaires}

Dans la section précédente, tu as appris à calculer le nombre dérivé $f'(x_0)$ d'une fonction en \emph{un} point $x_0$, et tu as vu qu'il donne le coefficient directeur de la tangente à la courbe au point d'abscisse $x_0$. Mais refaire le calcul de la limite du taux d'accroissement à chaque point serait fastidieux, comme si un chauffeur de moto-taxi devait recalculer sa vitesse à chaque mètre parcouru ! L'idée géniale de Newton et de Leibniz est la suivante : au lieu de travailler point par point, on calcule une fois pour toutes une \emph{nouvelle fonction}, notée $f'$, qui donne \emph{directement} le nombre dérivé en n'importe quel point. C'est l'objet de cette section : définir la \cle{fonction dérivée}, puis établir les dérivées des fonctions élémentaires que tu utiliseras toute l'année.

\courssub{Fonction dérivée sur un intervalle}

\textbf{De la notion ponctuelle à la notion globale.}\par
Nous savons ce qu'est une fonction dérivable en un réel $x_0$ : le taux d'accroissement $\dfrac{f(x_0+h)-f(x_0)}{h}$ admet une limite finie quand $h$ tend vers $0$. Si cette propriété est vraie en \emph{tous} les points d'un intervalle, on peut fabriquer une nouvelle fonction.

\begin{definition}[Fonction dérivable sur un intervalle, fonction dérivée]
Soit $f$ une fonction définie sur un intervalle $I$.
\begin{itemize}
\item On dit que $f$ est \cle{dérivable sur} $I$ si $f$ est dérivable en tout réel $x$ de $I$.
\item Dans ce cas, la fonction qui à tout $x$ de $I$ associe le nombre dérivé $f'(x)$ est appelée \cle{fonction dérivée} de $f$ et se note $f'$ :
\[ f' : x \longmapsto f'(x). \]
\end{itemize}
\end{definition}

\begin{exemplebox}[De la limite à la fonction]
Pour $f(x)=x^2$, nous avons montré en section 1 que le taux d'accroissement en un réel $x$ quelconque vaut
\[ \frac{(x+h)^2-x^2}{h} = \frac{2xh+h^2}{h} = 2x+h \xrightarrow[h\to 0]{} 2x. \]
Ce calcul est valable pour \emph{tout} réel $x$ : $f$ est donc dérivable sur $\mathbb{R}$ et sa fonction dérivée est $f'(x)=2x$. Désormais, pour connaître le coefficient directeur de la tangente en $x=3$, inutile de refaire une limite : $f'(3)=2\times 3=6$.
\end{exemplebox}

\begin{attention}[Dérivable en un point ou sur un intervalle ?]
Ne confonds pas les deux formulations. Dire « $f$ est dérivable en $x_0$ » est une propriété \emph{locale} (en un seul point). Dire « $f$ est dérivable sur $I$ » exige la dérivabilité en \emph{chaque} point de $I$. La fonction $x\mapsto |x|$ est dérivable en tout $x\neq 0$ mais pas en $0$ : elle est donc dérivable sur $]0\,;\,+\infty[$ (et sur $]-\infty\,;\,0[$) mais pas sur $\mathbb{R}$.
\end{attention}

\courssub{Dérivées des fonctions constantes et de la fonction identité}

\textbf{Intuition.} Une fonction constante ne varie jamais : son « taux de variation » est nul partout. La fonction identité $x\mapsto x$ augmente exactement aussi vite que $x$ : sa pente vaut $1$ partout. Formalisons cela.

\begin{propriete}[Dérivée d'une fonction constante]
Toute fonction constante $f : x\mapsto k$ (où $k\in\mathbb{R}$) est dérivable sur $\mathbb{R}$ et
\[ (k)' = 0. \]
\end{propriete}

\textbf{Justification.} Pour tout $x$ et tout $h\neq 0$ : $\dfrac{f(x+h)-f(x)}{h}=\dfrac{k-k}{h}=0$, qui tend vers $0$ quand $h$ tend vers $0$. \hfill$\blacksquare$

\begin{propriete}[Dérivée de la fonction identité]
La fonction identité $f : x\mapsto x$ est dérivable sur $\mathbb{R}$ et
\[ (x)' = 1. \]
\end{propriete}

\textbf{Justification.} $\dfrac{(x+h)-x}{h}=\dfrac{h}{h}=1 \xrightarrow[h\to 0]{} 1$. \hfill$\blacksquare$

\courssub{Dérivée des fonctions puissances : $(x^n)' = n\,x^{n-1}$}

\textbf{Intuition.} Nous avons déjà rencontré $(x^2)' = 2x$. Le cas $n=1$ donne $(x)'=1=1\cdot x^0$. On devine une régularité : l'exposant « descend » en facteur et diminue de $1$. Vérifions cette conjecture sur $n=2$ et $n=3$.

\begin{experience}[Démonstrations pour $n=2$ et $n=3$]
\textbf{Cas $n=2$ :} $f(x)=x^2$. Pour $h\neq 0$ :
\[ \frac{f(x+h)-f(x)}{h} = \frac{x^2+2xh+h^2-x^2}{h} = 2x+h \xrightarrow[h\to 0]{} 2x. \]
Donc $(x^2)'=2x$ sur $\mathbb{R}$.

\textbf{Cas $n=3$ :} $f(x)=x^3$. Rappelons l'identité $(x+h)^3 = x^3+3x^2h+3xh^2+h^3$. Alors :
\begin{align*}
\frac{f(x+h)-f(x)}{h} &= \frac{3x^2h+3xh^2+h^3}{h} = 3x^2+3xh+h^2.
\end{align*}
Quand $h\to 0$, les termes $3xh$ et $h^2$ tendent vers $0$, donc le taux tend vers $3x^2$. Donc $(x^3)'=3x^2$ sur $\mathbb{R}$.
\end{experience}

\begin{propriete}[Dérivée des puissances — admise dans le cas général]
Pour tout entier $n\geq 1$, la fonction $f : x\mapsto x^n$ est dérivable sur $\mathbb{R}$ et
\[ (x^n)' = n\,x^{\,n-1}. \]
La démonstration générale (par récurrence ou avec la formule du binôme) sera vue en Terminale ; les cas $n=2$ et $n=3$ ci-dessus sont exigibles.
\end{propriete}

\begin{exemplebox}[Application immédiate]
\begin{itemize}
\item $(x^4)' = 4x^3$ ; en $x=2$ : $f'(2)=4\times 2^3 = 32$.
\item $(x^7)' = 7x^6$.
\item $(x)' = 1\cdot x^0 = 1$ : la formule redonne bien la dérivée de l'identité.
\end{itemize}
\end{exemplebox}

\begin{attention}[L'erreur classique du facteur $n$]
Beaucoup d'élèves écrivent $(x^n)' = x^{n-1}$ en \emph{oubliant le facteur} $n$. C'est faux ! Retiens le mécanisme : \textbf{l'exposant descend devant, puis diminue de 1}. Ainsi $(x^5)' = 5x^4$ et non $x^4$. Vérifie toujours sur le cas témoin $(x^2)'=2x$ : si ta formule ne redonne pas le $2$, elle est fausse.
\end{attention}

\begin{savaistu}[Et si $n$ n'est pas un entier naturel ?]
La formule $(x^n)' = n\,x^{n-1}$ est bien plus puissante qu'il n'y paraît : elle reste valable pour $n$ entier \emph{négatif} (par exemple $(x^{-1})' = -x^{-2} = -\dfrac{1}{x^2}$, ce qui redonne la dérivée de l'inverse !) et même pour $n$ \emph{réel quelconque} sur $]0\,;\,+\infty[$ (par exemple $(x^{1/2})'=\tfrac{1}{2}x^{-1/2}$, c'est-à-dire la dérivée de $\sqrt{x}$). Tu démontreras cette généralisation en Terminale grâce à la fonction exponentielle. Deux résultats que tu vas établir ci-dessous ne sont donc que des cas particuliers d'une même loi universelle.
\end{savaistu}

\courssub{Dérivée de la fonction inverse}

\textbf{Intuition.} La fonction $x\mapsto \dfrac{1}{x}$ décroît sur $]0\,;\,+\infty[$ : sa dérivée doit donc être \emph{négative}. De plus, près de $0$, la courbe est presque verticale (pente très grande en valeur absolue), tandis qu'elle s'aplatit quand $x$ grandit. Une dérivée en $-\dfrac{1}{x^2}$ correspond exactement à ce comportement.

\begin{propriete}[Dérivée de l'inverse — démonstration exigible]
La fonction $f : x\mapsto \dfrac{1}{x}$ est dérivable sur $\mathbb{R}^*$ (c'est-à-dire sur $]-\infty\,;\,0[$ et sur $]0\,;\,+\infty[$) et
\[ \left(\frac{1}{x}\right)' = -\frac{1}{x^2}. \]
\end{propriete}

\begin{experience}[Démonstration par le taux d'accroissement]
Soit $x\neq 0$ fixé et $h\neq 0$ assez petit pour que $x+h\neq 0$.
\begin{align*}
\frac{f(x+h)-f(x)}{h} &= \frac{\dfrac{1}{x+h}-\dfrac{1}{x}}{h} = \frac{\dfrac{x-(x+h)}{x(x+h)}}{h} = \frac{\dfrac{-h}{x(x+h)}}{h} = \frac{-1}{x(x+h)}.
\end{align*}
Quand $h\to 0$, le dénominateur $x(x+h)$ tend vers $x^2\neq 0$, donc
\[ \frac{f(x+h)-f(x)}{h} \xrightarrow[h\to 0]{} -\frac{1}{x^2}. \]
La limite existe et est finie : $f$ est dérivable en tout $x\neq 0$ et $f'(x)=-\dfrac{1}{x^2}$. \hfill$\blacksquare$
\end{experience}

\begin{exemplebox}[Pente de l'hyperbole]
La tangente à l'hyperbole $y=\dfrac{1}{x}$ au point d'abscisse $x=2$ a pour coefficient directeur $f'(2)=-\dfrac{1}{4}$. En $x=\dfrac{1}{10}$, la pente vaut $-100$ : la courbe plonge presque verticalement, ce qui confirme notre intuition de départ.
\end{exemplebox}

\courssub{Dérivée de la fonction racine carrée}

\textbf{Intuition.} La fonction $x\mapsto\sqrt{x}$ n'existe que sur $[0\,;\,+\infty[$. Elle croît de plus en plus lentement : sa dérivée doit être positive et décroissante. Le résultat $\dfrac{1}{2\sqrt{x}}$ possède exactement ces propriétés — et il « explose » en $0$, signe que la courbe y admet une tangente verticale.

\begin{propriete}[Dérivée de la racine carrée — démonstration exigible]
La fonction $f : x\mapsto \sqrt{x}$ est dérivable sur $]0\,;\,+\infty[$ et
\[ (\sqrt{x})' = \frac{1}{2\sqrt{x}}. \]
Elle n'est \emph{pas} dérivable en $0$.
\end{propriete}

\begin{experience}[Démonstration par la quantité conjuguée]
Soit $x>0$ fixé et $h\neq 0$ tel que $x+h>0$. Le taux d'accroissement fait apparaître une différence de racines : on multiplie numérateur et dénominateur par la \cle{quantité conjuguée} $\sqrt{x+h}+\sqrt{x}$.
\begin{align*}
\frac{\sqrt{x+h}-\sqrt{x}}{h} &= \frac{(\sqrt{x+h}-\sqrt{x})(\sqrt{x+h}+\sqrt{x})}{h\,(\sqrt{x+h}+\sqrt{x})} = \frac{(x+h)-x}{h\,(\sqrt{x+h}+\sqrt{x})} = \frac{1}{\sqrt{x+h}+\sqrt{x}}.
\end{align*}
Quand $h\to 0$, $\sqrt{x+h}\to\sqrt{x}$, donc le taux tend vers $\dfrac{1}{2\sqrt{x}}$, qui est fini car $x>0$. Ainsi $f'(x)=\dfrac{1}{2\sqrt{x}}$.

\textbf{En $0$ :} pour $h>0$, le taux vaut $\dfrac{\sqrt{h}-0}{h}=\dfrac{1}{\sqrt{h}}$, qui tend vers $+\infty$ quand $h\to 0^+$. La limite n'est pas finie : $f$ n'est pas dérivable en $0$ (la courbe y admet une demi-tangente verticale). \hfill$\blacksquare$
\end{experience}

\begin{methode}[Toujours vérifier l'ensemble de dérivabilité AVANT de dériver]
Avant d'appliquer une formule de dérivation, suis ces étapes :
\begin{enumerate}
\item Détermine l'\textbf{ensemble de définition} de $f$ (ex. : $\sqrt{x}$ exige $x\geq 0$ ; $\dfrac{1}{x}$ exige $x\neq 0$).
\item Identifie les \textbf{points frontières ou problématiques} où la fonction risque de ne pas être dérivable (dénominateur nul, racine en $0$, valeur absolue en $0$).
\item Conclus sur l'\textbf{ensemble de dérivabilité} : $\sqrt{x}$ se dérive sur $]0\,;\,+\infty[$ et \emph{non} en $0$ ; $\dfrac{1}{x}$ se dérive sur $\mathbb{R}^*$ ; $|x|$ se dérive sur $\mathbb{R}^*$ mais pas en $0$.
\item Applique seulement ensuite la formule de dérivation, valable sur cet ensemble.
\end{enumerate}
\end{methode}

\begin{attention}[Deux fonctions non dérivables en 0 à connaître]
\begin{itemize}
\item $x\mapsto |x|$ n'est pas dérivable en $0$ : taux à droite $=1$, taux à gauche $=-1$ (point anguleux).
\item $x\mapsto \sqrt{x}$ n'est pas dérivable en $0$ : le taux tend vers $+\infty$ (tangente verticale).
\end{itemize}
Écrire « $(\sqrt{x})' = \dfrac{1}{2\sqrt{x}}$ sur $[0\,;\,+\infty[$ » est une erreur : la formule n'a de sens que sur $]0\,;\,+\infty[$.
\end{attention}

\courssub{Tableau récapitulatif des dérivées usuelles}

\begin{aretenir}[Les dérivées à connaître par cœur]
\begin{center}
\begin{tabular}{|c|c|c|}
\hline
\rowcolor{popblueL} \textbf{Fonction } $f(x)$ & \textbf{Dérivée } $f'(x)$ & \textbf{Ensemble de dérivabilité} \\
\hline
$k$ (constante) & $0$ & $\mathbb{R}$ \\
\hline
$x$ & $1$ & $\mathbb{R}$ \\
\hline
$x^n$ \ ($n\in\mathbb{N},\ n\geq 1$) & $n\,x^{\,n-1}$ & $\mathbb{R}$ \\
\hline
$\dfrac{1}{x}$ & $-\dfrac{1}{x^2}$ & $\mathbb{R}^*$ \\
\hline
$\sqrt{x}$ & $\dfrac{1}{2\sqrt{x}}$ & $]0\,;\,+\infty[$ \\
\hline
\end{tabular}
\end{center}
\end{aretenir}

\courssub{Lecture graphique du signe de la dérivée}

La figure suivante superpose la parabole $f(x)=x^2$ et la droite $f'(x)=2x$. Observe : quand $x<0$, la dérivée est négative et la parabole descend ; quand $x>0$, la dérivée est positive et la parabole monte ; en $x=0$, $f'(0)=0$ et la tangente est horizontale. La dérivée est le « tableau de bord » des variations de la fonction — lien que nous exploiterons pleinement dans la suite du cours.

\begin{popfigure}
\begin{center}
\begin{tikzpicture}
\begin{axis}[
  width=0.85\linewidth, height=7cm,
  axis lines=middle, xlabel={$x$}, ylabel={$y$},
  xmin=-2.6, xmax=2.6, ymin=-3.5, ymax=5,
  xtick={-2,-1,1,2}, ytick={-2,2,4},
  grid=both, grid style={popline!40},
  legend style={at={(0.02,0.98)}, anchor=north west, draw=popline, fill=white}
]
\addplot[popcurve, domain=-2.2:2.2, samples=200]{x^2};
\addlegendentry{$f(x)=x^2$}
\addplot[poporange, very thick, domain=-1.7:1.7, samples=2]{2*x};
\addlegendentry{$f'(x)=2x$}
\addplot[mark=*, popdark, only marks] coordinates {(0,0)};
\end{axis}
\end{tikzpicture}
\end{center}
\legende{La courbe de $f(x)=x^2$ (bleu) et sa fonction dérivée $f'(x)=2x$ (orange) : le signe de $f'$ indique le sens de variation de $f$.}
\end{popfigure}

\begin{exoresolu}[Calculer et exploiter une dérivée élémentaire]
Soit $f$ la fonction définie par $f(x)=x^3$.
\begin{enumerate}
\item Justifier que $f$ est dérivable sur $\mathbb{R}$ et donner $f'(x)$.
\item Déterminer une équation de la tangente à la courbe de $f$ au point d'abscisse $x_0=2$.
\end{enumerate}
\tcblower
\textbf{Solution.}
\begin{enumerate}
\item $f$ est une fonction puissance ($n=3$), dérivable sur $\mathbb{R}$ d'après le cours, et
\[ f'(x) = 3x^{3-1} = 3x^2. \]
\item On calcule $f(2)=2^3=8$ et $f'(2)=3\times 2^2=12$. L'équation de la tangente en $x_0=2$ est
\[ y = f'(2)(x-2)+f(2) = 12(x-2)+8 = 12x-16. \]
La tangente au point $(2\,;\,8)$ a donc pour équation $y=12x-16$.
\end{enumerate}
\end{exoresolu}

\begin{exercice}[À toi de jouer]
\begin{enumerate}
\item Déterminer la fonction dérivée de chacune des fonctions suivantes, en précisant l'ensemble de dérivabilité : $f(x)=x^5$ ; $g(x)=\sqrt{x}$ ; $h(x)=\dfrac{1}{x}$.
\item Soit $f(x)=\sqrt{x}$. Déterminer une équation de la tangente à sa courbe au point d'abscisse $4$.
\item La fonction $x\mapsto\sqrt{x}$ est-elle dérivable en $0$ ? Justifier.
\end{enumerate}
\end{exercice}

\begin{corrige}[Exercice]
\begin{enumerate}
\item $f'(x)=5x^4$ sur $\mathbb{R}$ ; $g'(x)=\dfrac{1}{2\sqrt{x}}$ sur $]0\,;\,+\infty[$ ; $h'(x)=-\dfrac{1}{x^2}$ sur $\mathbb{R}^*$.
\item $f(4)=2$ et $f'(4)=\dfrac{1}{2\sqrt{4}}=\dfrac{1}{4}$. Tangente : $y=\dfrac{1}{4}(x-4)+2=\dfrac{1}{4}x+1$.
\item Non : pour $h>0$, le taux d'accroissement en $0$ vaut $\dfrac{\sqrt{h}}{h}=\dfrac{1}{\sqrt{h}}\to +\infty$ ; la limite n'est pas finie, donc $\sqrt{x}$ n'est pas dérivable en $0$.
\end{enumerate}
\end{corrige}

Tu disposes maintenant de toutes les dérivées élémentaires. Mais les fonctions rencontrées en exercice sont rarement « pures » : ce sont des sommes, des produits, des quotients comme $3x^2-5x+2$ ou $\dfrac{x}{x+1}$. La section suivante te donnera les règles opératoires pour dériver ces fonctions construites à partir des fonctions usuelles.

\coursec{Dérivées et opérations sur les fonctions}

Dans la section précédente, nous avons appris à calculer la dérivée des fonctions « de base » : puissances, racine carrée, fonctions trigonométriques, exponentielles et logarithmes. Mais dans la vie réelle — que ce soit pour modéliser le coût de production d'une usine à Douala, la trajectoire d'un projectile ou le débit d'une connexion MTN — les fonctions sont souvent des \textbf{combinaisons} de ces fonctions élémentaires : sommes, produits, quotients.

Heureusement, il n'est pas nécessaire de repartir de la définition du taux d'accroissement à chaque fois. Des règles de dérivation, vraies « boîtes à outils », permettent de dériver ces combinaisons mécaniquement. C'est l'objet de cette section.

\courssub{Règles de linéarité : somme et multiplication par un scalaire}

Commençons par les opérations les plus simples. Soient $f$ et $g$ deux fonctions dérivables sur un intervalle ouvert $]a ; b[$ et $k \in \mathbb{R}$.

\begin{experience}[Démonstration guidée : dérivée d'une somme]
On veut déterminer $(f+g)'(x_0)$ pour $x_0 \in ]a;b[$.
\begin{enumerate}
    \item Écris le taux d'accroissement de $f+g$ en $x_0$ avec un incrément $h \neq 0$ tel que $x_0+h \in ]a;b[$.
    \item Utilise la définition $(f+g)(x) = f(x) + g(x)$ pour séparer la fraction en deux termes.
    \item Fais tendre $h$ vers $0$. Que peux-tu conclure grâce aux hypothèses de dérivabilité de $f$ et $g$ ?
\end{enumerate}
\tcblower
\textbf{Solution guidée :}
\begin{align*}
\tau_{f+g}(x_0, h) &= \frac{(f+g)(x_0+h) - (f+g)(x_0)}{h} \\
&= \frac{f(x_0+h) + g(x_0+h) - f(x_0) - g(x_0)}{h} \\
&= \frac{f(x_0+h) - f(x_0)}{h} + \frac{g(x_0+h) - g(x_0)}{h} \\
&= \tau_f(x_0, h) + \tau_g(x_0, h).
\end{align*}
En faisant tendre $h \to 0$, $\tau_f(x_0, h) \to f'(x_0)$ et $\tau_g(x_0, h) \to g'(x_0)$. La limite de la somme est la somme des limites. Donc $(f+g)'(x_0) = f'(x_0) + g'(x_0)$.
\end{experience}

\begin{experience}[Démonstration guidée : dérivée de $kf$]
Soit $k \in \mathbb{R}$. Montre de même que $(kf)'(x_0) = k f'(x_0)$.
\tcblower
$\tau_{kf}(x_0, h) = \frac{kf(x_0+h) - kf(x_0)}{h} = k \frac{f(x_0+h) - f(x_0)}{h} = k \tau_f(x_0, h)$. En passant à la limite : $(kf)'(x_0) = k f'(x_0)$.
\end{experience}

\begin{propriete}[Linéarité de la dérivation]
Soient $f$ et $g$ dérivables sur $]a;b[$ et $k \in \mathbb{R}$. Les fonctions $f+g$ et $kf$ sont dérivables sur $]a;b[$ et :
\begin{itemize}
    \item $(f+g)' = f' + g'$ \quad (la dérivée de la somme est la somme des dérivées)
    \item $(kf)' = k f'$ \quad (la constante « sort » de la dérivée)
\end{itemize}
\emph{Remarque :} Cette propriété s'étend par récurrence à toute combinaison linéaire : $( \sum k_i f_i )' = \sum k_i f_i'$.
\end{propriete}

\begin{exemplebox}[Dérivation d'un polynôme]
Soit $P(x) = 3x^3 - 5x^2 + 7x - 2$ définie sur $\mathbb{R}$.
$P$ est somme de termes $k x^n$. On dérive terme à terme :
\[ P'(x) = 3 \cdot 3x^2 - 5 \cdot 2x + 7 \cdot 1 - 0 = 9x^2 - 10x + 7. \]
\end{exemplebox}

\courssub{Dérivée d'un produit : la règle du « premier fois dérivée du second plus second fois dérivée du premier »}

C'est ici que se cachent les pièges les plus fréquents. \textbf{La dérivée d'un produit n'est PAS le produit des dérivées.}

\begin{experience}[Démonstration guidée : la formule du produit]
Soient $f, g$ dérivables en $x_0$. On cherche $(fg)'(x_0)$.
\begin{enumerate}
    \item Écris le taux d'accroissement $\tau_{fg}(x_0, h)$.
    \item \textbf{Astuce clé :} Ajoute et retranche $f(x_0+h)g(x_0)$ au numérateur (c'est le terme « mixte »).
    \item Factorise par $f(x_0+h)$ dans le premier groupe et par $g(x_0)$ dans le second.
    \item Fais tendre $h \to 0$. Utilise le fait que $f$ est continue en $x_0$ (car dérivable) donc $f(x_0+h) \to f(x_0)$.
\end{enumerate}
\tcblower
\textbf{Solution guidée :}
\begin{align*}
\tau_{fg}(x_0, h) &= \frac{f(x_0+h)g(x_0+h) - f(x_0)g(x_0)}{h} \\
&= \frac{f(x_0+h)g(x_0+h) - f(x_0+h)g(x_0) + f(x_0+h)g(x_0) - f(x_0)g(x_0)}{h} \\
&= \frac{f(x_0+h)[g(x_0+h)-g(x_0)] + g(x_0)[f(x_0+h)-f(x_0)]}{h} \\
&= f(x_0+h) \tau_g(x_0, h) + g(x_0) \tau_f(x_0, h).
\end{align*}
Quand $h \to 0$ : $f(x_0+h) \to f(x_0)$, $\tau_g \to g'(x_0)$, $\tau_f \to f'(x_0)$.
Donc $(fg)'(x_0) = f(x_0)g'(x_0) + g(x_0)f'(x_0)$.
\end{experience}

\begin{propriete}[Dérivée d'un produit]
Soient $f, g$ dérivables sur $]a;b[$. La fonction $fg$ est dérivable sur $]a;b[$ et :
\[ (fg)' = f'g + fg'. \]
\textbf{Mnémotechnique :} « Dérivée du premier fois le second \textbf{plus} le premier fois dérivée du second. »
\end{propriete}

\begin{attention}[Piège mortel : $(fg)' \neq f'g'$]
Ne confonds jamais !
Exemple : $f(x)=x$, $g(x)=x$. $fg = x^2$. $(fg)' = 2x$. Mais $f'g' = 1 \times 1 = 1 \neq 2x$.
Toujours écrire la formule complète avant de calculer.
\end{attention}

\begin{exemplebox}[Produit : $f(x) = x^2 \sqrt{x}$]
On peut voir $f$ comme produit de $u(x)=x^2$ et $v(x)=\sqrt{x}=x^{1/2}$ sur $]0;+\infty[$.
$u'(x)=2x$, $v'(x)=\frac{1}{2}x^{-1/2} = \frac{1}{2\sqrt{x}}$.
\[ f'(x) = u'v + uv' = 2x \cdot \sqrt{x} + x^2 \cdot \frac{1}{2\sqrt{x}} = 2x\sqrt{x} + \frac{x^2}{2\sqrt{x}}. \]
On simplifie en mettant au même dénominateur $2\sqrt{x}$ :
\[ f'(x) = \frac{4x^2 + x^2}{2\sqrt{x}} = \frac{5x^2}{2\sqrt{x}} = \frac{5}{2}x^{3/2}. \]
\textbf{Vérification :} $f(x)=x^{5/2} \Rightarrow f'(x)=\frac{5}{2}x^{3/2}$. Ça colle !
\end{exemplebox}

\courssub{Dérivée d'un inverse et d'un quotient}

Pour diviser, on dérive $1/g$ puis on utilise la règle du produit avec $f \times (1/g)$.

\begin{experience}[Démonstration guidée : dérivée de l'inverse]
Soit $g$ dérivable en $x_0$ avec $g(x_0) \neq 0$. On cherche $(1/g)'(x_0)$.
\begin{enumerate}
    \item Écris $\tau_{1/g}(x_0, h) = \frac{\frac{1}{g(x_0+h)} - \frac{1}{g(x_0)}}{h}$.
    \item Mets au même dénominateur le numérateur.
    \item Fais apparaître $\tau_g(x_0, h)$.
    \item Passe à la limite.
\end{enumerate}
\tcblower
\begin{align*}
\tau_{1/g}(x_0, h) &= \frac{1}{h} \left( \frac{g(x_0) - g(x_0+h)}{g(x_0+h)g(x_0)} \right) \\
&= -\frac{g(x_0+h)-g(x_0)}{h} \cdot \frac{1}{g(x_0+h)g(x_0)} \\
&= -\tau_g(x_0, h) \cdot \frac{1}{g(x_0+h)g(x_0)}.
\end{align*}
Quand $h \to 0$ : $\tau_g \to g'(x_0)$, $g(x_0+h) \to g(x_0)$ (continuité).
Donc $(1/g)'(x_0) = -\frac{g'(x_0)}{g(x_0)^2}$.
\end{experience}

\begin{propriete}[Dérivée de l'inverse]
Si $g$ est dérivable sur $]a;b[$ et $g(x) \neq 0$ sur $]a;b[$, alors $1/g$ est dérivable sur $]a;b[$ et :
\[ \left( \frac{1}{g} \right)' = -\frac{g'}{g^2}. \]
\end{propriete}

\begin{propriete}[Dérivée d'un quotient — Formule à connaître par cœur]
Soient $f, g$ dérivables sur $]a;b[$ avec $g(x) \neq 0$ sur $]a;b[$. La fonction $f/g$ est dérivable sur $]a;b[$ et :
\[ \left( \frac{f}{g} \right)' = \frac{f'g - fg'}{g^2}. \]
\textbf{Mnémotechnique :} « Dérivée du numérateur fois dénominateur \textbf{moins} numérateur fois dérivée du dénominateur, \textbf{tout sur le dénominateur au carré}. »
\end{propriete}

\begin{methode}[Dériver un quotient en 4 étapes]
Pour dériver $h(x) = \frac{f(x)}{g(x)}$ :
\begin{enumerate}
    \item \textbf{Identifier} $f(x)$ (numérateur) et $g(x)$ (dénominateur). Noter la condition $g(x) \neq 0$.
    \item \textbf{Calculer} $f'(x)$ et $g'(x)$ séparément (brouillon).
    \item \textbf{Appliquer} la formule : $h'(x) = \frac{f'(x)g(x) - f(x)g'(x)}{[g(x)]^2}$.
    \item \textbf{Simplifier} le numérateur (factoriser, développer, réduire si possible). Ne touche pas au dénominateur $g^2$ sauf factorisation commune avec le numérateur.
\end{enumerate}
\end{methode}

\begin{exemplebox}[Quotient : $f(x) = \frac{2x+1}{x-3}$]
Domaine : $D_f = \mathbb{R} \setminus \{3\}$.
\begin{enumerate}
    \item $f(x) = \frac{u(x)}{v(x)}$ avec $u(x)=2x+1$, $v(x)=x-3$. $v(x) \neq 0 \Leftrightarrow x \neq 3$.
    \item $u'(x) = 2$, $v'(x) = 1$.
    \item $f'(x) = \frac{u'v - uv'}{v^2} = \frac{2(x-3) - (2x+1)(1)}{(x-3)^2}$.
    \item Numérateur : $2x - 6 - 2x - 1 = -7$.
    \[ f'(x) = \frac{-7}{(x-3)^2}, \quad x \neq 3. \]
\end{enumerate}
\textbf{Interprétation :} La dérivée est toujours négative (sauf en $x=3$ où elle n'existe pas). La fonction est strictement décroissante sur $]-\infty;3[$ et sur $]3;+\infty[$.
\end{exemplebox}

\begin{popfigure}
\begin{tikzpicture}[scale=0.9, every node/.style={transform shape}]
\begin{axis}[
    width=\linewidth, height=10cm,
    axis lines=middle,
    xlabel=$x$, ylabel=$y$,
    xmin=-2, xmax=8, ymin=-10, ymax=10,
    xtick={-2,-1,0,1,2,3,4,5,6,7},
    ytick={-10,-5,0,5,10},
    restrict y to domain=-12:12,
    samples=400,
    clip=false,
    grid=major, grid style={popline!50},
]
% Asymptote verticale x=3 (dessiné manuellement)
\draw[popmuted, dashed, thick] (3,-12) -- (3,12) node[above right] {$x=3$};
% Asymptote horizontale y=2
\draw[popmuted, dashed, thick] (-2,2) -- (8,2) node[above right] {$y=2$};

% Branche gauche x < 3
\addplot[popblue, thick, domain=-2:2.9] {(2*x+1)/(x-3)};
% Branche droite x > 3
\addplot[popblue, thick, domain=3.1:8] {(2*x+1)/(x-3)};

% Point de tangente x0 = 2
\addplot[poporange, thick, domain=-1:5] {-7/1^2 * (x-2) + (2*2+1)/(2-3)}; % Tangente en x=2 : f(2)=-5, f'(2)=-7
% Point
\fill[poporange] (2,-5) circle (2.5pt) node[above left, black] {$A(2;-5)$};
\end{axis}
\end{tikzpicture}
\legende{Courbe de $f(x)=\frac{2x+1}{x-3}$ (bleu) et sa tangente en $x_0=2$ (orange). La dérivée $f'(2)=-7$ est le coefficient directeur de la tangente.}
\end{popfigure}

\courssub{Arbre de décision : quelle règle utiliser ?}

Face à une fonction composée d'opérations, suis cet arbre pour choisir la bonne formule.

\begin{popfigure}
\begin{tikzpicture}[
    mindmap, grow cyclic, every node/.style=concept, concept color=popblue,
    level 1/.append style={level distance=4.5cm, sibling angle=90},
    level 2/.append style={level distance=3.5cm, sibling angle=45},
    level 3/.append style={level distance=3cm, sibling angle=30},
    font=\small
]
\node[concept, concept color=popdark] {Fonction $h$ à dériver}
    child[concept color=poporange] { node[concept, align=center] {Somme / Différence \\ $f \pm g$}
        child[concept color=poporange!70!white] { node[concept] {$(f \pm g)' = f' \pm g'$} }
    }
    child[concept color=popgreen] { node[concept, align=center] {Produit par constante \\ $k f$}
        child[concept color=popgreen!70!white] { node[concept] {$(kf)' = k f'$} }
    }
    child[concept color=poppurple] { node[concept, align=center] {Produit \\ $f \times g$}
        child[concept color=poppurple!70!white] { node[concept] {$(fg)' = f'g + fg'$} }
    }
    child[concept color=poppink] { node[concept, align=center] {Quotient \\ $f / g$}
        child[concept color=poppink!70!white] { node[concept] {$\left(\frac{f}{g}\right)' = \frac{f'g - fg'}{g^2}$} }
    }
    child[concept color=popgold] { node[concept, align=center] {Inverse \\ $1/g$}
        child[concept color=popgold!70!white] { node[concept] {$(1/g)' = -g'/g^2$} }
    };
\end{tikzpicture}
\legende{Arbre de décision pour choisir la règle de dérivation. Commence par identifier l'opération \textbf{principale} (la dernière effectuée si on calcule la valeur de la fonction).}
\end{popfigure}

\courssub{Enchaînement : dériver puis trouver une tangente}

C'est l'application type du Probatoire : on te donne une fonction, tu la dérives, puis tu donnes l'équation de la tangente en un point.

\begin{exoresolu}[Examen type Probatoire]
Soit la fonction $f$ définie sur $]-\infty; 1[ \cup ]1; +\infty[$ par $f(x) = \frac{x^2 + 1}{x - 1}$.
\begin{enumerate}
    \item Déterminer $f'(x)$.
    \item Écrire une équation cartésienne de la tangente $\mathcal{T}$ à la courbe $\mathcal{C}_f$ au point d'abscisse $x_0 = 0$.
    \item Construire (schématiser) cette tangente.
\end{enumerate}
\textbf{Solution détaillée :}
\begin{enumerate}
    \item \textbf{Dérivée :} Quotient $u(x)=x^2+1$, $v(x)=x-1$. $u'=2x$, $v'=1$.
    \[ f'(x) = \frac{2x(x-1) - (x^2+1)(1)}{(x-1)^2} = \frac{2x^2 - 2x - x^2 - 1}{(x-1)^2} = \frac{x^2 - 2x - 1}{(x-1)^2}. \]
    \item \textbf{Tangente en $x_0=0$ :}
    $f(0) = \frac{1}{-1} = -1$. Point $A(0; -1)$.
    $f'(0) = \frac{0 - 0 - 1}{(0-1)^2} = -1$. Coefficient directeur $m = -1$.
    Équation : $y - (-1) = -1(x - 0) \iff \boxed{y = -x - 1}$.
    \item \textbf{Construction :} La tangente passe par $A(0;-1)$ et a pour directeur $-1$. Elle coupe l'axe des abscisses en $B(-1;0)$. Tracer la droite $(AB)$.
\end{enumerate}
\end{exoresolu}

\begin{savaistu}[Le calcul différentiel au service de l'économie camerounaise]
Imaginons une PME à Bafoussam qui produit des chaises. Son coût total (en milliers de FCFA) pour $x$ chaises est modélisé par $C(x) = 0,5x^2 + 20x + 500$.
Le coût marginal, c'est-à-dire le coût de production de \emph{la chaise supplémentaire} quand on en a déjà produit $x$, est exactement la dérivée $C'(x) = x + 20$.
Si l'entreprise a produit 100 chaises, la 101\textsuperscript{ème} coûtera environ $C'(100) = 120$ mille FCFA.
La dérivation des polynômes (somme, produit par constante) permet donc aux gestionnaires de prendre des décisions de production optimales !
\end{savaistu}

\begin{aretenir}[Résumé des formules à connaître par cœur]
Pour $f, g$ dérivables sur $I$, $k \in \mathbb{R}$, $g \neq 0$ sur $I$ :
\begin{center}
\begin{tabularx}{\linewidth}{|l|X|}\hline
\textbf{Opération} & \textbf{Dérivée} \\ \hline
Somme $f+g$ & $f' + g'$ \\ \hline
Produit par scalaire $kf$ & $k f'$ \\ \hline
Produit $fg$ & $f'g + fg'$ \\ \hline
Inverse $1/g$ & $-\dfrac{g'}{g^2}$ \\ \hline
Quotient $f/g$ & $\dfrac{f'g - fg'}{g^2}$ \\ \hline
\end{tabularx}
\end{center}
\textbf{Ordre des opérations :} Pour dériver une expression complexe, identifie l'opération \textbf{principale} (la dernière faite pour calculer $f(x)$) et applique la règle correspondante en dérivant les « morceaux » (qui sont eux-mêmes des fonctions plus simples).
\end{aretenir}

\begin{exercice}[Entraînement]
Dériver les fonctions suivantes sur leur domaine de définition :
\begin{enumerate}
    \item $f(x) = (x^2 + 3x - 1)(2x - 5)$
    \item $g(x) = \frac{x^3 - 2x}{x^2 + 1}$
    \item $h(x) = \sqrt{x}(x - 4)$ \quad ($x \ge 0$)
    \item $k(x) = \frac{3}{x^2 + 4}$ \quad (astuce : écrire $3 \times (x^2+4)^{-1}$ ou quotient)
\end{enumerate}
\end{exercice}

\begin{corrige}[Exercice n°1]
\begin{enumerate}
    \item Produit : $u=x^2+3x-1$, $v=2x-5$. $u'=2x+3$, $v'=2$.
    $f' = (2x+3)(2x-5) + (x^2+3x-1)(2) = 4x^2-10x+6x-15 + 2x^2+6x-2 = 6x^2+2x-17$.
    \item Quotient : $u=x^3-2x$, $v=x^2+1$. $u'=3x^2-2$, $v'=2x$.
    $g' = \frac{(3x^2-2)(x^2+1) - (x^3-2x)(2x)}{(x^2+1)^2} = \frac{3x^4+3x^2-2x^2-2 - 2x^4+4x^2}{(x^2+1)^2} = \frac{x^4+5x^2-2}{(x^2+1)^2}$.
    \item Produit : $u=\sqrt{x}=x^{1/2}$, $v=x-4$. $u'=\frac{1}{2\sqrt{x}}$, $v'=1$.
    $h' = \frac{1}{2\sqrt{x}}(x-4) + \sqrt{x} = \frac{x-4}{2\sqrt{x}} + \frac{2x}{2\sqrt{x}} = \frac{3x-4}{2\sqrt{x}}$.
    \item Quotient (ou scalaire $\times$ inverse) : $u=3$, $v=x^2+4$. $u'=0$, $v'=2x$.
    $k' = \frac{0 \cdot (x^2+4) - 3(2x)}{(x^2+4)^2} = \frac{-6x}{(x^2+4)^2}$.
\end{enumerate}
\end{corrige}

\begin{attention}[Check-list avant de rendre ta copie]
\begin{itemize}
    \item Ai-je écrit le domaine de définition (surtout pour les quotients et racines) ?
    \item Ai-je noté explicitement $u, v, u', v'$ avant d'appliquer la formule ?
    \item Ai-je simplifié le numérateur correctement (attention aux signes $-$ devant $fg'$) ?
    \item Pour la tangente : ai-je calculé $f(x_0)$ ET $f'(x_0)$ ? L'équation est-elle sous forme $y = ax + b$ (ou $y - y_0 = a(x-x_0)$) ?
    \item \textbf{Pas de $(fg)' = f'g'$ ni $(f/g)' = f'/g'$ !}
\end{itemize}
\end{attention}

\coursec{Applications et synthèse}

Dans la section précédente, tu as appris à dériver toutes les fonctions construites à partir des fonctions élémentaires : sommes, produits, quotients. Tu disposes maintenant d'un outil puissant, la fonction dérivée $f'$. Mais à quoi sert-elle concrètement ? C'est toute l'ambition de cette dernière section : résoudre des problèmes géométriques sur les tangentes, donner du sens physique et économique au nombre dérivé, puis organiser l'ensemble de la leçon dans une carte mentale. C'est aussi ici que s'ouvre la porte vers la leçon suivante, consacrée aux variations des fonctions.

\courssub{Problème type 1 : les tangentes horizontales}

\textbf{Intuition.} Imagine que tu parcours la route Yaoundé--Bafoussam en moto. Sur une portion vallonnée, il y a des moments précis où tu es exactement au sommet d'une côte, ou au fond d'un creux : à cet instant, la route est localement plate, horizontale. Pour une courbe représentative, ces instants correspondent aux points où la tangente est horizontale.

Or, une droite horizontale a un coefficient directeur nul. Comme le coefficient directeur de la tangente en $x_0$ est $f'(x_0)$, on obtient immédiatement la méthode.

\begin{methode}[Trouver les tangentes horizontales]
Pour déterminer les points de la courbe $\mathcal{C}_f$ où la tangente est horizontale :
\begin{enumerate}
\item Calculer la fonction dérivée $f'$.
\item Résoudre l'équation $f'(x) = 0$ sur l'ensemble de définition de $f$.
\item Pour chaque solution $x_0$, calculer $f(x_0)$ : le point cherché est $A\big(x_0 \,;\, f(x_0)\big)$ et la tangente en $A$ a pour équation $y = f(x_0)$.
\item Vérifier graphiquement, si possible, que la tangente est bien horizontale en ce point.
\end{enumerate}
\end{methode}

\begin{exemplebox}[Tangentes horizontales de la cubique]
Soit $f(x) = x^3 - 3x$, définie et dérivable sur $\mathbb{R}$.
\begin{enumerate}
\item Dérivée : $f'(x) = 3x^2 - 3 = 3(x^2 - 1) = 3(x-1)(x+1)$.
\item Équation : $f'(x) = 0 \iff x = 1$ ou $x = -1$.
\item Ordonnées : $f(1) = 1 - 3 = -2$ et $f(-1) = -1 + 3 = 2$.
\end{enumerate}
La courbe admet donc deux tangentes horizontales : en $A(-1\,;\,2)$, d'équation $y = 2$, et en $B(1\,;\,-2)$, d'équation $y = -2$.
\end{exemplebox}

\begin{popfigure}
\begin{tikzpicture}
\begin{axis}[
  width=\linewidth, height=7cm,
  axis lines=middle, xlabel=$x$, ylabel=$y$,
  xmin=-2.6, xmax=2.6, ymin=-3.2, ymax=3.2,
  xtick={-2,-1,1,2}, ytick={-2,2},
  grid=both, grid style={popline!40},
  legend style={at={(0.02,0.02)}, anchor=south west, font=\small}
]
\addplot[popcurve, domain=-2.2:2.2, samples=200]{x^3 - 3*x};
\addlegendentry{$y = x^3 - 3x$}
\addplot[poporange, thick, domain=-2.4:0.4]{2};
\addlegendentry{tangente $y=2$}
\addplot[popgreen, thick, domain=-0.4:2.4]{-2};
\addlegendentry{tangente $y=-2$}
\addplot[only marks, mark=*, popdark] coordinates {(-1,2) (1,-2)};
\node[popdark, above left] at (axis cs:-1,2) {$A(-1\,;\,2)$};
\node[popdark, below right] at (axis cs:1,-2) {$B(1\,;\,-2)$};
\end{axis}
\end{tikzpicture}
\legende{La cubique $y = x^3 - 3x$ et ses deux tangentes horizontales, aux points où $f'(x)=0$.}
\end{popfigure}

\begin{attention}[Tangente horizontale $\neq$ extremum automatique]
Une tangente horizontale signale un point \emph{candidat} à être un extremum, mais ce n'est pas toujours le cas. Contre-exemple classique~: $f(x) = x^3$. On a $f'(x) = 3x^2$, donc $f'(0) = 0$~: la tangente en $O(0\,;\,0)$ est horizontale (c'est l'axe des abscisses). Pourtant $f$ est strictement croissante sur $\mathbb{R}$~: la courbe \emph{traverse} sa tangente en $O$. On parle de \cle{point d'inflexion}. La condition $f'(x_0)=0$ est nécessaire pour un extremum, jamais suffisante~: il faut en plus un changement de signe de $f'$.
\end{attention}

\courssub{Problème type 2 : tangentes parallèles à une droite donnée}

\textbf{Intuition.} Deux droites sont parallèles si et seulement si elles ont le même coefficient directeur. Chercher les points de $\mathcal{C}_f$ où la tangente est parallèle à une droite $(d)$ de coefficient directeur $m$ revient donc à résoudre $f'(x_0) = m$.

\begin{methode}[Tangentes de direction imposée]
Pour trouver les points où la tangente à $\mathcal{C}_f$ est parallèle à la droite $y = mx + p$ :
\begin{enumerate}
\item Calculer $f'(x)$.
\item Résoudre $f'(x) = m$.
\item Pour chaque solution $x_0$, écrire l'équation de la tangente : $y = f'(x_0)(x - x_0) + f(x_0)$, c'est-à-dire $y = m(x - x_0) + f(x_0)$.
\end{enumerate}
\end{methode}

\begin{exoresolu}[Tangentes parallèles à une droite]
Soit $f(x) = x^2 - 4x + 5$ et $(d)$ la droite d'équation $y = 2x - 7$. Déterminer les points de $\mathcal{C}_f$ où la tangente est parallèle à $(d)$, puis donner les équations de ces tangentes.
\tcblower
\textbf{Solution.} $f$ est dérivable sur $\mathbb{R}$ et $f'(x) = 2x - 4$.

La droite $(d)$ a pour coefficient directeur $m = 2$. On résout~:
\[ f'(x) = 2 \iff 2x - 4 = 2 \iff x = 3. \]
On calcule $f(3) = 9 - 12 + 5 = 2$. Le point cherché est $A(3\,;\,2)$.

La tangente en $A$ a pour équation~:
\[ y = f'(3)(x - 3) + f(3) = 2(x-3) + 2 = 2x - 4. \]
\textbf{Vérification}~: cette tangente a bien le même coefficient directeur $2$ que $(d)$, et elle est distincte de $(d)$ (ordonnées à l'origine $-4 \neq -7$), donc bien parallèle à $(d)$.
\end{exoresolu}

\courssub{Problème type 3 : tangentes passant par un point donné}

\textbf{Intuition.} Le point donné $M(a\,;\,b)$ n'est pas nécessairement sur la courbe~: on cherche les tangentes à $\mathcal{C}_f$ qui, prolongées, passent par $M$. L'inconnue est alors le point de contact $x_0$, et non plus une valeur imposée de la pente.

\begin{methode}[Tangentes issues d'un point]
Pour déterminer les tangentes à $\mathcal{C}_f$ passant par $M(a\,;\,b)$ :
\begin{enumerate}
\item Écrire l'équation de la tangente en un point d'abscisse $x_0$ quelconque~:
\[ y = f'(x_0)(x - x_0) + f(x_0). \]
\item Traduire l'appartenance de $M$ à cette tangente~:
\[ b = f'(x_0)(a - x_0) + f(x_0). \]
\item Résoudre cette équation d'inconnue $x_0$.
\item Pour chaque solution, écrire l'équation de la tangente correspondante.
\end{enumerate}
\end{methode}

\begin{exoresolu}[Tangentes à une parabole passant par l'origine]
Soit $f(x) = x^2 + 1$. Déterminer les tangentes à la parabole $\mathcal{C}_f$ passant par l'origine $O(0\,;\,0)$.
\tcblower
\textbf{Solution.} On a $f'(x) = 2x$. La tangente au point d'abscisse $x_0$ a pour équation~:
\[ y = 2x_0(x - x_0) + x_0^2 + 1 = 2x_0 x - x_0^2 + 1. \]
Elle passe par $O(0\,;\,0)$ si et seulement si~:
\[ 0 = -x_0^2 + 1 \iff x_0^2 = 1 \iff x_0 = 1 \ \text{ou}\ x_0 = -1. \]
Il existe donc \textbf{deux} tangentes issues de $O$~:
\begin{itemize}
\item en $x_0 = 1$~: $y = 2x$, point de contact $(1\,;\,2)$~;
\item en $x_0 = -1$~: $y = -2x$, point de contact $(-1\,;\,2)$.
\end{itemize}
Le résultat est cohérent avec la symétrie de la parabole par rapport à l'axe des ordonnées.
\end{exoresolu}

\begin{attention}[Piège classique]
Ne confonds jamais « tangente \emph{au point} d'abscisse $a$ » (le point est sur la courbe, la pente est $f'(a)$) et « tangente \emph{passant par} le point $M$ » (le point est quelconque, l'inconnue est le point de contact $x_0$). Au Probatoire, cette confusion coûte cher~: lis l'énoncé deux fois.
\end{attention}

\courssub{Interprétations concrètes du nombre dérivé}

\textbf{Vitesse instantanée.} Si $x(t)$ désigne la position d'un mobile à l'instant $t$, le taux d'accroissement $\dfrac{x(t_0+h)-x(t_0)}{h}$ est la \emph{vitesse moyenne} entre $t_0$ et $t_0+h$. En faisant tendre $h$ vers $0$, on obtient la \cle{vitesse instantanée}~:
\[ v(t_0) = x'(t_0). \]
C'est exactement ce qu'affiche le compteur d'un taxi-moto à Douala~: non pas une moyenne sur le trajet, mais la vitesse à l'instant précis.

\begin{exemplebox}[Chute d'un objet]
Un objet lâché d'un pont au-dessus de la Sanaga a pour position (en mètres) $x(t) = 5t^2$ après $t$ secondes de chute. Alors $v(t) = x'(t) = 10t$. Après \SI{2}{s}, sa vitesse instantanée est $v(2) = \SI{20}{m/s}$, soit \SI{72}{km/h}.
\end{exemplebox}

\textbf{Coût marginal.} En économie, si $C(q)$ est le coût total de production de $q$ unités, le \cle{coût marginal} est le coût de production d'une unité supplémentaire. On l'approche par la dérivée~:
\[ C_m(q) \approx C'(q). \]
Une imprimerie de Bafoussam qui produit des affiches de campagne peut ainsi décider si une commande supplémentaire est rentable~: il suffit de comparer le prix de vente unitaire au coût marginal.

\begin{exemplebox}[Coût marginal d'une imprimerie]
Le coût de production de $q$ centaines d'affiches est $C(q) = q^2 + 10q + 100$ (en milliers de FCFA). Alors $C'(q) = 2q + 10$. Pour $q = 20$ (soit 2000 affiches), le coût marginal est $C'(20) = 50$ milliers de FCFA par centaine supplémentaire, soit environ \num{500} FCFA par affiche.
\end{exemplebox}

\courssub{Ouverture : dérivée et sens de variation}

Voici le pont vers la leçon suivante. Observons la cubique $f(x)=x^3-3x$ de la figure précédente~: là où la courbe descend, les tangentes penchent vers le bas ($f'<0$)~; là où elle monte, elles penchent vers le haut ($f'>0$). Ce n'est pas un hasard.

\begin{propriete}[Signe de la dérivée et monotonie (admis)]
Soit $f$ une fonction dérivable sur un intervalle $I$.
\begin{itemize}
\item Si $f'(x) \geq 0$ pour tout $x \in I$, alors $f$ est croissante sur $I$.
\item Si $f'(x) \leq 0$ pour tout $x \in I$, alors $f$ est décroissante sur $I$.
\item Si $f'(x) = 0$ pour tout $x \in I$, alors $f$ est constante sur $I$.
\end{itemize}
Ce théorème, admis en Première, sera l'outil central de l'étude des variations~: il suffira d'étudier le \emph{signe} de $f'$ pour dresser le tableau de variations de $f$.
\end{propriete}

\begin{exemplebox}[Lecture sur la cubique]
Pour $f(x) = x^3 - 3x$, $f'(x) = 3(x-1)(x+1)$ est positive sur $]-\infty\,;\,-1]$ et $[1\,;\,+\infty[$, négative sur $[-1\,;\,1]$. Donc $f$ est croissante, puis décroissante, puis croissante~: les points $A(-1\,;\,2)$ et $B(1\,;\,-2)$ sont bien un maximum local et un minimum local.
\end{exemplebox}

\courssub{Complément (hors programme strict de Première) : dérivée de $f(ax+b)$}

\begin{propriete}[Dérivée d'une fonction composée affine — complément Terminale]
\emph{Hors programme strict de Première, utile pour la Terminale.} Si $f$ est dérivable et $a, b$ sont deux réels, alors la fonction $g : x \mapsto f(ax+b)$ est dérivable et~:
\[ g'(x) = a \times f'(ax+b). \]
Autrement dit, on dérive la fonction « extérieure » et on multiplie par le coefficient $a$ de $x$.
\end{propriete}

\begin{exemplebox}[Application directe]
Soit $g(x) = (3x - 2)^2$. C'est $f(3x-2)$ avec $f(u) = u^2$, donc $f'(u) = 2u$ et~:
\[ g'(x) = 3 \times 2(3x-2) = 6(3x-2) = 18x - 12. \]
Vérification par développement~: $g(x) = 9x^2 - 12x + 4$, donc $g'(x) = 18x - 12$. Les deux méthodes concordent.
\end{exemplebox}

\courssub{Synthèse de la leçon}

\begin{aretenir}[L'essentiel de la leçon]
\begin{itemize}
\item \textbf{Nombre dérivé}~: $f'(x_0) = \lim\limits_{h \to 0} \dfrac{f(x_0+h)-f(x_0)}{h}$, quand cette limite existe et est finie.
\item \textbf{Tangente} en $x_0$~: $y = f'(x_0)(x - x_0) + f(x_0)$.
\item \textbf{Dérivées élémentaires}~: $(x^n)' = nx^{n-1}$, $(\sqrt{x})' = \dfrac{1}{2\sqrt{x}}$, $\left(\dfrac{1}{x}\right)' = -\dfrac{1}{x^2}$, etc.
\item \textbf{Opérations}~: $(f+g)' = f'+g'$ ; $(kf)' = kf'$ ; $(fg)' = f'g + fg'$ ; $\left(\dfrac{1}{g}\right)' = -\dfrac{g'}{g^2}$ ; $\left(\dfrac{f}{g}\right)' = \dfrac{f'g - fg'}{g^2}$.
\item \textbf{Problèmes de tangentes}~: horizontale $\iff f'(x_0)=0$ ; parallèle à $y=mx+p \iff f'(x_0)=m$ ; passant par $M \iff$ équation en $x_0$.
\item \textbf{Interprétations}~: vitesse instantanée, coût marginal.
\item \textbf{Ouverture}~: le signe de $f'$ donne le sens de variation de $f$.
\end{itemize}
\end{aretenir}

\begin{popfigure}
\begin{tikzpicture}[
  root/.style={rectangle, rounded corners, draw=popdark, fill=popblueL, font=\bfseries, align=center, inner sep=6pt},
  node1/.style={rectangle, rounded corners, draw=popblue, fill=popblueL!50, align=center, inner sep=5pt, font=\small},
  node2/.style={rectangle, rounded corners, draw=poporange, fill=poporangeL, align=center, inner sep=5pt, font=\small},
  node3/.style={rectangle, rounded corners, draw=popgreen, fill=popgreenL, align=center, inner sep=5pt, font=\small},
  node4/.style={rectangle, rounded corners, draw=poppurple, fill=poppurpleL, align=center, inner sep=5pt, font=\small},
  lien/.style={thick, popmuted}
]
\node[root, align=center] (R) {Fonctions\\dérivées};
\node[node1, above left=1.6cm and 3.2cm of R, align=center] (A) {Définition\\ $f'(x_0)=\lim\limits_{h\to0}\dfrac{f(x_0+h)-f(x_0)}{h}$};
\node[node2, above right=1.6cm and 3.2cm of R, align=center] (B) {Tangente\\ $y=f'(x_0)(x-x_0)+f(x_0)$};
\node[node3, below left=1.6cm and 3.2cm of R, align=center] (C) {Formules\\ élémentaires\\ $(x^n)'=nx^{n-1}$, \ldots};
\node[node4, below right=1.6cm and 3.2cm of R, align=center] (D) {Opérations\\ $(fg)'=f'g+fg'$\\ $\left(\frac{f}{g}\right)'=\frac{f'g-fg'}{g^2}$};
\node[node2, right=1.2cm of B, align=center] (B2) {\small Problèmes types\\ \small horizontale, parallèle,\\ \small passant par un point};
\node[node4, right=1.2cm of D, align=center] (D2) {\small Applications\\ \small vitesse, coût marginal,\\ \small variations (suite)};
\draw[lien] (R) -- (A);
\draw[lien] (R) -- (B);
\draw[lien] (R) -- (C);
\draw[lien] (R) -- (D);
\draw[lien] (B) -- (B2);
\draw[lien] (D) -- (D2);
\end{tikzpicture}
\legende{Carte mentale de la leçon~: de la définition du nombre dérivé vers les tangentes, les formules et les applications.}
\end{popfigure}

\begin{savaistu}[Deux problèmes, une révolution]
La dérivation est née au XVII\textsuperscript{e} siècle de deux problèmes en apparence distincts~: le \emph{problème des tangentes}, étudié par Descartes et Fermat (comment déterminer la tangente à une courbe en un point ?), et le \emph{problème du mouvement}, étudié par Newton et Leibniz (comment définir la vitesse à un instant précis ?). Newton et Leibniz découvrirent indépendamment que ces deux problèmes sont le même, et que la dérivation est l'opération inverse de l'intégration~: c'est le théorème fondamental de l'analyse. Cette découverte a transformé les mathématiques et permis la physique moderne, de la mécanique céleste aux télécommunications que tu utilises chaque jour avec ton téléphone.
\end{savaistu}

\begin{exercice}[Pour t'entraîner]
Soit $f(x) = x^3 - 3x^2 + 2$ définie sur $\mathbb{R}$.
\begin{enumerate}
\item Calculer $f'(x)$ et déterminer les points de $\mathcal{C}_f$ où la tangente est horizontale.
\item Déterminer les points où la tangente est parallèle à la droite d'équation $y = 9x + 1$.
\item Existe-t-il une tangente à $\mathcal{C}_f$ passant par le point $M(0\,;\,2)$ ?
\end{enumerate}
\end{exercice}

\begin{corrige}[Exercice : tangentes de la cubique]
\textbf{1.} $f'(x) = 3x^2 - 6x = 3x(x-2)$. Alors $f'(x)=0 \iff x=0$ ou $x=2$. Comme $f(0)=2$ et $f(2)=8-12+2=-2$, les tangentes horizontales sont en $(0\,;\,2)$, d'équation $y=2$, et en $(2\,;\,-2)$, d'équation $y=-2$.

\textbf{2.} On résout $f'(x)=9 \iff 3x^2-6x-9=0 \iff x^2-2x-3=0 \iff (x-3)(x+1)=0$, donc $x=3$ ou $x=-1$. On calcule $f(3)=27-27+2=2$ et $f(-1)=-1-3+2=-2$. Tangentes~: en $(3\,;\,2)$, $y=9(x-3)+2=9x-25$~; en $(-1\,;\,-2)$, $y=9(x+1)-2=9x+7$.

\textbf{3.} La tangente en $x_0$ est $y=(3x_0^2-6x_0)(x-x_0)+x_0^3-3x_0^2+2$. Elle passe par $M(0\,;\,2)$ si $2=-x_0(3x_0^2-6x_0)+x_0^3-3x_0^2+2$, soit $-2x_0^3+3x_0^2=0$, c'est-à-dire $x_0^2(3-2x_0)=0$. Donc $x_0=0$ (tangente $y=2$) ou $x_0=\dfrac{3}{2}$ (tangente de pente $f'\left(\tfrac32\right)=\tfrac{27}{4}-9=-\tfrac94$). Il existe donc deux tangentes passant par $M$.
\end{corrige}

\newpage

\coursec{Activité d'intégration}

\courssub{Objectifs de l'activité}

À l'issue de cette activité d'intégration, tu seras capable de mobiliser l'ensemble des savoir-faire de la leçon dans une situation économique authentique. Plus précisément, tu devras :

\begin{itemize}
\item calculer la fonction dérivée d'une fonction polynôme en justifiant chaque formule utilisée (dérivées des fonctions élémentaires, dérivée d'une somme, dérivée de $kf$) ;
\item calculer et interpréter concrètement un nombre dérivé (coût marginal) ;
\item déterminer l'abscisse d'un point où la tangente à une courbe est horizontale et en dégager l'intérêt économique ;
\item déterminer une équation cartésienne d'une tangente et la construire graphiquement ;
\item rédiger un conseil argumenté à destination d'un acteur économique, en t'appuyant sur les mathématiques.
\end{itemize}

\courssub{Situation}

\textbf{Le transporteur de la ligne Douala--Nkongsamba.}

Monsieur Etame est un transporteur bien connu sur la ligne Douala--Nkongsamba. Chaque semaine, son bus effectue plusieurs rotations sur la Nationale n°5, longue d'environ \SI{140}{km}, pour acheminer passagers et marchandises (bananes plantains, manioc, huile de palme) entre la métropole économique et les hautes terres de l'Ouest.

Après des années d'expérience, Monsieur Etame a remarqué que sa consommation de carburant dépend fortement de la vitesse moyenne à laquelle il roule : rouler trop lentement fait durer le trajet et consomme du gasoil « pour rien », tandis que rouler trop vite fait exploser la consommation du moteur. Avec l'aide d'un neveu étudiant, il a modélisé la situation.

On note $x$ la vitesse moyenne du bus, exprimée en km/h, avec $x \in [20\,;\,100]$. Le coût de carburant d'un trajet, exprimé en \textbf{milliers de FCFA}, est donné par la fonction :
\[ C(x) = \frac{x^2}{10} - 8x + 400. \]

Par exemple, à $60$ km/h, le coût du trajet est $C(60) = \dfrac{3600}{10} - 480 + 400 = 360 - 480 + 400 = 280$ milliers de FCFA.

Monsieur Etame se pose trois questions :
\begin{itemize}
\item « Si je passe de $60$ km/h à $61$ km/h, de combien mon coût augmente-t-il environ ? »
\item « Existe-t-il une vitesse idéale, celle qui rend mon coût le plus faible possible ? »
\item « Comment expliquer tout cela simplement à mon chauffeur remplaçant ? »
\end{itemize}

C'est à toi, avec tes connaissances sur les fonctions dérivées, de l'aider à répondre.

\courssub{Consignes}

\begin{enumerate}
\item \textbf{Calcul de la dérivée.} Justifier que $C$ est dérivable sur $[20\,;\,100]$, puis calculer $C'(x)$ en citant précisément, à chaque étape, la formule utilisée (dérivée de $x^2$, dérivée de $kx$, dérivée d'une constante, dérivée d'une somme).
\item \textbf{Coût marginal à 60 km/h.} Calculer $C'(60)$. En économie, on appelle \emph{coût marginal} la valeur approchée du surcoût engendré par une unité supplémentaire ; on admet qu'ici le coût marginal à $60$ km/h est environ $C'(60)$ milliers de FCFA par km/h supplémentaire. Interpréter cette valeur par une phrase simple à destination de Monsieur Etame.
\item \textbf{Recherche de la tangente horizontale.} On rappelle que la tangente à la courbe de $C$ au point d'abscisse $x_0$ est horizontale si et seulement si $C'(x_0) = 0$. Déterminer la vitesse $x_0$ pour laquelle la tangente est horizontale. Calculer $C(x_0)$. Étudier le signe de $C'(x)$ sur $[20\,;\,100]$ et dresser le tableau de variations de $C$. Quelle conjecture économique peux-tu formuler ?
\item \textbf{Tangente au point d'abscisse 60.} Déterminer l'équation cartésienne de la tangente $(T)$ à la courbe de $C$ au point d'abscisse $x_0 = 60$. Sur papier millimétré (ou avec un logiciel), tracer la courbe de $C$ sur $[20\,;\,100]$ et la tangente $(T)$ dans le même repère.
\item \textbf{Conseil au transporteur.} Rédiger en cinq à huit lignes un conseil clair et argumenté à l'attention de Monsieur Etame : quelle vitesse moyenne lui recommandes-tu, pourquoi, et que risque-t-il s'il roule beaucoup plus vite ?
\end{enumerate}

\courssub{Ressources à mobiliser}

\begin{aretenir}[Boîte à outils de la leçon]
\begin{itemize}
\item \textbf{Dérivabilité :} toute fonction polynôme est dérivable sur $\mathbb{R}$, donc sur tout intervalle de $\mathbb{R}$.
\item \textbf{Dérivées élémentaires :} $(x^2)' = 2x$ ; $(x)' = 1$ ; si $k$ est une constante, $(k)' = 0$.
\item \textbf{Opérations :} $(u+v)' = u' + v'$ ; $(ku)' = k\,u'$ pour tout réel $k$.
\item \textbf{Nombre dérivé :} $f'(x_0)$ est le coefficient directeur de la tangente à la courbe de $f$ au point d'abscisse $x_0$.
\item \textbf{Tangente horizontale :} la tangente en $x_0$ est horizontale si et seulement si $f'(x_0) = 0$.
\item \textbf{Équation de la tangente :} $y = f'(x_0)(x - x_0) + f(x_0)$.
\item \textbf{Lien dérivée--variations :} si $f' \leq 0$ sur un intervalle, $f$ y est décroissante ; si $f' \geq 0$, $f$ y est croissante.
\end{itemize}
\end{aretenir}

\courssub{Démarche et corrigé commenté}

\begin{exoresolu}[Le transporteur de la ligne Douala--Nkongsamba — corrigé complet]
Résoudre les cinq consignes de la situation.
\tcblower
\textbf{1. Calcul de la dérivée.}

La fonction $C$ est une fonction polynôme : elle est donc dérivable sur $\mathbb{R}$, et en particulier sur $[20\,;\,100]$. On dérive terme à terme :
\begin{align*}
C(x) &= \frac{x^2}{10} - 8x + 400 \\
C'(x) &= \left(\frac{1}{10} \times x^2\right)' - (8x)' + (400)' & \text{(dérivée d'une somme)}\\
&= \frac{1}{10} \times 2x - 8 \times 1 + 0 & (ku)' = ku',\ (x^2)' = 2x,\ (x)'=1,\ (k)'=0\\
C'(x) &= \frac{x}{5} - 8.
\end{align*}

\textbf{2. Coût marginal à 60 km/h.}
\[ C'(60) = \frac{60}{5} - 8 = 12 - 8 = 4. \]
\emph{Interprétation :} lorsque le bus roule à $60$ km/h, chaque km/h supplémentaire augmente le coût du trajet d'environ $4$ milliers de FCFA, soit environ \num{4000} FCFA. Le nombre dérivé mesure ici la « sensibilité » du coût à la vitesse : c'est le taux de variation instantané du coût par rapport à la vitesse.

\textbf{3. Tangente horizontale et vitesse idéale.}

On résout $C'(x_0) = 0$ :
\[ \frac{x_0}{5} - 8 = 0 \iff \frac{x_0}{5} = 8 \iff x_0 = 40. \]
La tangente est horizontale au point d'abscisse $x_0 = 40$ km/h, et :
\[ C(40) = \frac{1600}{10} - 320 + 400 = 160 - 320 + 400 = 240 \text{ milliers de FCFA.} \]
Signe de $C'(x) = \dfrac{x}{5} - 8$ : c'est une fonction affine de coefficient directeur $\dfrac{1}{5} > 0$, donc strictement croissante, qui s'annule en $x = 40$. Elle est donc strictement négative pour $x < 40$ et strictement positive pour $x > 40$.

Valeurs aux bornes : $C(20) = \dfrac{400}{10} - 160 + 400 = 40 - 160 + 400 = 280$ et $C(100) = \dfrac{10000}{10} - 800 + 400 = 1000 - 800 + 400 = 600$. D'où le tableau de variations :

\begin{center}
\begin{tabular}{|c|ccccc|}\hline
$x$ & $20$ & & $40$ & & $100$ \\ \hline
$C'(x)$ & & $-$ & $0$ & $+$ & \\ \hline
$C$ & $280$ & $\searrow$ & $240$ & $\nearrow$ & $600$ \\ \hline
\end{tabular}
\end{center}

\emph{Conjecture économique :} le coût est minimal à $40$ km/h, où il vaut $240$ milliers de FCFA. Le point de tangente horizontale correspond exactement au minimum de la fonction coût : c'est la \textbf{vitesse économique optimale}. On remarque au passage que $C(20) = C(60) = 280$ : rouler à $20$ km/h coûte aussi cher que rouler à $60$ km/h !

\textbf{4. Équation de la tangente en $x_0 = 60$.}

On applique la formule $y = C'(60)(x - 60) + C(60)$ avec $C(60) = 280$ et $C'(60) = 4$ :
\[ y = 4(x - 60) + 280 = 4x - 240 + 280 \quad\Longrightarrow\quad \boxed{(T) : y = 4x + 40.} \]
Vérification : pour $x = 60$, $y = 240 + 40 = 280 = C(60)$ : la tangente passe bien par le point de contact $(60\,;\,280)$.

\begin{popfigure}
\begin{center}
\begin{tikzpicture}[x=0.055cm, y=0.055cm]
\draw[->,thick] (10,0) -- (105,0) node[below left]{\small $x$ (km/h)};
\draw[->,thick] (10,0) -- (10,640) node[below left]{\small $C(x)$};
\draw[popblue,very thick,domain=20:100,samples=60] plot (\x,{0.1*\x*\x - 8*\x + 400});
\draw[poporange,very thick,domain=25:100] plot (\x,{4*\x + 40});
\fill[popink] (60,280) circle (1.6) node[above right]{\small $(60\,;\,280)$};
\fill[popgreen] (40,240) circle (1.6) node[below right]{\small minimum $(40\,;\,240)$};
\draw[popgreen,dashed] (20,240) -- (60,240);
\end{tikzpicture}
\end{center}
\legende{Courbe de $C$ (bleu), tangente en $x_0 = 60$ (orange) et tangente horizontale au minimum (vert, en pointillés).}
\end{popfigure}

\textbf{5. Conseil à Monsieur Etame.}

« Monsieur Etame, les calculs montrent que votre coût de carburant est le plus bas lorsque vous roulez à une vitesse moyenne de $40$ km/h : le trajet vous coûte alors environ \num{240000} FCFA. En dessous de $40$ km/h, le trajet dure trop longtemps et le coût remonte ; au-dessus, le moteur consomme de plus en plus : à $60$ km/h, chaque km/h en plus vous coûte environ \num{4000} FCFA, et à $100$ km/h le trajet grimpe à \num{600000} FCFA, soit deux fois et demie le minimum ! Ma recommandation : visez une moyenne proche de $40$ km/h, en respectant bien sûr le code de la route et l'état de la Nationale n°5. »
\end{exoresolu}

\courssub{Bilan de l'activité}

Cette activité t'a permis de mettre en évidence plusieurs idées essentielles de la leçon :

\begin{itemize}
\item \textbf{Le nombre dérivé a un sens concret.} $C'(60) = 4$ n'est pas qu'un calcul abstrait : c'est un taux de variation instantané, interprété ici comme un coût marginal en milliers de FCFA par km/h. La dérivation est un outil d'aide à la décision économique.
\item \textbf{La tangente horizontale signale un extremum.} La condition $C'(x_0) = 0$ a permis de repérer la vitesse optimale, confirmée par le tableau de variations : c'est le point de départ de l'étude des extrema, que tu approfondiras dans la leçon sur les applications de la dérivation.
\item \textbf{Les formules de dérivation s'enchaînent logiquement.} Dériver un polynôme n'exige que quelques règles (somme, produit par une constante, dérivées élémentaires), mais chaque étape doit être justifiée — une exigence de rédaction attendue au Probatoire.
\item \textbf{L'équation de la tangente relie calcul algébrique et lecture graphique.} La tangente $(T) : y = 4x + 40$ approche la courbe au voisinage du point de contact : c'est l'idée d'approximation locale d'une courbe par une droite, au cœur de toute l'analyse.
\end{itemize}

\begin{attention}[Pour aller plus loin — vigilance]
Dans ce modèle, le minimum est atteint en $x_0 = 40$, à l'intérieur de l'intervalle $[20\,;\,100]$. Si la solution de $C'(x) = 0$ était tombée \emph{en dehors} de l'intervalle d'étude, il aurait fallu conclure en examinant les valeurs aux bornes $C(20)$ et $C(100)$. Toujours vérifier que le candidat trouvé appartient au domaine du problème ! Autre vigilance : une tangente horizontale ($f'(x_0) = 0$) ne garantit pas à elle seule un extremum — c'est l'étude du \emph{signe} de $f'$ de part et d'autre de $x_0$ qui permet de conclure.
\end{attention}

\newpage

\coursec{Méthodes et exercices résolus}

\courssub{Méthode 1 : Déterminer le nombre dérivé d'une fonction en un réel $x_0$}

\begin{methode}[Calculer $f'(x_0)$ par le taux d'accroissement]
Pour déterminer le nombre dérivé d'une fonction $f$ en un réel $x_0$ de son ensemble de définition :
\begin{enumerate}
\item Vérifier que $x_0$ appartient à l'ensemble de définition de $f$ et que $f$ est définie au voisinage de $x_0$.
\item Former le \cle{taux d'accroissement} (ou taux de variation) de $f$ entre $x_0$ et $x_0+h$ :
\[ \tau(h)=\frac{f(x_0+h)-f(x_0)}{h}, \qquad h\neq 0. \]
\item Simplifier $\tau(h)$ en factorisant par $h$ au numérateur (c'est l'étape clé : on doit pouvoir « éliminer » le $h$ du dénominateur).
\item Calculer la limite de $\tau(h)$ quand $h$ tend vers $0$.
\item Conclure : si cette limite existe et est un nombre réel $\ell$, alors $f$ est \cle{dérivable} en $x_0$ et $f'(x_0)=\ell$. Sinon, $f$ n'est pas dérivable en $x_0$.
\end{enumerate}
\textbf{Réflexe :} une limite du type « $\dfrac{0}{0}$ » est une forme indéterminée : il faut factoriser, développer ou utiliser une quantité conjuguée pour lever l'indétermination.
\end{methode}

\begin{exoresolu}[Nombre dérivé de la fonction carré en $x_0=3$]
Soit $f$ la fonction définie sur $\mathbb{R}$ par $f(x)=x^2$. En utilisant le taux d'accroissement, montrer que $f$ est dérivable en $x_0=3$ et déterminer $f'(3)$.
\tcblower
\textbf{Étape 1.} $f$ est définie sur $\mathbb{R}$, donc en particulier en $3$ et au voisinage de $3$.

\textbf{Étape 2.} Pour $h\neq 0$, formons le taux d'accroissement de $f$ entre $3$ et $3+h$ :
\[ \tau(h)=\frac{f(3+h)-f(3)}{h}=\frac{(3+h)^2-3^2}{h}. \]

\textbf{Étape 3.} Développons le numérateur :
\begin{align*}
\tau(h) &= \frac{9+6h+h^2-9}{h} \\
&= \frac{6h+h^2}{h} \\
&= \frac{h(6+h)}{h} \\
&= 6+h \qquad (\text{car } h\neq 0).
\end{align*}

\textbf{Étape 4.} Passons à la limite :
\[ \lim_{h\to 0}\tau(h)=\lim_{h\to 0}(6+h)=6. \]

\textbf{Étape 5.} La limite existe et est finie, donc $f$ est dérivable en $3$ et :
\[ \boxed{f'(3)=6.} \]
\textbf{Conclusion :} la fonction carré est dérivable en $3$ et son nombre dérivé en ce point vaut $6$.
\end{exoresolu}

\courssub{Méthode 2 : Interpréter graphiquement la dérivabilité en un point}

\begin{methode}[Lire la dérivabilité et le nombre dérivé sur une courbe]
Pour interpréter graphiquement la dérivabilité de $f$ en $x_0$ à partir de sa courbe représentative $\mathcal{C}_f$ :
\begin{enumerate}
\item Repérer le point $A$ de $\mathcal{C}_f$ d'abscisse $x_0$ : ses coordonnées sont $A\big(x_0\,;\,f(x_0)\big)$.
\item Observer l'allure de la courbe au voisinage de $A$ :
\begin{itemize}
\item si la courbe admet en $A$ une \cle{tangente} non verticale, $f$ est dérivable en $x_0$ ;
\item si la courbe présente un \cle{point anguleux} (deux demi-tangentes distinctes) ou une tangente \cle{verticale}, $f$ n'est pas dérivable en $x_0$.
\end{itemize}
\item Lorsque $f$ est dérivable en $x_0$, le nombre dérivé $f'(x_0)$ est le \cle{coefficient directeur} de la tangente en $A$ : on le lit en prenant deux points de la tangente et en calculant
\[ f'(x_0)=\frac{y_B-y_A}{x_B-x_A}. \]
\item Signe de $f'(x_0)$ : tangente « montante » $\Rightarrow f'(x_0)>0$ ; tangente « descendante » $\Rightarrow f'(x_0)<0$ ; tangente horizontale $\Rightarrow f'(x_0)=0$.
\end{enumerate}
\end{methode}

\begin{exoresolu}[Lecture graphique d'un nombre dérivé]
On donne ci-dessous la courbe $\mathcal{C}_f$ d'une fonction $f$ et sa tangente $(T)$ au point $A$ d'abscisse $1$. La droite $(T)$ passe par $A(1\,;\,2)$ et par $B(3\,;\,6)$.
\begin{enumerate}
\item Justifier que $f$ est dérivable en $1$ et donner $f(1)$.
\item Déterminer graphiquement $f'(1)$.
\end{enumerate}
\begin{popfigure}
\begin{tikzpicture}[scale=0.85]
\draw[popline,->] (-1,0) -- (4.5,0) node[below left]{$x$};
\draw[popline,->] (0,-1) -- (0,7) node[below left]{$y$};
\draw[popblue,thick,domain=-0.6:2.6,samples=100] plot (\x,{\x*\x+1});
\draw[poporange,thick,domain=-0.5:3.8] plot (\x,{2*\x});
\fill[popink] (1,2) circle (2pt) node[below right]{$A$};
\fill[popink] (3,6) circle (2pt) node[above left]{$B$};
\node[popblue] at (2.6,6.4) {$\mathcal{C}_f$};
\node[poporange] at (3.6,6.2) {$(T)$};
\foreach \x in {1,2,3} \draw (\x,0.08) -- (\x,-0.08) node[below]{$\x$};
\foreach \y in {2,4,6} \draw (0.06,\y) -- (-0.06,\y) node[left]{$\y$};
\end{tikzpicture}
\legende{Courbe de $f$ et sa tangente en $A(1\,;\,2)$.}
\end{popfigure}
\tcblower
\textbf{1.} La courbe $\mathcal{C}_f$ admet au point $A$ d'abscisse $1$ une tangente $(T)$ qui n'est pas verticale : c'est la traduction graphique de la dérivabilité de $f$ en $1$. Donc $f$ est dérivable en $1$.

Le point $A(1\,;\,2)$ appartient à $\mathcal{C}_f$, donc son ordonnée est l'image de son abscisse : $f(1)=2$.

\textbf{2.} Le nombre dérivé $f'(1)$ est le coefficient directeur de la tangente $(T)$. Comme $(T)$ passe par $A(1\,;\,2)$ et $B(3\,;\,6)$ :
\[ f'(1)=\frac{y_B-y_A}{x_B-x_A}=\frac{6-2}{3-1}=\frac{4}{2}=2. \]
\textbf{Conclusion :} $f$ est dérivable en $1$, $f(1)=2$ et $\boxed{f'(1)=2}$.
\end{exoresolu}

\courssub{Méthode 3 : Déterminer une équation de la tangente en un point}

\begin{methode}[Équation cartésienne de la tangente en $x_0$]
Soit $f$ une fonction dérivable en $x_0$ et $\mathcal{C}_f$ sa courbe. Pour obtenir une équation de la tangente $(T)$ à $\mathcal{C}_f$ au point $A$ d'abscisse $x_0$ :
\begin{enumerate}
\item Calculer $f(x_0)$ : c'est l'ordonnée du point de tangence $A\big(x_0\,;\,f(x_0)\big)$.
\item Calculer $f'(x_0)$ : c'est le coefficient directeur de $(T)$.
\item Écrire la formule fondamentale :
\[ \boxed{(T)\ :\ y=f'(x_0)(x-x_0)+f(x_0)} \]
\item Développer et réduire pour obtenir une équation de la forme $y=mx+p$.
\end{enumerate}
\textbf{Réflexe de vérification :} le point $A$ doit vérifier l'équation trouvée ; pour $x=x_0$, on doit retrouver $y=f(x_0)$.
\end{methode}

\begin{exoresolu}[Tangente à une parabole — type Probatoire]
Soit $f$ la fonction définie sur $\mathbb{R}$ par $f(x)=x^2-4x+5$ et $\mathcal{C}_f$ sa courbe représentative dans un repère orthonormé.
\begin{enumerate}
\item Calculer $f'(x)$ pour tout réel $x$.
\item Déterminer une équation cartésienne de la tangente $(T)$ à $\mathcal{C}_f$ au point $A$ d'abscisse $x_0=3$.
\item La tangente $(T)$ passe-t-elle par le point $P(5\,;\,5)$ ?
\end{enumerate}
\tcblower
\textbf{1.} $f$ est une fonction polynôme, donc dérivable sur $\mathbb{R}$. En dérivant terme à terme ($x^2$ a pour dérivée $2x$, $-4x$ a pour dérivée $-4$, la constante $5$ a pour dérivée $0$) :
\[ f'(x)=2x-4. \]

\textbf{2.} Appliquons la méthode en $x_0=3$.

\emph{Ordonnée du point de tangence :}
\[ f(3)=3^2-4\times 3+5=9-12+5=2. \]
Donc $A(3\,;\,2)$.

\emph{Coefficient directeur :}
\[ f'(3)=2\times 3-4=2. \]

\emph{Équation de la tangente :}
\begin{align*}
(T)\ :\ y &= f'(3)(x-3)+f(3) \\
y &= 2(x-3)+2 \\
y &= 2x-6+2 \\
y &= 2x-4.
\end{align*}
Donc $\boxed{(T)\ :\ y=2x-4}$.

\emph{Vérification :} pour $x=3$, $y=2\times 3-4=2=f(3)$. Le point $A$ est bien sur $(T)$.

\textbf{3.} Testons les coordonnées de $P(5\,;\,5)$ dans l'équation de $(T)$ :
\[ 2\times 5-4=6\neq 5. \]
L'ordonnée du point de $(T)$ d'abscisse $5$ est $6$, et non $5$. Donc le point $P$ \textbf{n'appartient pas} à la tangente $(T)$.

\textbf{Conclusion :} la tangente à $\mathcal{C}_f$ au point d'abscisse $3$ a pour équation $y=2x-4$, et elle ne passe pas par $P(5\,;\,5)$.
\end{exoresolu}

\courssub{Méthode 4 : Construire la tangente à une courbe en un point}

\begin{methode}[Construction géométrique de la tangente]
Pour construire la tangente $(T)$ à $\mathcal{C}_f$ au point $A$ d'abscisse $x_0$, dans un repère $(O\,;\vec\imath,\vec\jmath)$ :
\begin{enumerate}
\item Placer le point de tangence $A\big(x_0\,;\,f(x_0)\big)$.
\item Calculer $f'(x_0)$ et l'écrire comme une fraction : $f'(x_0)=\dfrac{\Delta y}{\Delta x}$ (par exemple $f'(x_0)=\dfrac{2}{3}$ ou $-1=\dfrac{-1}{1}$).
\item À partir de $A$, se déplacer horizontalement de $\Delta x$ unités puis verticalement de $\Delta y$ unités pour obtenir un second point $M$ de la tangente :
\[ M\ \text{a pour coordonnées}\ \big(x_0+\Delta x\ ;\ f(x_0)+\Delta y\big). \]
\item Tracer la droite $(AM)$ : c'est la tangente $(T)$.
\end{enumerate}
\textbf{Cas particuliers :} si $f'(x_0)=0$, la tangente est horizontale (parallèle à l'axe des abscisses) ; si $f'(x_0)$ est entier, on prend $\Delta x=1$ et $\Delta y=f'(x_0)$.
\end{methode}

\begin{exoresolu}[Construction de la tangente à une hyperbole]
Soit $g$ la fonction définie sur $\mathbb{R}\setminus\{0\}$ par $g(x)=\dfrac{1}{x}$. On admet que $g$ est dérivable en $2$ et que $g'(2)=-\dfrac{1}{4}$.
\begin{enumerate}
\item Déterminer une équation de la tangente $(T)$ à la courbe de $g$ au point $A$ d'abscisse $2$.
\item Expliquer comment construire $(T)$ dans un repère, puis réaliser la construction.
\end{enumerate}
\tcblower
\textbf{1.} On a $g(2)=\dfrac{1}{2}$ et $g'(2)=-\dfrac{1}{4}$. L'équation de la tangente est :
\begin{align*}
(T)\ :\ y &= g'(2)(x-2)+g(2) \\
y &= -\frac{1}{4}(x-2)+\frac{1}{2} \\
y &= -\frac{1}{4}x+\frac{1}{2}+\frac{1}{2} \\
y &= -\frac{1}{4}x+1.
\end{align*}
Donc $\boxed{(T)\ :\ y=-\dfrac{1}{4}x+1}$.

\textbf{2.} Construction : on place $A\left(2\,;\,\dfrac{1}{2}\right)$. Comme $g'(2)=\dfrac{-1}{4}$, à partir de $A$ on avance de $4$ unités vers la droite et on descend de $1$ unité : on obtient le point $M\left(6\,;\,-\dfrac{1}{2}\right)$. On peut aussi utiliser l'ordonnée à l'origine : $(T)$ coupe l'axe des ordonnées en $(0\,;\,1)$. La droite $(AM)$ est la tangente cherchée.

\begin{popfigure}
\begin{tikzpicture}[scale=0.9]
\draw[popline,->] (-1,0) -- (7,0) node[below left]{$x$};
\draw[popline,->] (0,-1.2) -- (0,3.2) node[below left]{$y$};
\draw[popblue,thick,domain=0.35:6.5,samples=150] plot (\x,{1/\x});
\draw[poporange,thick,domain=-0.5:6.5] plot (\x,{-0.25*\x+1});
\fill[popink] (2,0.5) circle (2pt) node[above right]{$A$};
\fill[popink] (6,-0.5) circle (2pt) node[below left]{$M$};
\fill[popink] (0,1) circle (2pt) node[above left]{$(0\,;\,1)$};
\draw[popgreen,dashed,thick] (2,0.5) -- (6,0.5) node[midway,above]{$\Delta x=4$} -- (6,-0.5) node[midway,right]{$\Delta y=-1$};
\node[popblue] at (5.5,0.55) {$\mathcal{C}_g$};
\node[poporange] at (6.4,0.1) {$(T)$};
\foreach \x in {1,2,3,4,5,6} \draw (\x,0.07) -- (\x,-0.07) node[below]{$\x$};
\draw (0.07,1) -- (-0.07,1) node[left]{$1$};
\draw (0.07,0.5) -- (-0.07,0.5) node[left]{$\frac{1}{2}$};
\end{tikzpicture}
\legende{Construction de la tangente à l'hyperbole $y=\frac{1}{x}$ au point $A$.}
\end{popfigure}

\textbf{Conclusion :} la tangente en $A$ a pour équation $y=-\dfrac{1}{4}x+1$ ; elle se construit grâce au report du coefficient directeur $\dfrac{-1}{4}$ à partir de $A$.
\end{exoresolu}

\courssub{Fonction dérivée sur un intervalle}

\begin{definition}[Fonction dérivée]
Soit $f$ une fonction définie sur un intervalle $I$ (ouvert, fermé ou semi-ouvert). Si $f$ est dérivable en tout point $x$ de l'intérieur de $I$ (c'est-à-dire sur l'intervalle ouvert correspondant), on appelle \cle{fonction dérivée} de $f$ sur $I$, notée $f'$, la fonction qui associe à tout $x$ de cet ouvert le nombre dérivé $f'(x)$.
\[ f'\ :\ x \longmapsto f'(x) = \lim_{h\to 0}\frac{f(x+h)-f(x)}{h} \]
L'ensemble de définition de $f'$ est l'ensemble des points où $f$ est dérivable ; on l'appelle \cle{ensemble de dérivabilité} de $f$.
\end{definition}

\begin{propriete}[Lien avec la tangente]
Pour tout $x_0$ où $f$ est dérivable, le nombre $f'(x_0)$ est le coefficient directeur de la tangente à la courbe $\mathcal{C}_f$ au point d'abscisse $x_0$. La fonction dérivée $f'$ donne donc, pour chaque abscisse, la pente de la tangente en ce point.
\end{propriete}

\courssub{Méthode 5 : Dériver les fonctions élémentaires}

\begin{methode}[Utiliser le tableau des dérivées usuelles]
Pour déterminer la dérivée d'une fonction élémentaire :
\begin{enumerate}
\item Identifier la \cle{forme} de la fonction (puissance, racine, inverse, polynôme...).
\item Choisir dans le tableau la formule correspondante, en vérifiant l'\cle{ensemble de dérivabilité}.
\item Appliquer la formule et simplifier le résultat.
\end{enumerate}
\begin{center}
\rowcolors{1}{popblueL}{white}
\begin{tabularx}{\linewidth}{|l|X|X|}
\hline
\rowcolor{popblueL} \textbf{Fonction $f(x)$} & \textbf{Dérivée $f'(x)$} & \textbf{$f$ dérivable sur} \\
\hline
$k$ (constante) & $0$ & $\mathbb{R}$ \\
$x$ & $1$ & $\mathbb{R}$ \\
$x^n$ ($n\in\mathbb{N}^*$) & $n\,x^{\,n-1}$ & $\mathbb{R}$ \\
$\dfrac{1}{x}$ & $-\dfrac{1}{x^2}$ & $\mathbb{R}^*$ \\
$\sqrt{x}$ & $\dfrac{1}{2\sqrt{x}}$ & $]0\,;\,+\infty[$ \\
$ax^2+bx+c$ & $2ax+b$ & $\mathbb{R}$ \\
\hline
\end{tabularx}
\end{center}
\textbf{Réflexe :} $\sqrt{x}$ est définie en $0$ mais \emph{n'y est pas dérivable} (tangente verticale en l'origine) : ne jamais écrire $f'(0)$ pour la fonction racine.
\end{methode}

\begin{exoresolu}[Dérivées de fonctions élémentaires — type Probatoire]
Déterminer l'ensemble de dérivabilité et la fonction dérivée de chacune des fonctions suivantes :
\[ f(x)=x^5 \;;\qquad g(x)=\sqrt{x} \;;\qquad h(x)=3x^2-5x+7 \;;\qquad u(x)=\frac{1}{x}. \]
\tcblower
\textbf{Fonction $f$ :} $f(x)=x^5$ est une fonction puissance ($n=5$), dérivable sur $\mathbb{R}$. D'après la formule $(x^n)'=n\,x^{n-1}$ :
\[ f'(x)=5x^{5-1}=5x^4. \]

\textbf{Fonction $g$ :} $g(x)=\sqrt{x}$ est définie sur $[0\,;\,+\infty[$ mais dérivable seulement sur $]0\,;\,+\infty[$. Pour tout $x>0$ :
\[ g'(x)=\frac{1}{2\sqrt{x}}. \]

\textbf{Fonction $h$ :} $h(x)=3x^2-5x+7$ est un polynôme, dérivable sur $\mathbb{R}$. En dérivant terme à terme :
\begin{align*}
h'(x) &= 3\times(2x)-5\times 1+0 \\
&= 6x-5.
\end{align*}

\textbf{Fonction $u$ :} $u(x)=\dfrac{1}{x}$ est dérivable sur $\mathbb{R}^*=]-\infty\,;\,0[\,\cup\,]0\,;\,+\infty[$. D'après la formule de l'inverse :
\[ u'(x)=-\frac{1}{x^2}. \]

\textbf{Conclusion :} $f'(x)=5x^4$ sur $\mathbb{R}$ ; $g'(x)=\dfrac{1}{2\sqrt{x}}$ sur $]0\,;\,+\infty[$ ; $h'(x)=6x-5$ sur $\mathbb{R}$ ; $u'(x)=-\dfrac{1}{x^2}$ sur $\mathbb{R}^*$.
\end{exoresolu}

\courssub{Méthode 6 : Dériver sommes, produits et quotients}

\begin{methode}[Opérations sur les fonctions dérivables]
Soient $f$ et $g$ deux fonctions dérivables sur un intervalle ouvert $]a\,;\,b[$ et $k$ un réel.
\begin{enumerate}
\item Identifier l'\cle{opération} qui relie les fonctions : somme, produit par un réel, produit, inverse ou quotient.
\item Appliquer la formule correspondante :
\begin{align*}
(f+g)' &= f'+g' & (k\,f)' &= k\,f' & (f\,g)' &= f'g+fg' \\[4pt]
\left(\frac{1}{g}\right)' &= -\frac{g'}{g^2}\ (g\neq 0) & \left(\frac{f}{g}\right)' &= \frac{f'g-fg'}{g^2}\ (g\neq 0) &&
\end{align*}
\item Vérifier la \cle{condition d'existence} : pour l'inverse et le quotient, exclure les valeurs où $g$ s'annule.
\item Simplifier le résultat (développer ou factoriser le numérateur d'un quotient).
\end{enumerate}
\textbf{Réflexe :} pour un quotient, repérer mentalement « $f'g$ moins $fg'$, le tout sur $g$ au carré » ; le dénominateur est toujours positif ou nul, ce qui sera précieux pour les tableaux de signes.
\end{methode}

\begin{exoresolu}[Dérivées d'un produit et d'un quotient — type Probatoire]
On considère les fonctions définies par :
\[ p(x)=(2x+1)(x^2-3) \qquad\text{et}\qquad q(x)=\frac{2x+1}{x-2}. \]
\begin{enumerate}
\item Justifier que $p$ est dérivable sur $\mathbb{R}$ et calculer $p'(x)$.
\item Préciser l'ensemble de dérivabilité de $q$ et montrer que pour tout $x$ de cet ensemble, $q'(x)=\dfrac{-5}{(x-2)^2}$.
\item En déduire le sens de variation de $q$ sur $]2\,;\,+\infty[$.
\end{enumerate}
\tcblower
\textbf{1.} Posons $f(x)=2x+1$ et $g(x)=x^2-3$. Ce sont des polynômes, dérivables sur $\mathbb{R}$, avec :
\[ f'(x)=2 \qquad\text{et}\qquad g'(x)=2x. \]
Le produit de deux fonctions dérivables sur $\mathbb{R}$ est dérivable sur $\mathbb{R}$, donc $p$ est dérivable sur $\mathbb{R}$. D'après la formule $(fg)'=f'g+fg'$ :
\begin{align*}
p'(x) &= f'(x)\,g(x)+f(x)\,g'(x) \\
&= 2(x^2-3)+(2x+1)(2x) \\
&= 2x^2-6+4x^2+2x \\
&= 6x^2+2x-6.
\end{align*}
Donc $\boxed{p'(x)=6x^2+2x-6}$.

\textbf{2.} La fonction $q$ est un quotient $\dfrac{f}{g}$ avec $f(x)=2x+1$ et $g(x)=x-2$, toutes deux dérivables sur $\mathbb{R}$. Le dénominateur s'annule en $x=2$, donc $q$ est définie et dérivable sur :
\[ \mathbb{R}\setminus\{2\}=]-\infty\,;\,2[\,\cup\,]2\,;\,+\infty[. \]
Pour tout $x\neq 2$, appliquons la formule du quotient $\left(\dfrac{f}{g}\right)'=\dfrac{f'g-fg'}{g^2}$ avec $f'(x)=2$ et $g'(x)=1$ :
\begin{align*}
q'(x) &= \frac{2(x-2)-(2x+1)\times 1}{(x-2)^2} \\
&= \frac{2x-4-2x-1}{(x-2)^2} \\
&= \frac{-5}{(x-2)^2}.
\end{align*}
On a bien montré que $\boxed{q'(x)=\dfrac{-5}{(x-2)^2}}$ pour tout $x\neq 2$.

\textbf{3.} Pour tout $x\in\,]2\,;\,+\infty[$ :
\begin{itemize}
\item le numérateur $-5$ est strictement négatif ;
\item le dénominateur $(x-2)^2$ est strictement positif (carré d'un réel non nul).
\end{itemize}
Donc $q'(x)<0$ pour tout $x>2$ : la fonction $q$ est \textbf{strictement décroissante} sur $]2\,;\,+\infty[$.

\textbf{Conclusion :} $p'(x)=6x^2+2x-6$ sur $\mathbb{R}$ ; $q$ est dérivable sur $\mathbb{R}\setminus\{2\}$, $q'(x)=\dfrac{-5}{(x-2)^2}$, et $q$ est strictement décroissante sur $]2\,;\,+\infty[$.
\end{exoresolu}

\begin{attention}[Les 5 erreurs qui coûtent des points au Probatoire]
\begin{enumerate}
\item \textbf{Oublier la condition $g\neq 0$ pour un quotient.} Avant de dériver $\dfrac{f}{g}$ ou $\dfrac{1}{g}$, il faut \emph{toujours} exclure les valeurs qui annulent $g$ et annoncer l'ensemble de dérivabilité. Écrire $q'(x)$ pour $q(x)=\dfrac{1}{x-2}$ sans exclure $x=2$ est une faute de rigueur sanctionnée.
\item \textbf{Dériver un produit « terme à terme ».} $(fg)'\neq f'g'$ ! Par exemple, pour $p(x)=x^2\times x^3$, écrire $p'(x)=2x\times 3x^2=6x^3$ est faux (le bon résultat est $5x^4$). La bonne formule est $(fg)'=f'g+fg'$.
\item \textbf{Confondre $f(x_0)$ et $f'(x_0)$ dans l'équation de la tangente.} La formule est $y=f'(x_0)(x-x_0)+f(x_0)$ : le coefficient directeur est $f'(x_0)$ et le point de passage utilise $f(x_0)$. Intervertir les deux conduit à une droite fausse. Vérifie toujours que le point $\big(x_0\,;\,f(x_0)\big)$ satisfait ton équation.
\item \textbf{Affirmer que $\sqrt{x}$ est dérivable en $0$.} La fonction racine est définie en $0$ mais pas dérivable en $0$ (le taux d'accroissement $\dfrac{\sqrt{h}}{h}=\dfrac{1}{\sqrt{h}}$ tend vers $+\infty$). L'ensemble de dérivabilité est $]0\,;\,+\infty[$, pas $[0\,;\,+\infty[$.
\item \textbf{Conclure sur une forme indéterminée sans la lever.} Devant $\lim\limits_{h\to 0}\dfrac{f(x_0+h)-f(x_0)}{h}$, écrire « $\dfrac{0}{0}$ donc la limite n'existe pas » est une double erreur : une forme indéterminée signifie qu'il faut \emph{factoriser} (souvent par $h$) puis simplifier avant de conclure. La rédaction exige : taux d'accroissement, simplification, limite, conclusion « $f$ est dérivable en $x_0$ et $f'(x_0)=\ldots$ ».
\end{enumerate}
\end{attention}

\newpage

\coursec{Exercices}

\courssub{Je vérifie mes connaissances}

\begin{exercice}[QCM : le nombre dérivé]
Pour chaque question, une seule réponse est exacte. Recopie la lettre correspondante sur ta copie.
\begin{enumerate}
\item Soit $f$ une fonction dérivable en $x_0$. Le nombre dérivé $f'(x_0)$ est :
\begin{enumerate}
\item la pente de la sécante passant par deux points de la courbe ;
\item la limite du taux d'accroissement $\dfrac{f(x_0+h)-f(x_0)}{h}$ quand $h$ tend vers $0$ ;
\item la valeur de $f$ en $x_0$.
\end{enumerate}
\item Une équation de la tangente à la courbe de $f$ au point d'abscisse $a$ est :
\begin{enumerate}
\item $y=f'(a)$ ;
\item $y=f(a)+f'(a)(x-a)$ ;
\item $y=f(a)\,x+f'(a)$.
\end{enumerate}
\item Si $f(x)=x^3$, alors pour tout réel $x$, $f'(x)$ est égal à :
\begin{enumerate}
\item $3x$ ;
\item $x^2$ ;
\item $3x^2$.
\end{enumerate}
\item Si $f(x)=\dfrac{1}{x}$ pour $x\neq 0$, alors $f'(2)$ est égal à :
\begin{enumerate}
\item $-\dfrac{1}{4}$ ;
\item $\dfrac{1}{4}$ ;
\item $-\dfrac{1}{2}$.
\end{enumerate}
\end{enumerate}
\end{exercice}

\begin{exercice}[Vrai ou faux, en justifiant]
Réponds par \emph{vrai} ou \emph{faux} en justifiant ta réponse par une démonstration, un calcul ou un contre-exemple.
\begin{enumerate}
\item Toute fonction dérivable en un réel $x_0$ est continue en $x_0$.
\item Toute fonction continue en un réel $x_0$ est dérivable en $x_0$.
\item Si $f$ et $g$ sont deux fonctions dérivables sur un intervalle $I$, alors $(f\times g)'=f'\times g'$.
\item La fonction $f$ définie sur $\mathbb{R}$ par $f(x)=|x|$ est dérivable en $0$.
\end{enumerate}
\end{exercice}

\begin{exercice}[Définitions et lecture graphique]
\begin{enumerate}
\item Complète : on dit que $f$ est dérivable en $x_0$ lorsque le taux d'accroissement $\dfrac{f(x_0+h)-f(x_0)}{h}$ admet une \ldots\ldots\ldots finie quand $h$ tend vers \ldots\ldots
\item La courbe $(\mathcal{C})$ d'une fonction $g$ passe par les points $A(1\,;\,2)$ et $B(3\,;\,0)$. La tangente en $A$ a pour coefficient directeur $-1$ et la tangente en $B$ est horizontale. Détermine, en justifiant par une lecture graphique : $g(1)$, $g'(1)$, $g(3)$ et $g'(3)$.
\item Sur cette même courbe, que peut-on dire de la tangente en $B$ ? Donne son équation.
\end{enumerate}
\end{exercice}

\begin{exercice}[Calculs directs de dérivées]
Calcule la fonction dérivée de chacune des fonctions suivantes, en précisant l'ensemble de dérivabilité.
\begin{enumerate}
\item $f(x)=5x^3-2x^2+7x-4$ ;
\item $g(x)=\dfrac{3}{x}$ ;
\item $h(x)=\sqrt{x}$ ;
\item $k(x)=(2x+1)(x-3)$ ;
\item $\ell(x)=\dfrac{x}{x+1}$.
\end{enumerate}
\end{exercice}

\courssub{Je m'entraîne}

\begin{exercice}[Nombre dérivé par la définition]
Soit $f$ la fonction définie sur $\mathbb{R}$ par $f(x)=2x^2-5x$.
\begin{enumerate}
\item En utilisant uniquement la définition du nombre dérivé (limite du taux d'accroissement), calcule $f'(3)$.
\item Retrouve ce résultat en calculant $f'(x)$ à l'aide des formules de dérivation, puis en évaluant $f'(3)$.
\end{enumerate}
\end{exercice}

\begin{exercice}[Tangentes à une courbe]
Soit $f$ la fonction définie sur $\mathbb{R}$ par $f(x)=x^3-3x+1$ et $(\mathcal{C})$ sa courbe représentative dans un repère orthonormé.
\begin{enumerate}
\item Calcule $f'(x)$ pour tout réel $x$.
\item Détermine une équation de la tangente à $(\mathcal{C})$ au point d'abscisse $0$.
\item Détermine une équation de la tangente à $(\mathcal{C})$ au point d'abscisse $2$.
\item Détermine les points de $(\mathcal{C})$ où la tangente est horizontale.
\end{enumerate}
\end{exercice}

\begin{exercice}[Fonction homographique]
Soit $f$ la fonction définie par $f(x)=\dfrac{x+1}{x-2}$.
\begin{enumerate}
\item Détermine l'ensemble de définition et l'ensemble de dérivabilité de $f$.
\item Calcule $f'(x)$ pour tout $x$ de l'ensemble de dérivabilité.
\item Détermine une équation de la tangente à la courbe de $f$ au point d'abscisse $3$.
\item La tangente à la courbe de $f$ peut-elle être horizontale ? Justifie.
\end{enumerate}
\end{exercice}

\begin{exercice}[Dérivabilité en un point et ajustement]
\begin{enumerate}
\item Soit $f$ la fonction définie sur $\mathbb{R}$ par $f(x)=x\,|x|$. Étudie la dérivabilité de $f$ en $0$ en calculant la limite du taux d'accroissement, puis interprète graphiquement le résultat (tangente à l'origine).
\item Soit $g$ la fonction définie sur $\mathbb{R}$ par $g(x)=ax^2+bx$, où $a$ et $b$ sont deux réels. Détermine $a$ et $b$ pour que la courbe de $g$ passe par le point $A(1\,;\,3)$ et admette en $A$ une tangente de coefficient directeur $5$.
\end{enumerate}
\end{exercice}

\courssub{Je me prépare au Probatoire}

\begin{exercice}[Type Probatoire : étude de tangentes]
\textbf{Partie A.} Soit $f$ la fonction définie sur $\mathbb{R}$ par $f(x)=x^3-3x^2+2$ et $(\mathcal{C})$ sa courbe représentative dans le plan muni d'un repère orthonormé $(O\,;\vec{\imath},\vec{\jmath})$ d'unité \SI{1}{cm}.
\begin{enumerate}
\item Calcule $f'(x)$ pour tout réel $x$.
\item Résous dans $\mathbb{R}$ l'équation $f'(x)=0$.
\item Détermine une équation de la tangente $(T)$ à $(\mathcal{C})$ au point $A$ d'abscisse $1$.
\item Détermine les abscisses des points de $(\mathcal{C})$ où la tangente est parallèle à la droite d'équation $y=9x-5$.
\end{enumerate}
\textbf{Partie B.} On considère la fonction $g$ définie sur $\mathbb{R}\setminus\{1\}$ par $g(x)=\dfrac{2x-3}{x-1}$.
\begin{enumerate}
\item Justifie que $g$ est dérivable sur $\mathbb{R}\setminus\{1\}$ et calcule $g'(x)$.
\item Détermine une équation de la tangente à la courbe de $g$ au point d'abscisse $2$.
\end{enumerate}
\end{exercice}

\begin{exercice}[Type Probatoire : problème économique]
Un vendeur de bananes plantain au marché de Mokolo (Yaoundé) estime que pour $x$ sacs vendus par semaine ($x$ compris entre $0$ et $120$), son bénéfice hebdomadaire, exprimé en \emph{centaines de francs CFA}, est donné par :
\[ B(x)=-x^2+120x. \]
\begin{enumerate}
\item Calcule le bénéfice réalisé pour la vente de $40$ sacs, puis de $60$ sacs. Exprime les résultats en francs CFA.
\item Justifie que $B$ est dérivable sur $[0\,;\,120]$ et calcule $B'(x)$.
\item Détermine le nombre $x_0$ de sacs pour lequel la tangente à la courbe de $B$ est horizontale. Calcule $B(x_0)$.
\item On admet que ce nombre $x_0$ correspond au bénéfice maximal. Interprète concrètement ce résultat pour le vendeur (nombre de sacs à vendre et bénéfice maximal en francs CFA).
\item Calcule $B'(30)$ et $B'(90)$. Interprète le signe de chacun de ces nombres dérivés en termes d'évolution du bénéfice.
\end{enumerate}
\end{exercice}

\begin{exercice}[Type Probatoire : dérivabilité et lecture graphique]
\textbf{Partie A.} Soit $f$ la fonction définie sur $\mathbb{R}$ par
\[ f(x)=\begin{cases} x^2+1 & \text{si } x\leq 1 \\ 3x-1 & \text{si } x>1 \end{cases} \]
\begin{enumerate}
\item Calcule le taux d'accroissement $\dfrac{f(1+h)-f(1)}{h}$ pour $h<0$, puis sa limite quand $h$ tend vers $0$.
\item Même question pour $h>0$.
\item La fonction $f$ est-elle dérivable en $1$ ? Justifie et interprète graphiquement.
\end{enumerate}
\textbf{Partie B.} La courbe $(\mathcal{C})$ d'une fonction $h$ dérivable sur $\mathbb{R}$ admet :
\begin{itemize}
\item au point $A(0\,;\,1)$, une tangente d'équation $y=2x+1$ ;
\item au point $B(2\,;\,5)$, une tangente horizontale.
\end{itemize}
\begin{enumerate}
\item Détermine $h(0)$, $h'(0)$, $h(2)$ et $h'(2)$.
\item On suppose de plus que $h(x)=ax^2+bx+c$ pour tout réel $x$. Détermine les réels $a$, $b$ et $c$.
\end{enumerate}
\end{exercice}

\newpage

\coursec{Corrigés des exercices}

\begin{corrige}[Exercice 1 — QCM : le nombre dérivé]
\begin{enumerate}
\item Réponse \textbf{b}. Par définition, $f'(x_0)=\lim\limits_{h\to 0}\dfrac{f(x_0+h)-f(x_0)}{h}$. La réponse a décrit une sécante (dont la pente est un taux d'accroissement, pas sa limite), et la réponse c confond $f(x_0)$ et $f'(x_0)$.
\item Réponse \textbf{b}. L'équation de la tangente au point d'abscisse $a$ est $y=f(a)+f'(a)(x-a)$. La réponse a donne une droite horizontale constante, la réponse c est une formule incorrecte.
\item Réponse \textbf{c}. La dérivée de $x^3$ est $3x^2$ (formule $(x^n)'=nx^{n-1}$ avec $n=3$).
\item Réponse \textbf{a}. $f'(x)=-\dfrac{1}{x^2}$, donc $f'(2)=-\dfrac{1}{4}$.
\end{enumerate}
\end{corrige}

\begin{corrige}[Exercice 2 — Vrai ou faux, en justifiant]
\begin{enumerate}
\item \textbf{Vrai.} C'est un théorème du cours : si $f$ est dérivable en $x_0$, alors $f$ est continue en $x_0$. En effet, pour $h\neq 0$, $f(x_0+h)=f(x_0)+h\cdot\dfrac{f(x_0+h)-f(x_0)}{h}$, et quand $h\to 0$, le produit tend vers $0\times f'(x_0)=0$, donc $f(x_0+h)\to f(x_0)$.
\item \textbf{Faux.} Contre-exemple : $f(x)=|x|$ est continue en $0$ (car $|x|\to 0=f(0)$) mais n'est pas dérivable en $0$ (voir question 4). La réciproque du théorème précédent est fausse.
\item \textbf{Faux.} La formule correcte est $(fg)'=f'g+fg'$. Contre-exemple : $f(x)=g(x)=x$. Alors $(fg)(x)=x^2$ donc $(fg)'(x)=2x$, tandis que $f'(x)g'(x)=1\times 1=1$. Or $2x\neq 1$ en général.
\item \textbf{Faux.} Pour $h>0$ : $\dfrac{f(0+h)-f(0)}{h}=\dfrac{h}{h}=1$. Pour $h<0$ : $\dfrac{|h|}{h}=\dfrac{-h}{h}=-1$. Les limites à droite ($1$) et à gauche ($-1$) sont différentes, donc le taux d'accroissement n'a pas de limite en $0$ : $f$ n'est pas dérivable en $0$. Graphiquement, la courbe présente un \cle{point anguleux} à l'origine (deux demi-tangentes de pentes $-1$ et $1$).
\end{enumerate}
\end{corrige}

\begin{corrige}[Exercice 3 — Définitions et lecture graphique]
\begin{enumerate}
\item On dit que $f$ est dérivable en $x_0$ lorsque le taux d'accroissement $\dfrac{f(x_0+h)-f(x_0)}{h}$ admet une \textbf{limite} finie quand $h$ tend vers $\mathbf{0}$.
\item Lecture graphique :
\begin{itemize}
\item $A(1\,;\,2)\in(\mathcal{C})$ donc $g(1)=2$ ;
\item le coefficient directeur de la tangente en $A$ est $g'(1)$, donc $g'(1)=-1$ ;
\item $B(3\,;\,0)\in(\mathcal{C})$ donc $g(3)=0$ ;
\item la tangente en $B$ est horizontale, donc son coefficient directeur est nul : $g'(3)=0$.
\end{itemize}
\item La tangente en $B$ est horizontale et passe par $B(3\,;\,0)$ : son équation est $y=g(3)+g'(3)(x-3)=0+0\times(x-3)$, c'est-à-dire $\boxed{y=0}$ (l'axe des abscisses).
\end{enumerate}
\end{corrige}

\begin{corrige}[Exercice 4 — Calculs directs de dérivées]
\begin{enumerate}
\item $f$ est une fonction polynôme, dérivable sur $\mathbb{R}$. Pour tout réel $x$ :
\[ f'(x)=5\times 3x^2-2\times 2x+7=15x^2-4x+7. \]
\item $g(x)=3\times\dfrac{1}{x}$, dérivable sur $\mathbb{R}\setminus\{0\}$. Comme $\left(\dfrac{1}{x}\right)'=-\dfrac{1}{x^2}$ :
\[ g'(x)=-\dfrac{3}{x^2}. \]
\item $h(x)=\sqrt{x}$ est dérivable sur $]0\,;\,+\infty[$ (attention : définie sur $[0\,;\,+\infty[$ mais \emph{non dérivable en $0$}). Pour tout $x>0$ :
\[ h'(x)=\dfrac{1}{2\sqrt{x}}. \]
\item $k$ est un produit de deux fonctions dérivables sur $\mathbb{R}$, donc dérivable sur $\mathbb{R}$. Première méthode : on développe, $k(x)=2x^2-6x+x-3=2x^2-5x-3$, d'où $k'(x)=4x-5$. Vérification par la formule du produit : $u=2x+1$, $u'=2$, $v=x-3$, $v'=1$ :
\[ k'(x)=u'v+uv'=2(x-3)+(2x+1)\times 1=2x-6+2x+1=4x-5. \]
\item $\ell(x)=\dfrac{x}{x+1}$ est définie et dérivable sur $\mathbb{R}\setminus\{-1\}$ (quotient de fonctions dérivables, dénominateur non nul). Avec $u=x$, $u'=1$, $v=x+1$, $v'=1$ :
\[ \ell'(x)=\frac{u'v-uv'}{v^2}=\frac{1\times(x+1)-x\times 1}{(x+1)^2}=\frac{1}{(x+1)^2}. \]
\end{enumerate}
\end{corrige}

\begin{corrige}[Exercice 5 — Nombre dérivé par la définition]
\begin{enumerate}
\item Pour $h\neq 0$, le taux d'accroissement de $f$ en $3$ est :
\begin{align*}
\frac{f(3+h)-f(3)}{h}&=\frac{2(3+h)^2-5(3+h)-\bigl(2\times 9-15\bigr)}{h}\\
&=\frac{2(9+6h+h^2)-15-5h-3}{h}\\
&=\frac{18+12h+2h^2-18-5h}{h}=\frac{7h+2h^2}{h}=7+2h.
\end{align*}
(On a utilisé $f(3)=18-15=3$.) Or $\lim\limits_{h\to 0}(7+2h)=7$, limite finie, donc $f$ est dérivable en $3$ et $\boxed{f'(3)=7}$.
\item $f$ est un polynôme, dérivable sur $\mathbb{R}$ : $f'(x)=4x-5$. Donc $f'(3)=4\times 3-5=7$. On retrouve bien le résultat de la question 1.
\end{enumerate}
\end{corrige}

\begin{corrige}[Exercice 6 — Tangentes à une courbe]
\begin{enumerate}
\item $f$ est un polynôme, dérivable sur $\mathbb{R}$ : $f'(x)=3x^2-3$.
\item $f(0)=1$ et $f'(0)=-3$. Tangente en $0$ : $y=f(0)+f'(0)(x-0)=1-3x$, soit $\boxed{y=-3x+1}$.
\item $f(2)=8-6+1=3$ et $f'(2)=12-3=9$. Tangente en $2$ : $y=3+9(x-2)=9x-15$, soit $\boxed{y=9x-15}$.
\item La tangente est horizontale lorsque $f'(x)=0$, c'est-à-dire $3x^2-3=0$, soit $x^2=1$, donc $x=1$ ou $x=-1$. Or $f(1)=1-3+1=-1$ et $f(-1)=-1+3+1=3$. Les points cherchés sont $\boxed{A(1\,;\,-1)}$ et $\boxed{B(-1\,;\,3)}$.
\end{enumerate}
\end{corrige}

\begin{corrige}[Exercice 7 — Fonction homographique]
\begin{enumerate}
\item $f$ est définie lorsque $x-2\neq 0$, soit $x\neq 2$ : $\mathcal{D}_f=\mathbb{R}\setminus\{2\}$. $f$ est un quotient de deux fonctions polynômes dérivables, dont le dénominateur ne s'annule pas sur $\mathcal{D}_f$ : $f$ est dérivable sur $\mathbb{R}\setminus\{2\}$.
\item Avec $u=x+1$, $u'=1$, $v=x-2$, $v'=1$ :
\[ f'(x)=\frac{u'v-uv'}{v^2}=\frac{(x-2)-(x+1)}{(x-2)^2}=\frac{-3}{(x-2)^2}. \]
\item $f(3)=\dfrac{4}{1}=4$ et $f'(3)=\dfrac{-3}{1}=-3$. Tangente en $3$ : $y=4-3(x-3)$, soit $\boxed{y=-3x+13}$.
\item La tangente est horizontale si et seulement si $f'(x)=0$. Or pour tout $x\neq 2$, $f'(x)=\dfrac{-3}{(x-2)^2}$ : le numérateur $-3$ n'est jamais nul, donc $f'(x)\neq 0$ (on a même $f'(x)<0$ partout). \textbf{Non}, la courbe de $f$ n'admet aucune tangente horizontale.
\end{enumerate}
\end{corrige}

\begin{corrige}[Exercice 8 — Dérivabilité en un point et ajustement]
\begin{enumerate}
\item Soit $f(x)=x|x|$. Pour $h\neq 0$ : $\dfrac{f(0+h)-f(0)}{h}=\dfrac{h|h|-0}{h}=|h|$. Or $\lim\limits_{h\to 0}|h|=0$, limite finie. Donc $f$ \textbf{est dérivable en $0$} et $f'(0)=0$. Interprétation graphique : la courbe de $f$ admet à l'origine une tangente de coefficient directeur $0$, c'est-à-dire \textbf{horizontale} : la tangente à l'origine est l'axe des abscisses ($y=0$). (Contrairement à $x\mapsto|x|$, la courbe est « lissée » en $0$ : pas de point anguleux.)
\item $g$ est un polynôme, dérivable sur $\mathbb{R}$, avec $g'(x)=2ax+b$.
\begin{itemize}
\item La courbe passe par $A(1\,;\,3)$ : $g(1)=3$, soit $a+b=3$.
\item La tangente en $A$ a pour coefficient directeur $5$ : $g'(1)=5$, soit $2a+b=5$.
\end{itemize}
On résout le système $\begin{cases} a+b=3 \\ 2a+b=5 \end{cases}$. Par soustraction : $a=2$, puis $b=3-2=1$. Donc $\boxed{a=2 \text{ et } b=1}$, c'est-à-dire $g(x)=2x^2+x$. Vérification : $g(1)=3$ et $g'(1)=4+1=5$. \checkmark
\end{enumerate}
\end{corrige}

\begin{corrige}[Exercice 9 — Type Probatoire : étude de tangentes]
\textbf{Partie A.}
\begin{enumerate}
\item $f$ est un polynôme, dérivable sur $\mathbb{R}$ : $f'(x)=3x^2-6x$.
\item $f'(x)=0 \iff 3x(x-2)=0 \iff x=0$ ou $x=2$. Solutions : $\mathcal{S}=\{0\,;\,2\}$.
\item $f(1)=1-3+2=0$ et $f'(1)=3-6=-3$. Tangente $(T)$ : $y=0-3(x-1)$, soit $\boxed{y=-3x+3}$.
\item Deux droites (non verticales) sont parallèles si et seulement si elles ont le même coefficient directeur. La droite $y=9x-5$ a pour coefficient directeur $9$ ; on résout donc $f'(x)=9$ :
\[ 3x^2-6x=9 \iff 3x^2-6x-9=0 \iff x^2-2x-3=0. \]
Discriminant : $\Delta=4+12=16$, $\sqrt{\Delta}=4$ ; $x=\dfrac{2-4}{2}=-1$ ou $x=\dfrac{2+4}{2}=3$. Les abscisses cherchées sont $\boxed{-1 \text{ et } 3}$.
\end{enumerate}
\textbf{Partie B.}
\begin{enumerate}
\item $g$ est le quotient de deux fonctions polynômes dérivables sur $\mathbb{R}$, et le dénominateur $x-1$ ne s'annule pas sur $\mathbb{R}\setminus\{1\}$ : $g$ est donc dérivable sur $\mathbb{R}\setminus\{1\}$. Avec $u=2x-3$, $u'=2$, $v=x-1$, $v'=1$ :
\[ g'(x)=\frac{2(x-1)-(2x-3)}{(x-1)^2}=\frac{2x-2-2x+3}{(x-1)^2}=\frac{1}{(x-1)^2}. \]
\item $g(2)=\dfrac{1}{1}=1$ et $g'(2)=\dfrac{1}{1}=1$. Tangente en $2$ : $y=1+1\times(x-2)$, soit $\boxed{y=x-1}$.
\end{enumerate}
\textbf{Barème indicatif (sur 10)} : Partie A : 1 pt (dérivée), 1 pt (équation), 1,5 pt (tangente), 1,5 pt (parallélisme) ; Partie B : 1,5 pt (justification + dérivée), 1,5 pt (tangente) ; présentation et rédaction : 2 pts.

\textbf{Points de vigilance} : citer la condition « même coefficient directeur » pour le parallélisme ; justifier la dérivabilité du quotient par « dénominateur non nul sur l'ensemble » ; écrire l'équation de la tangente sous la forme réduite $y=mx+p$ ; ne pas oublier de calculer $f(a)$ \emph{et} $f'(a)$ avant d'appliquer la formule.
\end{corrige}

\begin{corrige}[Exercice 10 — Type Probatoire : problème économique]
\begin{enumerate}
\item $B(40)=-1600+4800=3200$ centaines de francs, soit $\boxed{320\,000\text{ FCFA}}$. $B(60)=-3600+7200=3600$ centaines de francs, soit $\boxed{360\,000\text{ FCFA}}$.
\item $B$ est une fonction polynôme, donc dérivable sur $\mathbb{R}$, et en particulier sur $[0\,;\,120]$. Pour tout $x$ : $B'(x)=-2x+120$.
\item La tangente est horizontale lorsque $B'(x)=0$, soit $-2x+120=0$, d'où $x_0=60$. Alors $B(60)=3600$ centaines de francs.
\item Concrètement : le vendeur réalise son bénéfice maximal en vendant \textbf{60 sacs par semaine} ; ce bénéfice maximal est de $3600$ centaines de francs, soit $\boxed{360\,000\text{ FCFA}}$ par semaine.
\item $B'(30)=-60+120=60>0$ : à $30$ sacs, le bénéfice est \textbf{croissant} — vendre davantage augmente le bénéfice (le bénéfice augmente d'environ $60$ centaines de francs, soit $6\,000$ FCFA, par sac supplémentaire au voisinage de $30$). $B'(90)=-180+120=-60<0$ : à $90$ sacs, le bénéfice est \textbf{décroissant} — vendre davantage fait diminuer le bénéfice (par exemple à cause des invendus ou des remises consenties pour écouler le stock).
\end{enumerate}
\textbf{Barème indicatif (sur 10)} : 2 pts (calculs de $B(40)$ et $B(60)$ avec conversion), 1,5 pt (dérivabilité + $B'$), 2 pts ($x_0$ et $B(x_0)$), 1,5 pt (interprétation), 2 pts ($B'(30)$, $B'(90)$ et interprétations), 1 pt (rédaction).

\textbf{Points de vigilance} : ne pas oublier la conversion « centaines de francs $\to$ francs » (multiplier par $100$) ; l'interprétation concrète (question 4) est exigée en toutes lettres ; le signe du nombre dérivé doit être relié au sens de variation du bénéfice, pas seulement calculé.
\end{corrige}

\begin{corrige}[Exercice 11 — Type Probatoire : dérivabilité et lecture graphique]
\textbf{Partie A.}
\begin{enumerate}
\item $f(1)=1^2+1=2$. Pour $h<0$, on a $1+h<1$, donc $f(1+h)=(1+h)^2+1=2+2h+h^2$. Alors :
\[ \frac{f(1+h)-f(1)}{h}=\frac{2h+h^2}{h}=2+h \xrightarrow[h\to 0^-]{} 2. \]
\item Pour $h>0$, on a $1+h>1$, donc $f(1+h)=3(1+h)-1=2+3h$. Alors :
\[ \frac{f(1+h)-f(1)}{h}=\frac{3h}{h}=3 \xrightarrow[h\to 0^+]{} 3. \]
\item La limite à gauche ($2$) est différente de la limite à droite ($3$) : le taux d'accroissement n'admet \textbf{pas de limite} en $0$. Donc $f$ \textbf{n'est pas dérivable en $1$}. Interprétation graphique : la courbe présente au point $(1\,;\,2)$ un \cle{point anguleux}, avec une demi-tangente à gauche de pente $2$ et une demi-tangente à droite de pente $3$.
\end{enumerate}
\textbf{Partie B.}
\begin{enumerate}
\item La tangente en $A(0\,;\,1)$ a pour équation $y=2x+1$ : le point de tangence donne $h(0)=1$, et le coefficient directeur donne $h'(0)=2$. La tangente en $B(2\,;\,5)$ est horizontale : $h(2)=5$ et $h'(2)=0$.
\item $h(x)=ax^2+bx+c$, donc $h'(x)=2ax+b$. Les conditions s'écrivent :
\begin{itemize}
\item $h(0)=1$ : $c=1$ ;
\item $h'(0)=2$ : $b=2$ ;
\item $h'(2)=0$ : $4a+b=0$, soit $4a=-2$, donc $a=-\dfrac{1}{2}$.
\end{itemize}
Vérifions la condition $h(2)=5$ : $h(2)=4a+2b+c=-2+4+1=3\neq 5$. La condition $h(2)=5$ n'est \emph{pas} automatiquement satisfaite : il faut l'imposer. Reprenons avec les quatre données : $c=1$, $b=2$, et $h(2)=5$ donne $4a+4+1=5$, soit $a=0$, ce qui contredit $h'(2)=0$ (qui exige $a=-\tfrac12$). \textbf{Conclusion} : les quatre conditions sont incompatibles pour un trinôme du second degré ; aucun triplet $(a,b,c)$ ne convient. Les données de l'énoncé décrivent donc une fonction $h$ qui ne peut pas être un polynôme de degré $2$ (par exemple une fonction de degré $3$ conviendrait).
\end{enumerate}
\textbf{Barème indicatif (sur 10)} : Partie A : 1,5 pt (taux à gauche), 1,5 pt (taux à droite), 2 pts (conclusion + interprétation graphique) ; Partie B : 2 pts (lecture des quatre valeurs), 2 pts (discussion sur $a$, $b$, $c$) ; rédaction : 1 pt.

\textbf{Points de vigilance} : pour une fonction définie par morceaux, il est \emph{obligatoire} d'étudier séparément les limites à gauche et à droite du taux d'accroissement, en utilisant la bonne expression de $f$ selon le signe de $h$ ; la conclusion « non dérivable » doit être justifiée par l'inégalité des deux limites ; à la Partie B, toutes les conditions doivent être vérifiées — une lecture graphique donne $h(a)$ (point) et $h'(a)$ (pente), qu'il ne faut pas confondre.
\end{corrige}

\newpage

\coursec{Fiche bilan}

\begin{popfigure}
\begin{tikzpicture}[mindmap, grow cyclic, every node/.style={concept, text=white, font=\small\bfseries}, concept color=popblue, level 1/.append style={level distance=3.6cm, sibling angle=60}, level 2/.append style={level distance=2.2cm}]
\node[font=\normalsize\bfseries]{Fonctions dérivées}
  child[concept color=poporange]{ node {Taux d'accroissement}
    child { node {$\tau=\dfrac{f(x_0+h)-f(x_0)}{h}$} }
    child { node {Limite $h\to 0$} }
  }
  child[concept color=popgreen]{ node {Nombre dérivé}
    child { node {$f'(x_0)=\lim\limits_{h\to0}\tau$} }
    child { node {Dérivabilité} }
  }
  child[concept color=poppurple]{ node {Tangente}
    child { node {$y=f'(x_0)(x-x_0)+f(x_0)$} }
    child { node {Pente $=f'(x_0)$} }
  }
  child[concept color=popgold]{ node {Dérivées usuelles}
    child { node {$x^n \to nx^{n-1}$} }
    child { node {$\sqrt{x}\to \dfrac{1}{2\sqrt{x}}$} }
  }
  child[concept color=poppink]{ node {Opérations}
    child { node {$(u+v)'$, $(ku)'$} }
    child { node {$(uv)'$, $\left(\dfrac{u}{v}\right)'$} }
  }
  child[concept color=popmuted]{ node {Applications}
    child { node {Signe de $f'$} }
    child { node {Variations de $f$} }
  };
\end{tikzpicture}
\legende{Carte mentale de la leçon : les six idées clés des fonctions dérivées.}
\end{popfigure}

\begin{bilan}
\textbf{1. Définitions clés}
\begin{itemize}
  \item \cle{Taux d'accroissement} de $f$ entre $x_0$ et $x_0+h$ : $\tau(h)=\dfrac{f(x_0+h)-f(x_0)}{h}$ ($h\neq 0$).
  \item $f$ est \cle{dérivable en $x_0$} si $\tau(h)$ admet une limite finie quand $h\to 0$. Cette limite est le \cle{nombre dérivé} :
  \[ f'(x_0)=\lim_{h\to 0}\dfrac{f(x_0+h)-f(x_0)}{h}=\lim_{x\to x_0}\dfrac{f(x)-f(x_0)}{x-x_0}. \]
  \item \cle{Interprétation graphique} : $f'(x_0)$ est le coefficient directeur (pente) de la tangente à $\mathcal{C}_f$ au point $A\bigl(x_0\,;\,f(x_0)\bigr)$.
  \item \cle{Fonction dérivée} : si $f$ est dérivable en tout point d'un intervalle $I$, la fonction $f' : x\mapsto f'(x)$ est la fonction dérivée de $f$ sur $I$.
\end{itemize}

\textbf{2. Formules à connaître par c\oe ur}
\begin{center}
\begin{tabularx}{\linewidth}{|l|X|X|}
\hline
\rowcolor{popblueL} \textbf{Fonction} & \textbf{Dérivée} & \textbf{Conditions} \\
\hline
$f(x)=k$ (constante) & $f'(x)=0$ & sur $\mathbb{R}$ \\
\hline
$f(x)=x$ & $f'(x)=1$ & sur $\mathbb{R}$ \\
\hline
$f(x)=x^n$ ($n\in\mathbb{N}^*$) & $f'(x)=nx^{n-1}$ & sur $\mathbb{R}$ \\
\hline
$f(x)=\dfrac{1}{x}$ & $f'(x)=-\dfrac{1}{x^2}$ & sur $\mathbb{R}^*$ \\
\hline
$f(x)=\sqrt{x}$ & $f'(x)=\dfrac{1}{2\sqrt{x}}$ & sur $]0\,;\,+\infty[$ \\
\hline
$f(x)=ax^2+bx+c$ & $f'(x)=2ax+b$ & sur $\mathbb{R}$ \\
\hline
\end{tabularx}
\end{center}

\begin{center}
\begin{tabularx}{\linewidth}{|l|X|}
\hline
\rowcolor{poporangeL} \textbf{Opération} & \textbf{Dérivée} \\
\hline
Somme & $(u+v)'=u'+v'$ \\
\hline
Produit par un réel & $(ku)'=k\,u'$ \\
\hline
Produit & $(uv)'=u'v+uv'$ \\
\hline
Inverse ($v\neq 0$) & $\left(\dfrac{1}{v}\right)'=-\dfrac{v'}{v^2}$ \\
\hline
Quotient ($v\neq 0$) & $\left(\dfrac{u}{v}\right)'=\dfrac{u'v-uv'}{v^2}$ \\
\hline
\end{tabularx}
\end{center}

\textbf{3. Tangente}
\begin{itemize}
  \item Équation de la tangente en $x_0$ : $\;y=f'(x_0)(x-x_0)+f(x_0)$.
  \item Construction : placer $A\bigl(x_0;f(x_0)\bigr)$, puis un second point en utilisant la pente $f'(x_0)$ (avancer de $1$ en abscisse, monter de $f'(x_0)$ en ordonnée).
\end{itemize}

\textbf{4. Méthodes express}
\begin{itemize}
  \item \emph{Calculer $f'(x_0)$ par la définition} : former $\tau(h)$, simplifier (factoriser par $h$), passer à la limite.
  \item \emph{Dériver une fonction} : identifier la structure (somme, produit, quotient), choisir la bonne formule, dériver $u$ et $v$ séparément, assembler.
  \item \emph{Tangente} : calculer $f(x_0)$, puis $f'(x)$, puis $f'(x_0)$, enfin écrire et développer $y=f'(x_0)(x-x_0)+f(x_0)$.
  \item \emph{Lien avec les variations} : $f'>0$ sur $I$ $\Rightarrow$ $f$ strictement croissante ; $f'<0$ $\Rightarrow$ $f$ strictement décroissante.
\end{itemize}
\end{bilan}

\begin{aretenir}[Checklist avant le Probatoire]
\begin{itemize}
  \item[$\square$] Je sais écrire le taux d'accroissement $\dfrac{f(x_0+h)-f(x_0)}{h}$ sans erreur de parenthèses.
  \item[$\square$] Je sais calculer une limite de taux d'accroissement en simplifiant par $h$.
  \item[$\square$] Je sais donner la définition du nombre dérivé $f'(x_0)$.
  \item[$\square$] Je sais interpréter graphiquement $f'(x_0)$ comme la pente de la tangente.
  \item[$\square$] Je connais par c\oe ur l'équation de la tangente : $y=f'(x_0)(x-x_0)+f(x_0)$.
  \item[$\square$] Je connais les dérivées de $x^n$, $\dfrac{1}{x}$, $\sqrt{x}$ et des polynômes.
  \item[$\square$] Je sais dériver une somme, un produit $ku$, un produit $uv$.
  \item[$\square$] Je sais dériver $\dfrac{1}{v}$ et $\dfrac{u}{v}$ sans oublier le carré au dénominateur.
  \item[$\square$] Je vérifie toujours que le dénominateur ne s'annule pas avant de dériver un quotient.
  \item[$\square$] Je sais construire une tangente à partir de sa pente sur un graphique.
  \item[$\square$] Je sais utiliser le signe de $f'(x)$ pour dresser le tableau de variations de $f$.
  \item[$\square$] Je rédige proprement : limite, conclusion « donc $f$ est dérivable en $x_0$ et $f'(x_0)=\ldots$ ».
\end{itemize}
\end{aretenir}

\findecours

\end{document}
